% ==================================================================
%  徐森林、金亚东、薛春华《数学分析》（第三册）
%  第 12 章  无穷级数
%  源：原书印刷页 1-76（PDF 物理页 = 印刷页 + 6）
% ==================================================================

\chapter{无穷级数}\label{ch:12}

前面接触到的函数主要是初等函数, 有相当多的自然现象和工程技术中的问题需要用这些函数来描述. 但是, 随着科学技术的发展, 人们对自然界的认识逐步深化, 发现许多自然现象不能用初等函数来描述, 特别有很多微分方程的解不能用初等函数来表达, 这就要求人们去构造一些新的函数.

19 世纪上半叶, 数学家普遍认为, 连续函数除了一些特殊点外都是可导的, 他们不能想像有处处连续处处不可导的函数存在. 1875 年 Weierstrass 首先构造出具有上述性质的函数, 使大家对连续与可导的概念在认识上前进了一大步. Weierstrass 构造的这个函数正是用无穷级数 $\sum_{n=1}^{\infty}u_{n}(x)$ 来表达的. 在例~\ref{ex:13.2.14} 中介绍的例子不是 Weierstrass 构造的, 而是由 Van der Waerden 在 1930 年构造的, 它在想法上更直观一些.

由此可见, 无穷级数是构造新函数的一个十分有用的工具. 当然, 随之会有很多新问题: 无穷多个函数如何求和? 如何研究和函数的性质?

% ==================================================================
%  徐森林、金亚东、薛春华《数学分析》（第三册）
%  \S12.1  数项级数
%  源：原书印刷页 1-9（PDF 物理页 7-15）
%  底本：PaddleOCR-VL-1.6，见 chapters/_ch12_12_1_raw.md
%  约定：全西文标点；行间公式用一对方括号；多行公式用 align/aligned；
%        跨册引用一律 \extref；数学模式内的中文必须包 \text{}
%  已订正的 OCR 错误见 tools/_make_ch12_1.py 头部说明
% ==================================================================
\section{数项级数}\label{sec:12.1}

\begin{definition}[无穷级数与部分和]\label{def:12.1.1}
设 $a_{1}, \cdots, a_{n}, \cdots$ 为一实数列, 称形式和 $\sum_{n=1}^{\infty}a_{n}$ 为无穷级数, 称 $a_{n}$ 为该无穷级数的第 $n$ 项(也称为通项), 称
\[
  S_{n}=\sum_{k=1}^{n}a_{k}
\]
为它的第 $n$ 个部分和. 如果
\[
  \lim_{n\to+\infty}S_{n}=\lim_{n\to+\infty}\sum_{k=1}^{n}a_{k}=S\in\mathbb{R},
\]
则称级数 $\sum_{n=1}^{\infty}a_{n}$ 是收敛的, 其和为 $S$, 记作
\[
  \sum_{n=1}^{\infty}a_{n}=S=\lim_{n\to+\infty}S_{n}=\lim_{n\to+\infty}\sum_{k=1}^{n}a_{k}.
\]
否则称 $\sum_{n=1}^{\infty}a_{n}$ 是发散的.

如果 $\lim_{n\to+\infty}S_{n}=+\infty$(或 $-\infty,\infty$), 则称级数 $\sum_{n=1}^{\infty}a_{n}$ 为和为 $+\infty$(或 $-\infty,\infty$), 但此时级数 $\sum_{n=1}^{\infty}a_{n}$ 是发散的. 有时也说级数 $\sum_{n=1}^{\infty}a_{n}$ 发散于 $+\infty$(或 $-\infty,\infty$, 记作 $\sum_{n=1}^{\infty}a_{n}=+\infty$(或 $-\infty,\infty$.)
\end{definition}

上述表明, 由数列 $\{a_{n}\}$ 得出一个新数列 $S_{n}=\sum_{k=1}^{n}a_{k}$; 反之, 任何一个数列 $\{S_{n}\}$, 令 $a_{1}=S_{1},\ a_{2}=S_{2}-S_{1},\cdots,\ a_{n}=S_{n}-S_{n-1}\ (n\geqslant2)$, 就有
\[
  \sum_{k=1}^{n}a_{k}=S_{1}+(S_{2}-S_{1})+\cdots+(S_{n}-S_{n-1})=S_{n},
\]
它是级数 $\sum_{n=1}^{\infty}a_{n}$ 的第 $n$ 个部分和. 于是,
\begin{center}
级数 $\sum_{n=1}^{\infty}a_{n}$ 收敛(发散)
$\Leftrightarrow$ 数列 $S_{n}=\sum_{k=1}^{n}a_{k}$ 收敛(发散).
\end{center}

由此可想象到, 级数 $\sum_{n=1}^{\infty}a_{n}$ 与数列 $\{S_{n}\}$ 之间具有对偶性, 也就是级数 $\sum_{n=1}^{\infty}a_{n}$ 具有的性质(问题)可以翻译成(化为)数列 $\{S_{n}\}$ 的性质(问题); 反之, 数列 $\{S_{n}\}$ 具有的性质(问题)可以翻译成(化为)级数 $\sum_{n=1}^{\infty}a_{n}$ 的性质(问题). 但是, 应该注意的是, 它们各有自己固有的特性和方便的论述, 因此, 即使对数列已积累了大量的结果, 我们还有必要研究级数的各种收敛判别法, 函数项级数一致收敛判别法以及级数的重要性质.

\begin{theorem}[级数收敛的 Cauchy 判别法——级数收敛的 Cauchy 准则]\label{thm:12.1.1}
下列结论是等价的:
\begin{enumerate}
  \item 级数 $\sum_{n=1}^{\infty}a_{n}$ 收敛;
  \item $\forall\varepsilon>0,\ \exists N\in\mathbb{N}$, 当 $n,m\in\mathbb{N},\ m>n$ 时, 有
  \[
    \left|\sum_{k=n+1}^{m}a_{k}\right|<\varepsilon;
\]
  \item $\forall\varepsilon>0,\ \exists N\in\mathbb{N}$, 当 $n>N$ 时, 有
  \[
    \left|\sum_{k=n+1}^{n+p}a_{k}\right|<\varepsilon,\quad \forall p\in\mathbb{N}.
\]
\end{enumerate}
\end{theorem}

\begin{proof}
(1) 级数 $\sum_{n=1}^{\infty}a_{n}$ 收敛, 即 $S_{n}=\sum_{k=1}^{n}a_{k}$ 收敛, 根据数列收敛的 Cauchy 判别法 $\forall\varepsilon>0,\ \exists N\in\mathbb{N}$, 当 $n,m\in\mathbb{N},\ m\geqslant n$ 时, 有
\[
  |S_{n}-S_{m}|<\varepsilon
\]
$\Leftrightarrow(2)$\ $\forall\varepsilon>0,\ \exists N\in\mathbb{N}$, 当 $n,m\in\mathbb{N},\ m>n$ 时, 有
\[
  \Big|\sum_{k=n+1}^{m}a_{k}\Big|\;=|S_{n}-S_{m}|<\varepsilon
\]
$\Leftrightarrow(3)$\ $\forall\varepsilon>0,\ \exists N\in\mathbb{N}$, 当 $n>N$ 时, 有
\[
  \left|\sum_{k=n+1}^{n+p}a_{k}\right|<\varepsilon,\quad \forall p\in\mathbb{N}.
\]
\end{proof}

现在, 我们来重述\extref{定理 1.4.1}.

\begin{theorem}[级数收敛的必要条件]\label{thm:12.1.2}
设级数 $\sum_{n=1}^{\infty}a_{n}$ 收敛, 则 $\lim_{n\to+\infty}a_{n}=0$. 但反之不真.
\end{theorem}

\begin{proof}
解法 1 显然,
\[
  \lim_{n\to+\infty}a_{n}=\lim_{n\to+\infty}(S_{n}-S_{n-1})=\lim_{n\to+\infty}S_{n}-\lim_{n\to+\infty}S_{n-1}=S-S=0.
\]

解法 2 $\forall\varepsilon>0$, 因为级数 $\sum_{n=1}^{\infty}a_{n}$ 收敛, 所以根据定理~\ref{thm:12.1.1} 知, $\exists N\in\mathbb{N}$, 当 $n-1>N$ 时, 取 $p=1$, 有
\[
  |a_{n}-0|=\left|\sum_{k=(n-1)+1}^{(n-1)+1}a_{k}\right|<\varepsilon,
\]
这就证明了 $\lim_{n\to+\infty}a_{n}=0$.

但是, 反之不真. 例如: $a_{n}=\frac{1}{n}$, $\lim_{n\to+\infty}a_{n}=\lim_{n\to+\infty}\frac{1}{n}=0$. 而由\extref{例 1.4.2}, $\sum_{n=1}^{\infty}\frac{1}{n}$ 发散.
\end{proof}

\begin{corollary}\label{cor:12.1.1}
如果 $\lim_{n\to+\infty}a_{n}\neq0$, 则 $\sum_{n=1}^{\infty}a_{n}$ 发散.
\end{corollary}

\begin{proof}
[反证] 假设 $\sum_{n=1}^{\infty}a_{n}$ 收敛, 由定理~\ref{thm:12.1.2}, $\lim_{n\to+\infty}a_{n}=0$, 这与已知 $\lim_{n\to+\infty}a_{n}\neq0$ 相矛盾.
\end{proof}

\begin{theorem}[级数收敛的简单性质]\label{thm:12.1.3}
设级数 $\sum_{n=1}^{\infty}a_{n}$, $\sum_{n=1}^{\infty}b_{n}$ 都收敛, $c$ 为常数, 则级数 $\sum_{n=1}^{\infty}ca_{n}$ 与 $\sum_{n=1}^{\infty}(a_{n}\pm b_{n})$ 也收敛, 且
\[
  \sum_{n=1}^{\infty}ca_{n}=c\sum_{n=1}^{\infty}a_{n},\quad\sum_{n=1}^{\infty}(a_{n}\pm b_{n})=\sum_{n=1}^{\infty}a_{n}\pm\sum_{n=1}^{\infty}b_{n}.
\]
\end{theorem}

\begin{proof}
$\sum_{n=1}^{\infty}ca_{n}=\lim_{n\to+\infty}\sum_{k=1}^{n}ca_{k}=\lim_{n\to+\infty}c\sum_{k=1}^{n}a_{k}=c\lim_{n\to+\infty}\sum_{k=1}^{n}a_{k}=c\sum_{n=1}^{\infty}a_{n}$,
\begin{align*}
\sum_{n=1}^{\infty}(a_{n}\pm b_{n})
  &=\lim_{n\to+\infty}\sum_{k=1}^{n}(a_{k}\pm b_{k})
   =\lim_{n\to+\infty}\left(\sum_{k=1}^{n}a_{k}\pm\sum_{k=1}^{n}b_{k}\right) \\
  &=\lim_{n\to+\infty}\sum_{k=1}^{n}a_{k}\pm\lim_{n\to+\infty}\sum_{k=1}^{n}b_{k}
   =\sum_{n=1}^{\infty}a_{n}\pm\sum_{n=1}^{\infty}b_{n}.
\end{align*}
\end{proof}

\begin{theorem}\label{thm:12.1.4}
级数 $\sum_{n=1}^{\infty}a_{n}$ 去掉或加上有限项, 不影响级数的敛散性.
\end{theorem}

\begin{proof}
设去掉级数 $\sum_{n=1}^{\infty}a_{n}$ 前面 $m$ 项所得级数为 $\sum_{n=m+1}^{\infty}a_{n}$. 记
\[
  S_{n}=\sum_{k=1}^{n}a_{k},\quad S_{n}^{\prime}=\sum_{k=m+1}^{m+n}a_{k}=S_{m+n}-S_{m}.
\]
于是
\[
  \sum_{n=1}^{\infty}a_{n}\text{ 敛散}\ \Leftrightarrow S_{n}\text{ 敛散}\ \Leftrightarrow S_{m+n}\text{ 敛散}\ (m\text{ 固定})\ \Leftrightarrow S_{n}^{\prime}=S_{m+n}-S_{m}\text{ 敛散}\ \Leftrightarrow \sum_{n=m+1}^{\infty}a_{n}\text{ 敛散}.
\]
\end{proof}

\begin{theorem}\label{thm:12.1.5}
设级数 $\sum_{n=1}^{\infty}a_{n}$ 收敛. 如果将级数的项任意归组, 但不改变先后次序, 所得新级数为
\[
  (a_{1}+\cdots+a_{n_{1}})+(a_{n_{1}+1}+\cdots+a_{n_{2}})+\cdots+(a_{n_{k-1}+1}+\cdots+a_{n_{k}})+\cdots,
\]
则新级数也收敛, 且与原级数有相同的和.

注意, 逆命题不成立.
\end{theorem}

\begin{proof}
设原级数的部分和为 $S_{n}$, 则新级数的部分和为 $S_{k}^{\prime}=S_{n_{k}}$, 它为原级数部分和一个子列. 由已知, $S_{n}$ 收敛, 故 $S_{n_{k}}$ 也收敛, 且
\begin{align*}
  (a_{1}+\cdots+a_{n_{1}})+(a_{n_{1}+1}+\cdots+a_{n_{2}})+\cdots+(a_{n_{k-1}+1}+\cdots+a_{n_{k}})+\cdots
  &\quad =\lim_{k\to+\infty}S_{n_{k}} \\
  &\quad =\lim_{n\to+\infty}S_{n}=\sum_{n=1}^{\infty}a_{n}.
\end{align*}

反例: 考察级数 $\sum_{n=1}^{\infty}(-1)^{n-1}=1-1+1-1+1-1+\cdots$, 由于
\[
  \lim_{k\to+\infty}S_{2k-1}=\lim_{k\to+\infty}1=1,\quad\lim_{k\to+\infty}S_{2k}=\lim_{k\to+\infty}0=0,
\]
$\lim_{n\to+\infty}S_{n}$ 不存在, 所以级数 $\sum_{n=1}^{\infty}(-1)^{n-1}$ 发散. 但新级数 $n\to+\infty$ 收敛.
\end{proof}

如果加一个条件, 定理~\ref{thm:12.1.5} 的逆命题就成立了.

\begin{theorem}\label{thm:12.1.6}
设级数
\[
  (a_{1}+\cdots+a_{n_{1}})+(a_{n_{1}+1}+\cdots+a_{n_{2}})+\cdots+(a_{n_{k-1}+1}+\cdots+a_{n_{k}})+\cdots
\]
在同一括号内有相同的符号, 且此级数收敛, 则原级数 $\sum_{n=1}^{\infty}a_{n}$ 也收敛, 且两级数之和相等.
\end{theorem}

\begin{proof}
设 $S_{k}^{\prime}=(a_{1}+\cdots+a_{n_{1}})+(a_{n_{1}+1}+\cdots+a_{n_{k}})+\cdots+(a_{n_{k-1}+1}+\cdots+a_{n_{k}})=S_{n_{k}}$, 由新级数收敛, 故 $\lim_{k\to+\infty}S_{k}^{\prime}=S$. 因为每个括号内是同号的, 所以当 $n$ 从 $n_{k-1}$ 变到 $n_{k}$ 时, 原级数的部分和将单调地在 $S_{k-1}^{\prime}=S_{n_{k-1}}$ 与 $S_{k}^{\prime}=S_{n_{k}}$ 之间变动, 即
\begin{align*}
  S_{k-1}^{\prime}=S_{n_{k-1}}&\leqslant S_{n_{k-1}+1}\leqslant\cdots\leqslant S_{n_{k-1}+1}\leqslant S_{n_{k}}=S_{k}^{\prime}
  &&\text{(当 }a_{n_{k-1}+1}\geqslant0,\cdots,a_{n_{k}}\geqslant0\text{)}, \\
  S_{k-1}^{\prime}=S_{n_{k-1}}&\geqslant S_{n_{k-1}+1}\geqslant\cdots\geqslant S_{n_{k-1}+1}\geqslant S_{n_{k}}=S_{k}^{\prime}
  &&\text{(当 }a_{n_{k-1}+1}\leqslant0,\cdots,a_{n_{k}}\leqslant0\text{)}.
\end{align*}
令 $n\to+\infty$, 则 $k\to+\infty$, 且 $\lim S_{k}^{\prime}=\lim S_{k-1}^{\prime}=S$. 根据极限的定义立知,
\[
  \lim_{n\to+\infty}S_{n}=S.
\]
\end{proof}

\begin{example}[等比级数的收敛性]\label{ex:12.1.1}
讨论等比级数的收敛性.
\end{example}

\begin{proof}
令 $S_{n}=\sum_{k=0}^{n}q^{k}=\frac{1-q^{n+1}}{1-q}$(当 $q\neq1$). 分情况讨论:
\begin{align*}
  q=1:\quad &S_{n}=\sum_{k=0}^{n}1=n+1\to+\infty;\\
  q>1:\quad &\lim_{n\to+\infty}S_{n}=+\infty;\\
  -1<q<1:\quad &\lim_{n\to+\infty}S_{n}=\frac{1}{1-q};\\
  q<-1:\quad &\lim_{n\to+\infty}S_{n}\ \text{不存在};\\
  q=-1:\quad &\lim_{n\to+\infty}\sum_{k=0}^{n}(-1)^{k}=\lim_{n\to+\infty}\frac{1+(-1)^{n}}{2}\ \text{不存在}.
\end{align*}
故等比级数 $\sum_{n=0}^{\infty}q^{n}$ 当且仅当 $|q|<1$ 时收敛, 其和为 $\frac{1}{1-q}$.
\end{proof}

\noindent\textbf{裂项相消法}\quad 将 $a_{n}$ 表示为 $V_{n+1}-V_{n}$(或 $V_{n}-V_{n-1}$)的形式, 如果 $\lim_{n\to+\infty}V_{n+1}$ 存在, 则
\[
  S_{n}=\sum_{k=1}^{n}a_{k}=\sum_{k=1}^{n}(V_{k+1}-V_{k})=V_{n+1}-V_{1}\to\lim_{n\to+\infty}V_{n+1}-V_{1}\quad(n\to+\infty)
\]
或
\[
  S_{n}=\sum_{k=1}^{n}a_{k}=\sum_{k=1}^{n}(V_{k}-V_{k-1})=V_{n}-V_{0}\to\lim_{n\to+\infty}V_{n}-V_{0}\quad(n\to+\infty).
\]

特别地, 若 $a_{n}=\frac{1}{\alpha_{n}\alpha_{n+1}\cdots\alpha_{n+m}}$, 其中 $\alpha_{k}$ 形成公差为 $d$ 的等差数列, 取
\[
  V_{n}=-\frac{1}{md}\frac{1}{\alpha_{n}\alpha_{n+1}\cdots\alpha_{n+m-1}},\quad m\in\mathbb{N},
\]
则有
\begin{align*}
  V_{n+1}-V_{n}
  &=-\frac{1}{md}\left(\frac{1}{\alpha_{n+1}\alpha_{n+2}\cdots\alpha_{n+m}}-\frac{1}{\alpha_{n}\alpha_{n+1}\cdots\alpha_{n+m-1}}\right) \\
  &\reason{通分}=-\frac{1}{md}\frac{\alpha_{n}-\alpha_{n+m}}{\alpha_{n}\alpha_{n+1}\cdots\alpha_{n+m}}
   =-\frac{1}{md}\frac{-md}{md\,\alpha_{n}\alpha_{n+1}\cdots\alpha_{n+m}} \\
  &=\frac{1}{\alpha_{n}\alpha_{n+1}\cdots\alpha_{n+m}}=a_{n}.
\end{align*}

\begin{example}\label{ex:12.1.2}
求下列级数的和:
\begin{enumerate}
  \item $\displaystyle\lim_{n=1}\frac{1}{n(n+1)}$;
  \item $\displaystyle\lim_{n=1}\frac{1}{4n^{2}-1}$;
  \item $\displaystyle\sum_{n=1}^{\infty}\frac{1}{n(n+1)(n+2)}$;
  \item $\displaystyle\sum_{n=1}^{\infty}\arctan\frac{1}{n^{2}+n+1}$.
\end{enumerate}
\end{example}

\begin{proof}
(1)
\begin{align*}
  \sum_{n=1}^{\infty}\frac{1}{n(n+1)}
  &=\lim_{n\to+\infty}\sum_{k=1}^{n}\left(\frac{1}{k}-\frac{1}{k+1}\right)
   \reason{裂项相消}\lim_{n\to+\infty}\left(1-\frac{1}{n+1}\right) \\
  &=1-0=1.
\end{align*}

(2)
\begin{align*}
  \sum_{n=1}^{\infty}\frac{1}{4n^{2}-1}
  &=\lim_{n\to+\infty}\sum_{k=1}^{n}\frac{1}{2}\left(\frac{1}{2k-1}-\frac{1}{2k+1}\right)
   \xlongequal{\text{裂项相消}}\lim_{n\to+\infty}\frac{1}{2}\left(1-\frac{1}{2n+1}\right) \\
  &=\frac{1}{2}(1-0)=\frac{1}{2}.
\end{align*}

(3)
\begin{align*}
  \sum_{n=1}^{\infty}\frac{1}{n(n+1)(n+2)}
  &=\lim_{n\to+\infty}\sum_{k=1}^{n}\frac{1}{k(k+1)(k+2)}
   =\lim_{n\to+\infty}\sum_{k=1}^{n}\frac{1}{2}\left[\frac{1}{k(k+1)}-\frac{1}{(k+1)(k+2)}\right] \\
  &\reason{裂项相消}\lim_{n\to+\infty}\frac{1}{2}\left[\frac{1}{1\times2}-\frac{1}{(n+1)(n+2)}\right] \\
  &=\frac{1}{2}\left(\frac{1}{2}-0\right)=\frac{1}{4}.
\end{align*}

(4)
\begin{align*}
  \sum_{n=1}^{\infty}\arctan\frac{1}{n^{2}+n+1}
  &=\sum_{n=1}^{\infty}\arctan\frac{(n+1)-n}{n^{2}+n+1}
   =\sum_{n=1}^{\infty}\bigl[\arctan(n+1)-\arctan n\bigr] \\
  &=\lim_{n\to+\infty}\sum_{k=1}^{n}\bigl[\arctan(k+1)-\arctan k\bigr]
   =\lim_{n\to+\infty}\bigl[\arctan(n+1)-\arctan 1\bigr] \\
  &=\frac{\pi}{2}-\frac{\pi}{4}=\frac{\pi}{4}.
\end{align*}
\end{proof}

\begin{example}\label{ex:12.1.3}
求级数 $\sum_{n=0}^{\infty}\frac{(-1)^{n}+2}{3^{n}}$ 的和.
\end{example}

\begin{proof}
\begin{align*}
  \sum_{n=0}^{\infty}\frac{(-1)^{n}+2}{3^{n}}
  &=\reason{定理~\ref{thm:12.1.3}}\sum_{n=0}^{\infty}\left(-\frac{1}{3}\right)^{n}+2\sum_{n=0}^{\infty}\left(\frac{1}{3}\right)^{n} \\
  &=\frac{1}{1-\left(-\frac{1}{3}\right)}+2\frac{1}{1-\frac{1}{3}}
   =\frac{3}{4}+3=\frac{15}{4}.
\end{align*}
\end{proof}

\begin{example}\label{ex:12.1.4}
求级数 $\sum_{n=1}^{\infty}\frac{n}{3^{n}}$ 的和.
\end{example}

\begin{proof}
解法 1 设 $S=\sum_{n=1}^{\infty}\frac{n}{3^{n}}$, 则
\[
  \frac{1}{3}S=\sum_{n=1}^{\infty}\frac{n}{3^{n+1}}=\sum_{n=2}^{\infty}\frac{n-1}{3^{n}}.
\]
两式相减得
\begin{align*}
  \frac{2}{3}S=S-\frac{1}{3}S
  &=\sum_{n=1}^{\infty}\frac{n}{3^{n}}-\sum_{n=2}^{\infty}\frac{n-1}{3^{n}}
   =\frac{1}{3}+\sum_{n=2}^{\infty}\frac{1}{3^{n}} \\
  &=\frac{\frac{1}{3}}{1-\frac{1}{3}}=\frac{1}{2},
\end{align*}
\[
  S=\frac{3}{2}\times\frac{1}{2}=\frac{3}{4}.
\]

解法 2 由
\begin{align*}
  \sum_{k=1}^{n}kx^{k-1}
  &=\left(\sum_{k=1}^{n}x^{k}\right)^{\prime}
   =\left(\frac{1-x^{n+1}}{1-x}\right)^{\prime} \\
  &=\frac{-(n+1)x^{n}(1-x)-(1-x^{n+1})(-1)}{(1-x)^{2}}
   =\frac{nx^{n+1}-(n+1)x^{n}+1}{(1-x)^{2}},
\end{align*}
所以
\begin{align*}
  \sum_{n=1}^{\infty}\frac{n}{3^{n}}
  &=\frac{1}{3}\sum_{n=1}^{\infty}n\left(\frac{1}{3}\right)^{n-1}
   =\frac{1}{3}\lim_{n\to+\infty}\sum_{k=1}^{n}k\left(\frac{1}{3}\right)^{k-1} \\
  &=\frac{1}{3}\lim_{n\to+\infty}\frac{n\left(\frac{1}{3}\right)^{n+1}-(n+1)\left(\frac{1}{3}\right)^{n}+1}{\left(1-\frac{1}{3}\right)^{2}}
   =\frac{1}{3}\frac{0+0+1}{\left(\frac{2}{3}\right)^{2}}
   =\frac{3}{4}.
\end{align*}
\end{proof}

\begin{example}\label{ex:12.1.5}
证明: 级数 $\sum_{n=1}^{\infty}n\sin\frac{1}{n}$ 与 $\sum_{n=1}^{\infty}\left(1+\frac{1}{n}\right)^{n}$ 均发散.
\end{example}

\begin{proof}
因为
\[
  \lim_{n\to+\infty}n\sin\frac{1}{n}=\lim_{n\to+\infty}\frac{\sin\frac{1}{n}}{\frac{1}{n}}=1\neq0,
\]
\[
  \lim_{n\to+\infty}\left(1+\frac{1}{n}\right)^{n}=\mathrm{e}\neq0,
\]
所以, 根据定理~\ref{thm:12.1.2}, 级数 $\sum_{n=1}^{\infty}n\sin\frac{1}{n}$ 与 $\sum_{n=1}^{\infty}\left(1+\frac{1}{n}\right)^{n}$ 均发散.
\end{proof}

\begin{example}\label{ex:12.1.6}
设 $a_{n}>0$, $\{a_{n}-a_{n+1}\}$ 为一个严格减的数列. 如果 $\sum_{n=1}^{\infty}a_{n}$ 收敛, 证明:
\[
  \lim_{n\to+\infty}\left(\frac{1}{a_{n+1}}-\frac{1}{a_{n}}\right)=+\infty.
\]
\end{example}

\begin{proof}
因为 $\sum_{n=1}^{\infty}a_{n}$ 收敛, 由定理~\ref{thm:12.1.2} 知, $\lim_{n\to+\infty}a_{n}=0$, 从而 $\lim_{n\to+\infty}(a_{n}-a_{n+1})=0$. 但因 $\{a_{n}-a_{n+1}\}$ 严格减, 故必有 $a_{n}-a_{n+1}>0$, 即 $\{a_{n}\}$ 为严格减的数列. 此外, 由于 $\{a_{n}-a_{n+1}\}$ 严格减, 故有
\begin{align*}
  a_{n}^{2}&=\sum_{k=n}^{\infty}\left(a_{k}^{2}-a_{k+1}^{2}\right)
   =\sum_{k=n}^{\infty}\left(a_{k}-a_{k+1}\right)\left(a_{k}+a_{k+1}\right) \\
  &<\left(a_{n}-a_{n+1}\right)\sum_{k=n}^{\infty}\left(a_{k}+a_{k+1}\right),
\end{align*}
于是
\[
  \frac{1}{a_{n+1}}-\frac{1}{a_{n}}=\frac{a_{n}-a_{n+1}}{a_{n}a_{n+1}}>\frac{a_{n}-a_{n+1}}{a_{n}^{2}}>\frac{1}{\sum_{k=n}^{\infty}(a_{k}+a_{k+1})}\to+\infty\quad(n\to+\infty).
\]
这是因为 $\sum_{k=1}^{\infty}(a_{k}+a_{k+1})$ 收敛, 所以 $\lim_{n\to+\infty}\sum_{k=n}^{\infty}(a_{k}+a_{k+1})=0$. 由此推得
\[
  \lim_{n\to+\infty}\left(\frac{1}{a_{n+1}}-\frac{1}{a_{n}}\right)=+\infty.
\]
\end{proof}

\begin{exblock}{练习题 12.1}
\begin{enumerate}
  \item 研究下列级数的敛散性:
  \[
    \sum_{n=1}^{\infty}\sqrt{n};\quad\sum_{n=1}^{\infty}\frac{1}{\sqrt{n}};\quad\sum_{n=1}^{\infty}\frac{1}{10^{n}};\quad\sum_{n=1}^{\infty}\frac{1}{4n+5};\quad\sum_{n=1}^{\infty}\frac{1}{\ln n}.
\]
  \item 证明下列级数发散:
  \begin{enumerate}
    \item $\displaystyle\sum_{n=1}^{\infty}\frac{n}{n+1}$;
    \item $\displaystyle\sum_{n=1}^{\infty}(-1)^{n}\frac{n^{2}+1}{3n^{2}-2}$;
    \item $\displaystyle\sum_{n=1}^{\infty}\frac{1}{\sqrt[n]{n}}$;
    \item $\displaystyle\sum_{n=1}^{\infty}\left(1-\frac{1}{n}\right)^{n}$;
    \item $\displaystyle\sum_{n=1}^{\infty}\left(1+\frac{1}{n}\right)^{n}$;
    \item $\displaystyle\sum_{n=1}^{\infty}(-1)^{n}n\sin\frac{1}{n}$.
  \end{enumerate}
  \item 求下列级数的和:
  \begin{enumerate}
    \item $1-\frac{1}{2}+\frac{1}{4}-\frac{1}{8}+\cdots+(-1)^{n-1}\frac{1}{2^{n-1}}+\cdots$;
    \item $\frac{1}{1\times3}+\frac{1}{3\times5}+\frac{1}{5\times7}+\cdots+\frac{1}{(2n-1)(2n+1)}+\cdots$;
    \item $\frac{1}{1\times4}+\frac{1}{4\times7}+\frac{1}{7\times10}+\cdots+\frac{1}{(3n-2)(3n+1)}+\cdots$;
    \item $\frac{1}{1\times4}+\frac{1}{2\times5}+\frac{1}{3\times6}+\cdots+\frac{1}{n(n+3)}+\cdots$;
    \item $\frac{1}{2}+\frac{3}{2^{2}}+\frac{5}{2^{3}}+\cdots+\frac{2n-1}{2^{n}}+\cdots$;
    \item $\displaystyle\sum_{n=1}^{\infty}\arctan\frac{2n+1}{n^{2}(n+1)^{2}+1}$.
  \end{enumerate}
  \item 证明下列等式:
  \begin{enumerate}
    \item $\displaystyle\sum_{n=1}^{\infty}\frac{2n+1}{n^{2}(n+1)^{2}}=1$;
    \item $\displaystyle\sum_{n=1}^{\infty}\left(\sqrt{n+2}-2\sqrt{n+1}+\sqrt{n}\right)=1-\sqrt{2}$;
    \item $\displaystyle\sum_{n=1}^{\infty}\ln\frac{n(2n+1)}{(n+1)(2n-1)}=\ln2$;
    \item $\displaystyle\sum_{n=1}^{\infty}\frac{1}{n(n+m)}=\frac{1}{m}\left(1+\frac{1}{2}+\cdots+\frac{1}{m}\right)$, 其中 $m$ 为自然数.
  \end{enumerate}
  \item 作一无穷级数 $\sum_{n=1}^{\infty}a_{n}$, 使其部分和 $S_{n}=\frac{1}{n},\ n=1,2,\cdots$.
  \item 设 $\sum_{n=1}^{\infty}a_{n}$ 收敛, 证明: $\sum_{n=1}^{\infty}(a_{n}+a_{n+1})$ 也收敛. 举例说明, 逆命题不成立. 但若 $a_{n}\geqslant0$, 则逆命题也成立, 试证之.
  \item 设 $\sum_{n=1}^{\infty}a_{n}$ 收敛, 证明: $\sum_{n=1}^{\infty}(a_{n}+b_{n})$ 与 $\sum_{n=1}^{\infty}b_{n}$ 同敛散.
  \item 设 $\sum_{n=1}^{\infty}a_{n}$ 与 $\sum_{n=1}^{\infty}b_{n}$ 都是发散级数, 举例说明下列级数
  \[
    \sum_{n=1}^{\infty}(a_{n}+b_{n}),\quad\sum_{n=1}^{\infty}(a_{n}-b_{n}),\quad\sum_{n=1}^{\infty}a_{n}b_{n},\quad\sum_{n=1}^{\infty}\frac{a_{n}}{b_{n}},
\]
  可能收敛, 也可能发散.
  \item 计算 $\sum_{n=1}^{\infty}n\mathrm{e}^{-nx}$(提示: 考虑 $(1-\mathrm{e}^{-x})S_{n}$).
  \item 求下列级数的和:
  \begin{enumerate}
    \item $\displaystyle\sum_{n=0}^{\infty}\frac{x^{2^{n}}}{1-x^{2^{n+1}}},\ |x|<1$;
    \item $\displaystyle\sum_{k=1}^{\infty}\arctan\frac{1}{2k^{2}}$;
    \item $\displaystyle\sum_{k=1}^{\infty}\arctan\frac{2}{4k^{2}-4k+1}$;
    \item $\displaystyle\sum_{n=1}^{\infty}\frac{1}{n(n+1)(n+1)}$.
  \end{enumerate}
\end{enumerate}
\end{exblock}

\begin{exblock}{思考题 12.1}
\begin{enumerate}
  \item 设 $r$ 为正整数, $pn+q\neq0,\ n=1,2,\cdots$, 证明:
  \[
    \sum_{n=1}^{\infty}\frac{1}{(pn+q)(pn+q+pr)}=\frac{1}{pr}\left(\frac{1}{q+p}+\frac{1}{q+2p}+\cdots+\frac{1}{q+rp}\right).
\]
  \item 设 $m$ 为给定的正整数, 证明:
  \[
    \sum_{n=1}^{\infty}\frac{1}{m^{2}-n^{2}}=-\frac{1}{4m^{2}}.
\]
  \item 应用不等式
  \[
    \ln\left(1+\frac{1}{k}\right)<\frac{1}{k}<\ln\left(1+\frac{1}{k-1}\right),\quad k=2,3,\cdots,
\]
  证明:
  \[
    10<1+\frac{1}{2}+\frac{1}{3}+\cdots+\frac{1}{10^{6}}<20.
\]
  \item 设数列 $\{na_{n}\}$ 与级数 $\sum_{n=1}^{\infty}n(a_{n}-a_{n+1})$ 都收敛. 证明: 级数 $\sum_{n=1}^{\infty}a_{n}$ 也收敛.
  \item (1) 设 $\sum_{n=1}^{\infty}a_{n}$ 为一个收敛级数, 其和为 $S$. 用 $\{S_{n}\}$ 记它的部分和数列. 令
  \[
    \sigma_{n}=\frac{S_{1}+S_{2}+\cdots+S_{n}}{n}.
\]
  证明: $\lim_{n\to+\infty}\sigma_{n}=S$.
  \item (2) 构造一个发散级数 $\sum_{n=1}^{\infty}a_{n}$, 使得由 (1) 所定义的 $\{\sigma_{n}\}$ 却是收敛的.
  \item (3) 设 $\{\sigma_{n}\}$ 是由 (1) 定义的数列. 如果 $\lim_{n\to+\infty}\sigma_{n}=S$, 证明: $a_{n}=o(n),\ n\to+\infty$.
  \item 证明: 级数 $\sum_{n=1}^{\infty}a_{n}$ 收敛 $\Leftrightarrow$ 对于任意的正整数序列 $p_{1},p_{2},\cdots,p_{k},\cdots$ 及自然数的任意子序列 $\{n_{k}\}$, 皆有
  \[
    \lim_{k\to+\infty}\left(a_{n_{k}+1}+a_{n_{k}+2}+\cdots+a_{n_{k}+p_{k}}\right)=0.
\]
\end{enumerate}
\end{exblock}

% ==================================================================
%  徐森林、金亚东、薛春华《数学分析》（第三册）
%  \S12.2  正项级数的判别法（一）判别法体系
%  源：原书印刷页 10-22（PDF 物理页 16-28）
%  底本：PaddleOCR-VL-1.6，见 chapters/_ch12_12_2_raw.md
%  约定：全西文标点；行间公式用一对方括号；多行公式用 align/aligned；
%        跨册引用一律 \extref；数学模式内的中文必须包 \text{}
%  已订正的 OCR 错误见 tools/_make_ch12_2a.py 头部说明
% ==================================================================
\section{正项级数的判别法}\label{sec:12.2}

我们先来给出正项级数的各种判别法. 它们都是根据级数的通项 $a_{n}$ 的性质来判定无穷级数 $\sum_{n=1}^{\infty}a_{n}$ 的敛散性, 从而得到了部分和 $S_{n}=\sum_{k=1}^{n}a_{k}$ 所组成的数列的敛散性. 这也是从 $a_{n}$ 这局部性态反映和 $\sum_{n=1}^{\infty}a_{n}$ 的整体性质.

\begin{definition}\label{def:12.2.1}
设 $a_{n}\geqslant0\ (a_{n}>0),\ n=1,2,\cdots$, 则称级数 $\sum_{n=1}^{\infty}a_{n}$ 为(严格)正项级数.

显然, (严格)正项级数的部分和数列 $S_{n}=S_{n-1}+a_{n}\geqslant S_{n-1}\ (>S_{n-1})$ 是(严格)单调增的. 因此, 有下面的定理.
\end{definition}

\begin{theorem}\label{thm:12.2.1}
正项级数 $\sum_{n=1}^{\infty}a_{n}$ 收敛 $\Leftrightarrow$ 部分和 $S_{n}=\sum_{k=1}^{n}a_{k}$ 有界.
\end{theorem}

\begin{proof}
($\Leftarrow$) 因 $S_{n}$ 单调增有界, 故 $S_{n}$ 收敛, 它等价于 $\sum_{n=1}^{\infty}a_{n}$ 收敛.

($\Rightarrow$) $\sum_{n=1}^{\infty}a_{n}$ 收敛, 它等价于 $S_{n}$ 收敛. 根据\extref{定理 1.2.2} 知, $S_{n}$ 有界.
\end{proof}

由定理~\ref{thm:12.2.1} 可推出下面几个重要的比较判别法.

\begin{theorem}[Cauchy 积分比较判别法]\label{thm:12.2.2}
设 $f(x)$ 为 $[a,+\infty)$ 上的非负单调减的连续函数, 其中 $a>0$. 则级数
\[
  \sum_{n=1}^{\infty}f(a+n)
\]
与广义积分
\[
  \int_{a}^{+\infty}f(x)\,\mathrm{d}x
\]
同敛散.
\end{theorem}

\begin{proof}
由 $f(x)$ 单调减, 故当 $a+k\leqslant x\leqslant a+k+1$ 时, $f(a+k+1)\leqslant f(x)\leqslant f(a+k)$. 于是,
\[
  f(a+k+1)=\int_{a+k}^{a+k+1}f(a+k+1)\,\mathrm{d}x
  \leqslant\int_{a+k}^{a+k+1}f(x)\,\mathrm{d}x
  \leqslant\int_{a+k}^{a+k+1}f(a+k)\,\mathrm{d}x=f(a+k),
  \quad k=0,1,2,\cdots,
\]
\[
  \sum_{k=0}^{n}f(a+k+1)\leqslant\sum_{k=0}^{n}\int_{a+k}^{a+k+1}f(x)\,\mathrm{d}x
  =\int_{a}^{a+n+1}f(x)\,\mathrm{d}x\leqslant\sum_{k=0}^{n}f(a+k).
\]
如果 $\int_{a}^{+\infty}f(x)\,\mathrm{d}x$ 收敛, 则
\[
  \sum_{k=0}^{n}f(a+k+1)\leqslant\int_{a}^{a+n+1}f(x)\,\mathrm{d}x
  \leqslant\int_{a}^{+\infty}f(x)\,\mathrm{d}x<+\infty.
\]
根据定理~\ref{thm:12.2.1} 知, $\sum_{k=0}^{\infty}f(a+k+1)$, 从而 $\sum_{n=1}^{\infty}f(a+n)$ 收敛.

如果 $\sum_{n=1}^{\infty}f(a+n)$ 收敛(即 $\sum_{k=0}^{\infty}f(a+k)$ 收敛), 则 $\forall A>a$, 有充分大的 $n$, 使得 $A<a+n+1$,
\[
  \int_{a}^{A}f(x)\,\mathrm{d}x\leqslant\int_{a}^{a+n+1}f(x)\,\mathrm{d}x
  \leqslant\sum_{k=0}^{n}f(a+k)\leqslant\sum_{k=0}^{\infty}f(a+k)<+\infty.
\]
从而 $\int_{a}^{+\infty}f(x)\,\mathrm{d}x$ 收敛.
\end{proof}

\begin{example}\label{ex:12.2.1}
对于级数
\[
  \sum_{n=1}^{\infty}\frac{1}{n^{x}}=1+\frac{1}{2^{x}}+\frac{1}{3^{x}}+\cdots,
\]
证明: 当 $x>1$ 时, 此级数收敛; 当 $x\leqslant1$ 时, 此级数发散. 于是, $\zeta(x)=\sum_{n=1}^{\infty}\frac{1}{n^{x}}$ 定义了 $(1,+\infty)$ 上的一个函数, 称为 Riemann $\zeta$ 函数.
\end{example}

\begin{proof}
证法 1 参阅\extref{例 1.4.2}.

证法 2 应用定理~\ref{thm:12.2.2} 知
\[
  \sum_{n=1}^{\infty}\frac{1}{n^{x}}\text{ 敛散}\ \Leftrightarrow\ \int_{1}^{+\infty}\frac{\mathrm{d}t}{t^{x}}\text{ 敛散}.
\]
而当 $x>1$ 时, $\int_{1}^{+\infty}\frac{\mathrm{d}t}{t^{x}}$ 收敛; 当 $x\leqslant1$ 时, $\int_{1}^{+\infty}\frac{\mathrm{d}t}{t^{x}}$ 发散.

因此, 当 $x>1$ 时, $\sum_{n=1}^{\infty}\frac{1}{n^{x}}$ 收敛; 当 $x\leqslant1$ 时, $\sum_{n=1}^{\infty}\frac{1}{n^{x}}$ 发散.
\end{proof}

\begin{example}\label{ex:12.2.2}
证明: 级数
\[
  \sum_{n=3}^{\infty}\frac{1}{n^{\alpha}\ln^{\beta}n\ln^{\gamma}\ln n}
\]
当 $\alpha>1$ 时收敛, 当 $\alpha<1$ 时发散; 当 $\alpha=1,\ \beta>1$ 时收敛, 当 $\alpha=1,\ \beta<1$ 时发散; 当 $\alpha=1,\ \beta=1,\ \gamma>1$ 时收敛, 当 $\alpha=1,\ \beta=1,\ \gamma\leqslant1$ 时发散.
\end{example}

\begin{proof}
由\extref{例 6.5.1} 知, 当 $\alpha>1$ 时, $\int_{3}^{+\infty}\frac{\mathrm{d}x}{x^{\alpha}}$ 收敛; 当 $\alpha\leqslant1$ 时, $\int_{3}^{+\infty}\frac{\mathrm{d}x}{x^{\alpha}}$ 发散.

当 $\alpha>1$ 时, 令 $\alpha=1+2\varepsilon,\ \varepsilon>0$, 则
\[
  \frac{1}{x^{\alpha}\ln^{\beta}x\ln^{\gamma}\ln x}
  =\frac{1}{x^{1+\varepsilon}}\frac{1}{x^{\varepsilon}\ln^{\beta}x\ln^{\gamma}\ln x},
\]
\[
  \lim_{x\to+\infty}\frac{\dfrac{1}{x^{\alpha}\ln^{\beta}x\ln^{\gamma}\ln x}}{\dfrac{1}{x^{1+\varepsilon}}}
  =\lim_{x\to+\infty}\frac{1}{x^{\varepsilon}\ln^{\beta}x\ln^{\gamma}\ln x}=0.
\]
因此, 由 $\int_{3}^{+\infty}\frac{\mathrm{d}x}{x^{1+\varepsilon}}$ 收敛知, $\sum_{n=3}^{\infty}\frac{1}{n^{\alpha}\ln^{\beta}n\ln^{\gamma}\ln n}$ 收敛.

当 $\alpha<1$ 时, 令 $\alpha=1-2\varepsilon,\ \varepsilon>0$, 则
\[
  \frac{1}{x^{\alpha}\ln^{\beta}x\ln^{\gamma}\ln x}
  =\frac{1}{x^{1-\varepsilon}}\frac{x^{\varepsilon}}{\ln^{\beta}x\ln^{\gamma}\ln x},
\]
\[
  \lim_{x\to+\infty}\frac{\dfrac{1}{x^{\alpha}\ln^{\beta}x\ln^{\gamma}\ln x}}{\dfrac{1}{x^{1-\varepsilon}}}
  =\lim_{x\to+\infty}\frac{x^{\varepsilon}}{\ln^{\beta}x\ln^{\gamma}\ln x}=+\infty.
\]
因此, 由 $\int_{3}^{+\infty}\frac{\mathrm{d}x}{x^{1-\varepsilon}}$ 发散知, $\sum_{n=3}^{\infty}\frac{1}{n^{\alpha}\ln^{\beta}n\ln^{\gamma}\ln n}$ 发散.

当 $\alpha=1$ 时, 因为
\[
  \int_{3}^{+\infty}\frac{\mathrm{d}x}{x\ln^{\beta}x\ln^{\gamma}\ln x}
  \xrightarrow{u=\ln x}\int_{\ln3}^{+\infty}\frac{\mathrm{d}u}{u^{\beta}\ln^{\gamma}u},
\]
所以, 当 $\beta>1$ 时, 积分收敛; 当 $\beta<1$ 时, 积分发散; 当 $\beta=1$ 时, 积分
\[
  \int_{\ln3}^{+\infty}\frac{\mathrm{d}u}{u\ln^{\gamma}u}
  =\int_{\ln\ln3}^{+\infty}\frac{\mathrm{d}v}{v^{\gamma}}.
\]
因此, 当 $\gamma>1$ 时, 此积分收敛; 当 $\gamma\leqslant1$ 时, 此积分发散.

根据上述结果以及积分判别法(定理~\ref{thm:12.2.2}) 得到级数
\[
  \sum_{n=3}^{\infty}\frac{1}{n^{\alpha}\ln^{\beta}n\ln^{\gamma}\ln n}
\]
当 $\alpha>1$ 时收敛; 当 $\alpha<1$ 时发散; 当 $\alpha=1,\ \beta>1$ 时收敛; 当 $\alpha=1,\ \beta<1$ 时发散; 当 $\alpha=1,\ \beta=1,\ \gamma>1$ 时收敛; 当 $\alpha=1,\ \beta=1,\ \gamma\leqslant1$ 时发散.
\end{proof}

\begin{theorem}[比较判别法]\label{thm:12.2.3}
设 $\sum_{n=1}^{\infty}a_{n}$ 与 $\sum_{n=1}^{\infty}b_{n}$ 都为正项级数. 如果 $\exists N\in\mathbb{N}$, 使得当 $n>N$ 时, 有 $a_{n}\leqslant b_{n}$. 则:
\begin{enumerate}
  \item $\sum_{n=1}^{\infty}b_{n}$ 收敛 $\Rightarrow$ $\sum_{n=1}^{\infty}a_{n}$ 也收敛;
  \item $\sum_{n=1}^{\infty}a_{n}$ 发散 $\Rightarrow$ $\sum_{n=1}^{\infty}b_{n}$ 也发散.
\end{enumerate}

大的级数收敛, 小的级数也收敛; 小的级数发散, 大的级数也发散.
\end{theorem}

\begin{proof}
从 $N+1$ 项开始, 有 $a_{n}\leqslant b_{n}$. 并设 $S_{n}=\sum_{k=1}^{n}a_{k},\ S_{n}^{\prime}=\sum_{k=1}^{n}b_{k}$.

(1) 由 $\sum_{k=1}^{\infty}b_{k}$ 收敛推得 $S_{n}^{\prime}$ 有界, 则当 $n>N+1$ 时, 有
\begin{align*}
  S_{n}&=a_{1}+\cdots+a_{n}=a_{1}+\cdots+a_{N}+a_{N+1}+\cdots+a_{n} \\
       &\leqslant a_{1}+\cdots+a_{N}+b_{N+1}+\cdots+b_{n} \\
       &=\sum_{k=1}^{n}b_{k}+\sum_{k=1}^{N}a_{k}-\sum_{k=1}^{N}b_{k}
        =S_{n}^{\prime}+\sum_{k=1}^{N}a_{k}-\sum_{k=1}^{N}b_{k}<+\infty,
\end{align*}
故 $S_{n}$ 也有界. 根据定理~\ref{thm:12.2.1}, $\sum_{n=1}^{\infty}a_{n}$ 收敛.

(2)(反证) 假设 $\sum_{n=1}^{\infty}b_{n}$ 收敛, 由 (1) 知, $\sum_{n=1}^{\infty}a_{n}$ 也收敛, 这与已知 $\sum_{n=1}^{\infty}b_{n}$ 发散相矛盾.
\end{proof}

\begin{theoremp}[比较判别法的极限形式]\label{thm:12.2.3p}
设 $\sum_{n=1}^{\infty}a_{n}$ 与 $\sum_{n=1}^{\infty}b_{n}$ 为两个正项级数, 且有极限
\[
  \lim_{n\to+\infty}\frac{a_{n}}{b_{n}}=l.
\]
则:
\begin{enumerate}
  \item 若 $0<l<+\infty$, 则 $\sum_{n=1}^{\infty}a_{n}$ 与 $\sum_{n=1}^{\infty}b_{n}$ 同敛散;
  \item 若 $l=0$, 则 $\sum_{n=1}^{\infty}b_{n}$ 收敛蕴涵着 $\sum_{n=1}^{\infty}a_{n}$ 收敛;
  \item 若 $l=+\infty$, 则 $\sum_{n=1}^{\infty}b_{n}$ 发散蕴涵着 $\sum_{n=1}^{\infty}a_{n}$ 发散.
\end{enumerate}
\end{theoremp}

\begin{proof}
(1) 因为 $\lim_{n\to+\infty}\frac{a_{n}}{b_{n}}=l,\ 0<l<+\infty$, 则对 $\varepsilon=\frac{l}{2}$, $\exists N\in\mathbb{N}$, 当 $n>N$ 时,
\[
  0<\frac{l}{2}=l-\frac{l}{2}<\frac{a_{n}}{b_{n}}<l+\frac{l}{2}=\frac{3l}{2},
\]
即
\[
  \frac{l}{2}b_{n}<a_{n}<\frac{3}{2}lb_{n}.
\]
如果 $\sum_{n=1}^{\infty}b_{n}$ 收敛, 则 $\sum_{n=1}^{\infty}\frac{3}{2}lb_{n}$ 也收敛. 根据定理~\ref{thm:12.2.3}(1), $\sum_{n=1}^{\infty}a_{n}$ 收敛. 如果 $\sum_{n=1}^{\infty}b_{n}$ 发散, 则 $\sum_{n=1}^{\infty}\frac{l}{2}b_{n}$ 也发散. 根据定理~\ref{thm:12.2.3}(2), $\sum_{n=1}^{\infty}a_{n}$ 发散. 这就证明了 $\sum_{n=1}^{\infty}a_{n}$ 与 $\sum_{n=1}^{\infty}b_{n}$ 同敛散.

(2) 因为 $\lim_{n\to+\infty}\frac{a_{n}}{b_{n}}=0$, 所以 $\exists N\in\mathbb{N}$, 当 $n>N$ 时,
\[
  0\leqslant\frac{a_{n}}{b_{n}}\leqslant1,\quad\text{即}\quad a_{n}\leqslant b_{n}.
\]
根据定理~\ref{thm:12.2.3}(1), $\sum_{n=1}^{\infty}b_{n}$ 收敛蕴涵着 $\sum_{n=1}^{\infty}a_{n}$ 收敛.

(3) 因为 $\lim_{n\to+\infty}\frac{a_{n}}{b_{n}}=+\infty$, 所以 $\exists N\in\mathbb{N}$, 当 $n>N$ 时,
\[
  \frac{a_{n}}{b_{n}}>1,\quad\text{即}\quad a_{n}>b_{n}>0.
\]
根据定理~\ref{thm:12.2.3}(2), $\sum_{n=1}^{\infty}b_{n}$ 发散蕴涵着 $\sum_{n=1}^{\infty}a_{n}$ 发散.
\end{proof}

应用比较判别法, 就可派生出几个常用的判别法.

将正项级数与几何(即等比)级数比较, 就派生出 Cauchy 判别法与 d'Alembert 判别法.

将正项级数与收敛较慢的级数 $\sum_{n=1}^{\infty}\frac{1}{n^{\alpha}}\ (\alpha>1)$ 作比较, 就派生出 Raabe 判别法.

将正项级数与收敛更慢的级数 $\sum_{n=2}^{\infty}\frac{1}{n\ln^{\alpha}n}\ (\alpha>1)$ 作比较, 就派生出 Gauss 判别法.

一般来说, 用于比较的级数收敛越慢, 所派生出的判别法就越精细, 它能判别的级数越广泛.

\begin{theorem}[Cauchy 判别法]\label{thm:12.2.4}
设 $\sum_{n=1}^{\infty}a_{n}$ 为正项级数.
\begin{enumerate}
  \item 如果 $\exists q\in[0,1)$, 使得
  \[
    \sqrt[n]{a_{n}}\leqslant q<1\quad(n\geqslant N,\ N\in\mathbb{N}\text{ 为某个自然数}),
\]
  则级数 $\sum_{n=1}^{\infty}a_{n}$ 收敛;
  \item 如果 $\sqrt[n]{a_{n}}\geqslant1$ 对无穷个 $n$ 成立, 则级数 $\sum_{n=1}^{\infty}a_{n}$ 发散.
\end{enumerate}
\end{theorem}

\begin{proof}
(1) 由于 $\sqrt[n]{a_{n}}\leqslant q<1\ (n\geqslant N)$, 故
\[
  0\leqslant a_{n}\leqslant q^{n}.
\]
再由等比级数 $\sum_{n=1}^{\infty}q^{n}$ 收敛及定理~\ref{thm:12.2.3}(1) 知, $\sum_{n=1}^{\infty}a_{n}$ 收敛.

(2) 由于 $\sqrt[n]{a_{n}}\geqslant1$ 对无穷个 $n$ 成立, 故 $a_{n}\geqslant1$ 对无穷个 $n$ 成立, 它蕴涵着 $\lim_{n\to+\infty}a_{n}\neq0$. 根据推论~\ref{cor:12.1.1}, 级数 $\sum_{n=1}^{\infty}a_{n}$ 发散.
\end{proof}

\begin{theoremp}[Cauchy 判别法的极限形式]\label{thm:12.2.4p}
设 $\sum_{n=1}^{\infty}a_{n}$ 为正项级数, 且
\[
  \varlimsup_{n\to+\infty}\sqrt[n]{a_{n}}=q,
\]
则:
\begin{enumerate}
  \item 当 $0\leqslant q<1$ 时, 级数 $\sum_{n=1}^{\infty}a_{n}$ 收敛;
  \item 当 $q>1$ 时, 级数 $\sum_{n=1}^{\infty}a_{n}$ 发散;
  \item 当 $q=1$ 时, 不能判定级数 $\sum_{n=1}^{\infty}a_{n}$ 敛散性.
\end{enumerate}
\end{theoremp}

\begin{proof}
(1) 由 $\lim_{n\to+\infty}\sqrt[n]{a_{n}}=q<1$. 取 $\varepsilon>0$, 使 $q+\varepsilon<1$. 于是, $\exists N\in\mathbb{N}$, 当 $n>N$ 时, $\sqrt[n]{a_{n}}<q+\varepsilon<1$. 根据定理~\ref{thm:12.2.4}, $\sum_{n=1}^{\infty}a_{n}$ 收敛.

(2) 由 $\lim_{n\to+\infty}\sqrt[n]{a_{n}}=q>1$, 则存在 $\{a_{n}\}$ 的子列 $\{a_{n_{k}}\}$ 使
\[
  \lim_{k\to+\infty}\sqrt[n_{k}]{a_{n_{k}}}>1.
\]
故有无穷个 $n$ 使 $\sqrt[n]{a_{n}}\geqslant1$. 根据定理~\ref{thm:12.2.4}(2) 知, $\sum_{n=1}^{\infty}a_{n}$ 发散.

(3) 对于级数 $\sum_{n=1}^{\infty}\frac{1}{n}$ 与 $\sum_{n=1}^{\infty}\frac{1}{n^{2}}$ 有
\[
  \lim_{n\to+\infty}\sqrt[n]{\frac{1}{n}}=\lim_{n\to+\infty}\frac{1}{\sqrt[n]{n}}=1,
  \quad\lim_{n\to+\infty}\sqrt[n]{\frac{1}{n^{2}}}
  =\lim_{n\to+\infty}\left(\frac{1}{\sqrt[n]{n}}\right)^{2}=1,
\]
但 $\sum_{n=1}^{\infty}\frac{1}{n}$ 发散, $\sum_{n=1}^{\infty}\frac{1}{n^{2}}$ 收敛.
\end{proof}

\begin{theorem}[d'Alembert 判别法]\label{thm:12.2.5}
设 $\sum_{n=1}^{\infty}a_{n}$ 为严格正项级数.
\begin{enumerate}
  \item 若 $\exists q\in(0,1)$, 使得
  \[
    \frac{a_{n+1}}{a_{n}}\leqslant q<1\quad(n\geqslant N,\ \text{某个 }N\in\mathbb{N}),
\]
  则级数 $\sum_{n=1}^{\infty}a_{n}$ 收敛;
  \item 若 $\frac{a_{n+1}}{a_{n}}\geqslant1\ (n\geqslant N$, 某个 $N\in\mathbb{N})$, 则级数 $\sum_{n=1}^{\infty}a_{n}$ 发散.
\end{enumerate}
\end{theorem}

\begin{proof}
(1) 由条件得
\[
  \frac{a_{N+1}}{a_{N}}\leqslant q,\quad\frac{a_{N+2}}{a_{N+1}}\leqslant q,\quad\cdots,\quad\frac{a_{n}}{a_{n-1}}\leqslant q,\quad\cdots,
\]
于是
\[
  \frac{a_{n}}{a_{N}}=\frac{a_{n}}{a_{n-1}}\cdots\frac{a_{N+2}}{a_{N+1}}\frac{a_{N+1}}{a_{N}}\leqslant q^{n-N},
\]
\[
  0\leqslant a_{n}\leqslant q^{n-N}a_{N}=\frac{a_{N}}{q^{N}}q^{n},\quad n\geqslant N,\quad q\in(0,1).
\]
由等比级数 $\sum_{n=1}^{\infty}q^{n}$ 从而 $\sum_{n=1}^{\infty}\frac{a_{N}}{q^{N}}q^{n}$ 收敛, 以及定理~\ref{thm:12.2.3}(1) 知, $\sum_{n=1}^{\infty}a_{n}$ 收敛.

(2) 由条件, 当 $n\geqslant N$ 时, 有
\[
  a_{n+1}\geqslant a_{n}\geqslant\cdots\geqslant a_{N+1}\geqslant a_{N}.
\]
由此知, $\{a_{n}\}\ (n\geqslant N)$ 单调增. 又因 $a_{n}>0$, 故 $\lim_{n\to+\infty}a_{n}\neq0$. 根据推论~\ref{cor:12.1.1}, $\sum_{n=1}^{\infty}a_{n}$ 发散.
\end{proof}

\begin{theoremp}[d'Alembert 判别法的极限形式]\label{thm:12.2.5p}
设 $\sum_{n=1}^{\infty}a_{n}$ 为严格正项级数.
\begin{enumerate}
  \item 若 $\lim_{n\to+\infty}\frac{a_{n+1}}{a_{n}}=q\in(0,1)$, 则 $\sum_{n=1}^{\infty}a_{n}$ 收敛;
  \item 若 $\lim_{n\to+\infty}\frac{a_{n+1}}{a_{n}}=q'>1$, 则 $\sum_{n=1}^{\infty}a_{n}$ 发散;
  \item 当 $q=1$ 或 $q'=1$ 时, $\sum_{n=1}^{\infty}a_{n}$ 的敛散性不能加以判定.
\end{enumerate}
\end{theoremp}

\begin{proof}
(1) 取 $\varepsilon>0$, 使 $0<q+\varepsilon<1$. 于是, 由 $\lim_{n\to+\infty}\frac{a_{n+1}}{a_{n}}=q$, $\exists N\in\mathbb{N}$, 当 $n\geqslant N$ 时,
\[
  \frac{a_{n+1}}{a_{n}}<q+\varepsilon<1,
\]
则由定理~\ref{thm:12.2.5}(1), $\sum_{n=1}^{\infty}a_{n}$ 收敛.

(2) 先取 $\varepsilon>0$ 使 $q'-\varepsilon>1$. 由下极限定义, $\exists N\in\mathbb{N}$, 当 $n>N$ 时,
\[
  \frac{a_{n+1}}{a_{n}}>q'-\varepsilon>1.
\]
则由定理~\ref{thm:12.2.5}(2) 知, $\sum_{n=1}^{\infty}a_{n}$ 发散.

(3) 对级数 $\sum_{n=1}^{\infty}\frac{1}{n}$ 与 $\sum_{n=1}^{\infty}\frac{1}{n^{2}}$ 有
\[
  \varlimsup_{n\to+\infty}\frac{a_{n+1}}{a_{n}}
  =\varliminf_{n\to+\infty}\frac{a_{n+1}}{a_{n}}
  =\lim_{n\to+\infty}\frac{a_{n+1}}{a_{n}}=1.
\]
但 $\sum_{n=1}^{\infty}\frac{1}{n}$ 发散, $\sum_{n=1}^{\infty}\frac{1}{n^{2}}$ 收敛. 因此, 当 $q=1$ 或 $q'=1$ 时, $\sum_{n=1}^{\infty}a_{n}$ 的敛散性不能判定.
\end{proof}

Cauchy 判别法与 d'Alembert 判别法哪个强些? 根据下面引理~\ref{lem:12.2.1} 可以看出: 凡用 d'Alembert 判别法可判定的, 用 Cauchy 判别法也一定能判定(若 $\lim_{n\to+\infty}\frac{a_{n+1}}{a_{n}}=q<1$, 则 $\lim_{n\to+\infty}\sqrt[n]{a_{n}}\leqslant\lim_{n\to+\infty}\frac{a_{n+1}}{a_{n}}=q<1$; 若 $\lim_{n\to+\infty}\frac{a_{n+1}}{a_{n}}=q'>1$, 则 $\lim_{n\to+\infty}\sqrt[n]{a_{n}}\geqslant\lim_{n\to+\infty}\frac{a_{n+1}}{a_{n}}=q'\geqslant1$). 但反之不然(参阅例~\ref{ex:12.2.14}). 由此可见, Cauchy 判别法比 d'Alembert 判别法适用的面要宽一些. 理论上, Cauchy 判别法比 d'Alembert 判别法要强! 但在有些场合下, 使用 d'Alembert 判别法要方便些, 特别是 $\frac{a_{n+1}}{a_{n}}$ 简单或 $\lim_{n\to+\infty}\frac{a_{n+1}}{a_{n}}$ 好算的情形! 因此, 它还是受欢迎的.

但总的来说, 这两个判别法的适用面都不宽, 原因是它们只能判别比某个几何(等比)级数收敛得还快的级数(参阅引理~\ref{lem:12.2.2}). 所谓级数 $\sum_{n=1}^{\infty}a_{n}$ 比 $\sum_{n=1}^{\infty}b_{n}$ 收敛得快, 是指
\[
  \lim_{n\to+\infty}\frac{a_{n}}{b_{n}}=0.
\]
例如: 当 $0<q<1$, $\alpha>1$, $\beta>1$ 时, $\sum_{n=1}^{\infty}q^{n}$ 比 $\sum_{n=1}^{\infty}\frac{1}{n^{\alpha}}$ 收敛得快; $\sum_{n=1}^{\infty}\frac{1}{n^{\alpha}}$ 比 $\sum_{n=1}^{\infty}\frac{1}{n(\ln n)^{\beta}}$ 收敛得快.

\begin{lemma}\label{lem:12.2.1}
设 $a_{n}>0\ (n=1,2,\cdots)$, 则
\[
  \varliminf_{n\to+\infty}\frac{a_{n+1}}{a_{n}}
  \leqslant\varliminf_{n\to+\infty}\sqrt[n]{a_{n}}
  \leqslant\varlimsup_{n\to+\infty}\sqrt[n]{a_{n}}
  \leqslant\varlimsup_{n\to+\infty}\frac{a_{n+1}}{a_{n}}.
\]
\end{lemma}

\begin{proof}
先证 $\varlimsup_{n\to+\infty}\sqrt[n]{a_{n}}\leqslant\varlimsup_{n\to+\infty}\frac{a_{n+1}}{a_{n}}$.

设 $\lim_{n\to+\infty}\frac{a_{n+1}}{a_{n}}=q$. 若 $q=+\infty$, 则不等式自然成立.

若 $0<q<+\infty$, 则 $\forall\varepsilon>0$, $\exists N\in\mathbb{N}$, 当 $n\geqslant N$ 时, 有
\[
  \frac{a_{n+1}}{a_{n}}<q+\varepsilon.
\]
因而,
\[
  \frac{a_{N+1}}{a_{N}}<q+\varepsilon,\quad\frac{a_{N+2}}{a_{N+1}}<q+\varepsilon,\quad\cdots,\quad\frac{a_{N+k}}{a_{N+k-1}}<q+\varepsilon.
\]
各式相乘得
\[
  \frac{a_{N+k}}{a_{N}}<(q+\varepsilon)^{k},\quad\text{即}\quad a_{N+k}<a_{N}(q+\varepsilon)^{k},
\]
或
\[
  a_{n}<a_{N}(q+\varepsilon)^{n-N}=a_{N}(q+\varepsilon)^{-N}(q+\varepsilon)^{n}.
\]
于是
\[
  \sqrt[n]{a_{n}}<\sqrt[n]{a_{N}(q+\varepsilon)^{-N}}\,(q+\varepsilon),
\]
所以, $\varlimsup_{n\to+\infty}\sqrt[n]{a_{n}}\leqslant q+\varepsilon$. 令 $\varepsilon\to0^{+}$ 得到
\[
  \varlimsup_{n\to+\infty}\sqrt[n]{a_{n}}\leqslant q=\varlimsup_{n\to+\infty}\frac{a_{n+1}}{a_{n}}.
\]

再证 $\varliminf_{n\to+\infty}\frac{a_{n+1}}{a_{n}}\leqslant\varliminf_{n\to+\infty}\sqrt[n]{a_{n}}$.

若 $\lim_{n\to+\infty}\frac{a_{n+1}}{a_{n}}=0$, 则不等式显然成立.

若 $\lim_{n\to+\infty}\frac{a_{n+1}}{a_{n}}=q>0$, 则 $\forall\varepsilon>0$ 且 $q-\varepsilon>0$, $\exists N\in\mathbb{N}$, 当 $n\geqslant N$ 时, 有
\[
  0<q-\varepsilon<\frac{a_{n+1}}{a_{n}}.
\]
特别地, 有
\[
  \frac{a_{N+1}}{a_{N}}>q-\varepsilon,\quad\frac{a_{N+2}}{a_{N+1}}>q-\varepsilon,\quad\cdots,\quad\frac{a_{n}}{a_{n-1}}>q-\varepsilon.
\]
各式相乘得
\[
  \frac{a_{n}}{a_{N}}>(q-\varepsilon)^{n-N},\quad\text{即}\quad a_{n}>a_{N}(q-\varepsilon)^{-N}(q-\varepsilon)^{n},
\]
于是
\[
  \sqrt[n]{a_{n}}>\sqrt[n]{a_{N}(q-\varepsilon)^{-N}}\,(q-\varepsilon),
\]
所以, $\varliminf_{n\to+\infty}\sqrt[n]{a_{n}}\geqslant q-\varepsilon$. 再令 $\varepsilon\to0^{+}$ 得到
\[
  \varliminf_{n\to+\infty}\sqrt[n]{a_{n}}\geqslant q=\varliminf_{n\to+\infty}\frac{a_{n+1}}{a_{n}}.
\]
\end{proof}

\begin{lemma}\label{lem:12.2.2}
设级数 $\sum_{n=1}^{\infty}a_{n}$ 满足 Cauchy 判别法或 d'Alembert 判别法判定收敛的条件($\overline{\lim}_{n\to+\infty}\sqrt[n]{a_{n}}=q<1$ 或 $\overline{\lim}_{n\to+\infty}\frac{a_{n+1}}{a_{n}}=q<1$), 则必 $\exists r\in(q,1)$, 使得
\[
  \lim_{n\to+\infty}\frac{a_{n}}{r^{n}}=0,
\]
即 $\sum_{n=1}^{\infty}a_{n}$ 比几何(等比)级数 $\sum_{n=1}^{\infty}r^{n}$ 收敛得快.
\end{lemma}

\begin{proof}
(1) 取 $r,\ \varepsilon>0$, 使 $q+\varepsilon<r<1$. 因为
\[
  \varlimsup_{n\to+\infty}\sqrt[n]{a_{n}}=q<1,
\]
所以, $\exists N\in\mathbb{N}$, 当 $n\geqslant N$ 时, $\sqrt[n]{a_{n}}<q+\varepsilon<r$, 即 $a_{n}<(q+\varepsilon)^{n}<r^{n}$. 于是,
\[
  \frac{a_{n}}{r^{n}}=\frac{a_{n}}{(q+\varepsilon)^{n}}\left(\frac{q+\varepsilon}{r}\right)^{n}\to0\quad(n\to+\infty).
\]
这就表明 $\sum_{n=1}^{\infty}a_{n}$ 比 $\sum_{n=1}^{\infty}r^{n}$ 收敛得快.

(2) $\lim_{n\to+\infty}\frac{a_{n+1}}{a_{n}}=q<1$, 则 $\exists N\in\mathbb{N}$, 当 $n\geqslant N$ 时, $\frac{a_{n+1}}{a_{n}}<q+\varepsilon<r$, 所以
\[
  a_{n}=\frac{a_{n}}{a_{n-1}}\frac{a_{n-1}}{a_{n-2}}\cdots\frac{a_{N+1}}{a_{N}}a_{N}
  <a_{N}(q+\varepsilon)^{n-N}=a_{N}(q+\varepsilon)^{-N}(q+\varepsilon)^{n}\leqslant a_{N}r^{-N}r^{n}.
\]
于是
\[
  \frac{a_{n}}{r^{n}}=\frac{a_{n}}{(q+\varepsilon)^{n}}\left(\frac{q+\varepsilon}{r}\right)^{n}\to0\quad(n\to+\infty).
\]
这就表明 $\sum_{n=1}^{\infty}a_{n}$ 比 $\sum_{n=1}^{\infty}r^{n}$ 收敛得快.
\end{proof}

如果 $\sum_{n=1}^{\infty}a_{n}$ 比任何几何(等比)级数收敛得慢(如 $\sum_{n=1}^{\infty}\frac{1}{n^{\alpha}},\ \alpha>1$), 必须有更精细的判别法. $\sum_{n=1}^{\infty}a_{n}$ 与 $\sum_{n=1}^{\infty}\frac{1}{n^{\alpha}}$ 相比较得到 Raabe 判别法.

\begin{theorem}[Raabe 判别法]\label{thm:12.2.6}
设 $\sum_{n=1}^{\infty}a_{n}$ 为严格正项级数.
\begin{enumerate}
  \item 若 $\exists r>1,\ \exists N_{0}\in\mathbb{N}$, 使得当 $n\geqslant N_{0}$ 时, 有
  \[
    n\left(\frac{a_{n}}{a_{n+1}}-1\right)\geqslant r,
\]
  则级数 $\sum_{n=1}^{\infty}a_{n}$ 收敛;
  \item 若 $\exists N_{0}\in\mathbb{N}$, 使得当 $n\geqslant N_{0}$ 时, 有
  \[
    n\left(\frac{a_{n}}{a_{n+1}}-1\right)\leqslant1,
\]
  则级数 $\sum_{n=1}^{\infty}a_{n}$ 发散.
\end{enumerate}
\end{theorem}

\begin{proof}
(1) 取 $\alpha\in(1,r)$. 易知,
\[
  \lim_{n\to+\infty}\frac{\left(1+\frac{1}{n}\right)^{\alpha}-1}{\frac{1}{n}}
  \xlongequal{\text{L'Hospital 法则}}\lim_{x\to0}\frac{(1+x)^{\alpha}-1}{x}
  =\lim_{x\to0}\frac{\alpha(1+x)^{\alpha-1}}{1}
  =\alpha<r.
\]
故 $\exists N\in\mathbb{N}$, 当 $n>N$ 时,
\[
  \frac{\left(1+\frac{1}{n}\right)^{\alpha}-1}{\frac{1}{n}}<r,
\]
\[
  \left(1+\frac{1}{n}\right)^{\alpha}<1+\frac{r}{n}.
\]
由条件, 当 $n>\max\{N_{0},N\}$ 时, 有
\[
  \frac{a_{n}}{a_{n+1}}\geqslant1+\frac{r}{n}>\left(1+\frac{1}{n}\right)^{\alpha}=\frac{(n+1)^{\alpha}}{n^{\alpha}},
\]
\[
  n^{\alpha}a_{n}>(n+1)^{\alpha}a_{n+1},
\]
即 $\{n^{\alpha}a_{n}\mid n>\max\{N_{0},N_{1}\}\}$ 单调减且大于零, 故 $\{n^{\alpha}a_{n}\}$ 有界, 从而 $\exists M>0$ 使
\[
  n^{\alpha}a_{n}>M\quad(n=1,2,\cdots),\qquad\text{即}\quad a_{n}<\frac{M}{n^{\alpha}}\quad(n=1,2,\cdots).
\]
由 $\alpha>1$ 知 $\sum_{n=1}^{\infty}\frac{M}{n^{\alpha}}$ 收敛, 根据定理~\ref{thm:12.2.3}(1) 得到 $\sum_{n=1}^{\infty}a_{n}$ 是收敛的.

(2) 由条件, 当 $n>N_{0}$ 时,
\[
  \frac{a_{n}}{a_{n+1}}\leqslant1+\frac{1}{n}=\frac{n+1}{n},
\]
\[
  na_{n}\leqslant(n+1)a_{n+1},
\]
即数列 $\{na_{n}\mid n\geqslant N_{0}\}$ 单调增. 于是,
\[
  na_{n}\geqslant N_{0}a_{N_{0}},\quad n\geqslant N_{0},
\]
\[
  a_{n}\geqslant\frac{N_{0}a_{N_{0}}}{n}>0,\quad n\geqslant N_{0}.
\]
从 $\sum_{n=1}^{\infty}\frac{N_{0}a_{N_{0}}}{n}$ 发散及定理~\ref{thm:12.2.3}(2) 知 $\sum_{n=1}^{\infty}a_{n}$ 发散.
\end{proof}

\begin{theoremp}[Raabe 判别法的极限形式]\label{thm:12.2.6p}
设 $\sum_{n=1}^{\infty}a_{n}$ 为严格正项级数.
\begin{enumerate}
  \item 若 $\lim_{n\to+\infty}n\left(\frac{a_{n}}{a_{n+1}}-1\right)=l>1$, 则 $\sum_{n=1}^{\infty}a_{n}$ 收敛;
  \item 若 $\lim_{n\to+\infty}n\left(\frac{a_{n}}{a_{n+1}}-1\right)=l'<1$, 则 $\sum_{n=1}^{\infty}a_{n}$ 发散;
  \item 若 $l=1$ 或 $l'=1$, 则 $\sum_{n=1}^{\infty}a_{n}$ 不能断定敛散.
\end{enumerate}
\end{theoremp}

\begin{proof}
(1) 由于 $\lim_{n\to+\infty}n\left(\frac{a_{n}}{a_{n+1}}-1\right)=l>1$, 取 $r\in(1,l)$, 则 $\exists N_{0}\in\mathbb{N}$, 当 $n\geqslant N_{0}$, 有
\[
  n\left(\frac{a_{n}}{a_{n+1}}-1\right)>r>1.
\]
由定理~\ref{thm:12.2.6}(1) 知, 级数 $\sum_{n=1}^{\infty}a_{n}$ 收敛.

(2) 因为 $\lim_{n\to+\infty}n\left(\frac{a_{n}}{a_{n+1}}-1\right)=l'<1$, 故 $\exists N_{0}\in\mathbb{N}$, 当 $n\geqslant N_{0}$ 时,
\[
  n\left(\frac{a_{n}}{a_{n+1}}-1\right)<1.
\]
由定理~\ref{thm:12.2.6}(2) 知, 级数 $\sum_{n=1}^{\infty}a_{n}$ 发散.

(3) 考虑级数 $\sum_{n=2}^{\infty}\frac{1}{n\ln^{2}n}$ 与 $\sum_{n=2}^{\infty}\frac{1}{n}$. 显然,
\begin{align*}
  n\left(\frac{1}{n\ln^{2}n}-\frac{1}{(n+1)\ln^{2}(n+1)}-1\right)
  &=n\left(\frac{(n+1)\ln^{2}(n+1)}{n\ln^{2}n}-1\right) \\
  &=\frac{(n+1)\ln^{2}(n+1)-n\ln^{2}n}{\ln^{2}n} \\
  &=\frac{\ln^{2}(n+1)}{\ln^{2}n}
   +\frac{n\bigl[\ln(n+1)+\ln n\bigr]\bigl[\ln(n+1)-\ln n\bigr]}{\ln^{2}n} \\
  &=\frac{\ln^{2}(n+1)}{\ln^{2}n}
   +\frac{\bigl[\ln(n+1)+\ln n\bigr]\ln\left(1+\frac{1}{n}\right)^{n}}{\ln^{2}n}
   \to1+0\cdot1 \\
  &=1\quad(n\to+\infty),
\end{align*}
\[
  n\left(\frac{1}{n}-\frac{1}{n+1}-1\right)=n\frac{(n+1)-n}{n}=1\to1\quad(n\to+\infty).
\]
但是, $\sum_{n=2}^{\infty}\frac{1}{n}$ 发散; 而由
\[
  \int_{2}^{+\infty}\frac{\mathrm{d}x}{x\ln^{2}x}
  =\int_{2}^{+\infty}\frac{\mathrm{d}\ln x}{\ln^{2}x}
  =-\left.\frac{1}{\ln x}\right|_{2}^{+\infty}
  =\frac{1}{\ln2}
\]
收敛及积分比较判别法可知, $\sum_{n=2}^{\infty}\frac{1}{n\ln^{2}n}$ 收敛.
\end{proof}

如果级数 $\sum_{n=1}^{\infty}a_{n}$ 与比 $\sum_{n=2}^{\infty}\frac{1}{n^{\alpha}}\ (\alpha>1)$ 收敛得更慢的级数 $\sum_{n=2}^{\infty}\frac{1}{n\ln^{\beta}n}\ (\beta>1)$ 相比较就得到 Gauss 判别法.

\begin{theorem}[Gauss 判别法]\label{thm:12.2.7}
设 $\sum_{n=1}^{\infty}a_{n}$ 为严格正项级数, 且
\[
  \frac{a_{n}}{a_{n+1}}=1+\frac{1}{n}+\frac{\delta}{n\ln n}+o\left(\frac{1}{n\ln n}\right),\quad n\to+\infty.
\]
则级数 $\sum_{n=1}^{\infty}a_{n}$:
\begin{enumerate}
  \item 当 $\delta>1$ 时, 收敛;
  \item 当 $\delta<1$ 时, 发散;
  \item 当 $\delta=1$ 时, 不能判定敛散性.
\end{enumerate}
\end{theorem}

\begin{proof}
(1) 因为 $\delta>1$, 故可取 $\beta$ 使 $\delta>\beta>1$. 下面证明: $\exists N\in\mathbb{N}$, 当 $n\geqslant N$ 时, 有 $a_{n}\leqslant\frac{M}{n\ln^{\beta}n}$($M$ 为正的常数), 从而由例~\ref{ex:12.2.2} 知 $\sum_{n=2}^{\infty}\frac{M}{n\ln^{\beta}n}$ 收敛, 再由比较判别法(定理~\ref{thm:12.2.3}) 推得 $\sum_{n=1}^{\infty}a_{n}$ 收敛.

因为
\[
  \frac{\ln(n+1)}{\ln n}
  =\frac{\ln n+\ln\left(1+\frac{1}{n}\right)}{\ln n}
  =1+\frac{\frac{1}{n}+o\left(\frac{1}{n}\right)}{\ln n}
  =1+\frac{1}{n\ln n}+o\left(\frac{1}{n\ln n}\right),
\]
所以,
\[
  \left[\frac{\ln(n+1)}{\ln n}\right]^{\beta}
  =\left[1+\frac{1}{n\ln n}+o\left(\frac{1}{n\ln n}\right)\right]^{\beta}
  =1+\frac{\beta}{n\ln n}+o\left(\frac{1}{n\ln n}\right).
\]
于是,
\begin{align*}
  \frac{n+1}{n}\left[\frac{\ln(n+1)}{\ln n}\right]^{\beta}
  &=\left(1+\frac{1}{n}\right)\left[1+\frac{\beta}{n\ln n}+o\left(\frac{1}{n\ln n}\right)\right] \\
  &=1+\frac{1}{n}+\frac{\beta}{n\ln n}+\frac{\beta}{n^{2}\ln n}+o\left(\frac{1}{n\ln n}\right) \\
  &=1+\frac{1}{n}+\frac{\beta}{n\ln n}+o\left(\frac{1}{n\ln n}\right).
\end{align*}
由此推得
\begin{align*}
  \frac{a_{n}}{a_{n+1}}-\frac{n+1}{n}\left[\frac{\ln(n+1)}{\ln n}\right]^{\beta}
  &\xlongequal{\text{题设}}\left[1+\frac{1}{n}+\frac{\delta}{n\ln n}
     +o\left(\frac{1}{n\ln n}\right)\right]
   -\left[1+\frac{1}{n}+\frac{\beta}{n\ln n}+o\left(\frac{1}{n\ln n}\right)\right] \\
  &=\frac{\delta-\beta}{n\ln n}+o\left(\frac{1}{n\ln n}\right)
   =\frac{1}{n\ln n}\left[\delta-\beta
     +\frac{o\left(\frac{1}{n\ln n}\right)}{n\ln n}\right]
   \quad(n\to+\infty).
\end{align*}
因为 $\delta-\beta>0$, 故 $\exists N\in\mathbb{N}$, 当 $n\geqslant N$ 时, 上式为正, 即
\[
  \frac{a_{n}}{a_{n+1}}-\frac{n+1}{n}\left[\frac{\ln(n+1)}{\ln n}\right]^{\beta}>0,
\]
于是
\[
  n\ln^{\beta}n\,a_{n}>(n+1)\ln^{\beta}(n+1)\,a_{n+1}.
\]
这就表明数列 $\{n\ln^{\beta}n\,a_{n}\mid n\geqslant N\}$ 单调减有下界 $0$, 所以它有界, 即 $\exists M>0$, 使得
\[
  n\ln^{\beta}n\,a_{n}<M,
\]
\[
  a_{n}<\frac{M}{n\ln^{\beta}n}.
\]
由例~\ref{ex:12.2.2}, $\sum_{n=2}^{\infty}\frac{M}{n\ln^{\beta}n}\ (\beta>1)$ 收敛. 根据比较判别法(定理~\ref{thm:12.2.3}(1)) 得到 $\sum_{n=1}^{\infty}a_{n}$ 收敛.

(2) 因为 $\delta<1$, 在 (1) 的式中取 $\beta=1$ 就有
\[
  \frac{a_{n}}{a_{n+1}}-\frac{n+1}{n}\frac{\ln(n+1)}{\ln n}
  =\frac{\delta-1}{n\ln n}+o\left(\frac{1}{n\ln n}\right)\quad(n\to+\infty).
\]
故 $\exists N\in\mathbb{N}$, 当 $n\geqslant N$ 时, 上式为负, 即
\[
  \frac{a_{n}}{a_{n+1}}-\frac{n+1}{n}\frac{\ln(n+1)}{\ln n}<0,
\]
于是
\[
  n\ln n\,a_{n}<(n+1)\ln(n+1)\,a_{n+1}.
\]
这就表明数列 $\{n\ln n\,a_{n}\mid n\geqslant N\}$ 单调增, 从而
\[
  n\ln n\,a_{n}\geqslant N\ln N\,a_{N}>0,
\]
\[
  a_{n}\geqslant\frac{N\ln N\,a_{N}}{n\ln n}.
\]
由 $\sum_{n=2}^{\infty}\frac{N\ln N\,a_{N}}{n\ln n}$ 发散与比较判别法(定理~\ref{thm:12.2.3}(2)) 得到 $\sum_{n=1}^{\infty}a_{n}$ 发散.

(3) 由例~\ref{ex:12.2.2} 知, $\sum_{n=2}^{\infty}\frac{1}{n\ln n}$ 发散, 而 $\sum_{n=3}^{\infty}\frac{1}{n\ln n(\ln\ln n)^{2}}$ 收敛.

另一方面, 由于
\[
  \frac{\ln(n+1)}{\ln n}
  =\frac{\ln n+\ln\left(1+\frac{1}{n}\right)}{\ln n}
  =\frac{\ln n+\frac{1}{n}+o\left(\frac{1}{n}\right)}{\ln n}
  =1+\frac{1}{n\ln n}+o\left(\frac{1}{n\ln n}\right),
\]
\begin{align*}
  \frac{\ln\ln(n+1)}{\ln\ln n}
  &=\frac{\ln\left[\ln n+\frac{1}{n}+o\left(\frac{1}{n}\right)\right]}{\ln\ln n} \\
  &=\frac{\ln\left\{\ln n\left[1+\frac{1}{n\ln n}+o\left(\frac{1}{n\ln n}\right)\right]\right\}}{\ln\ln n} \\
  &=\frac{\ln\ln n+\ln\left[1+\frac{1}{n\ln n}+o\left(\frac{1}{n\ln n}\right)\right]}{\ln\ln n} \\
  &=\frac{\ln\ln n+\frac{1}{n\ln n}+o\left(\frac{1}{n\ln n}\right)}{\ln\ln n}
   =1+o\left(\frac{1}{n\ln n}\right),
\end{align*}
我们有
\begin{align*}
  \frac{\dfrac{1}{n\ln n}}{\dfrac{1}{(n+1)\ln(n+1)}}
  &=\left(1+\frac{1}{n}\right)\frac{\ln(n+1)}{\ln n}
   =\left(1+\frac{1}{n}\right)\left[1+\frac{1}{n\ln n}+o\left(\frac{1}{n\ln n}\right)\right] \\
  &=1+\frac{1}{n}+\frac{1}{n\ln n}+o\left(\frac{1}{n\ln n}\right),
\end{align*}
\begin{align*}
  \frac{\dfrac{1}{n\ln n(\ln\ln n)^{2}}}
        {\dfrac{1}{(n+1)\ln(n+1)\left[\ln\ln(n+1)\right]^{2}}}
  &=\left(1+\frac{1}{n}\right)\frac{\ln(n+1)}{\ln n}
   \frac{\left[\ln\ln(n+1)\right]^{2}}{(\ln\ln n)^{2}} \\
  &=\left[1+\frac{1}{n}+\frac{1}{n\ln n}+o\left(\frac{1}{n\ln n}\right)\right]
   \left[1+o\left(\frac{1}{n\ln n}\right)\right]^{2} \\
  &=1+\frac{1}{n}+\frac{1}{n\ln n}+o\left(\frac{1}{n\ln n}\right).
\end{align*}
这表明级数 $\sum_{n=2}^{\infty}\frac{1}{n\ln n}$ 与 $\sum_{n=3}^{\infty}\frac{1}{n\ln n(\ln\ln n)^{2}}$ 都不能用 Gauss 判别法判定.
\end{proof}

完全类似的证明可以得到推广的 Gauss 判别法.

\begin{theoremp}[推广的 Gauss 判别法]\label{thm:12.2.7p}
设 $\sum_{n=1}^{\infty}a_{n}$ 为严格正项级数, 且
\[
  \frac{a_{n}}{a_{n+1}}=1+\frac{1}{n}+\frac{\delta_{n}}{n\ln n}
   +o\left(\frac{1}{n\ln n}\right),\quad n\to+\infty.
\]
如果 $\lim_{n\to+\infty}\delta_{n}=\delta\in\mathbb{R}$, 则级数 $\sum_{n=1}^{\infty}a_{n}$:
\begin{enumerate}
  \item 当 $\delta>1$ 时, 收敛;
  \item 当 $\delta<1$ 时, 发散;
  \item 当 $\delta=1$ 时, 不能判定敛散性.
\end{enumerate}
\end{theoremp}

\begin{remark}\label{rem:12.2.1}
如果定理~\ref{thm:12.2.7} 中, $\delta=1$, 则应用 Gauss 判别法不能判定级数 $\sum_{n=1}^{\infty}a_{n}$ 的敛散性. 若要用比较判别法, 则需找出一个比 $\sum_{n=2}^{\infty}\frac{1}{n\ln^{\beta}n}\ (\beta>1)$ 收敛得更慢的级数来比较, 例如 $\sum_{n=2}^{\infty}\frac{1}{n\ln n\ln^{\gamma}\ln n}\ (\gamma>1)$ 就是这样的级数.
\end{remark}

根据下面的例~\ref{ex:12.2.3} 与例~\ref{ex:12.2.17}, 收敛得最慢的级数是不存在的. 因此, 与一个已知级数作比较来建立判断一切级数的敛散性的"万能"判别法是不存在的. 不过上面介绍的这些判别法, 对于相当广泛的级数已经足够用了.
% ===== PART-B =====

\begin{example}\label{ex:12.2.3}
设 $\sum_{n=1}^{\infty}a_{n}$ 为收敛的严格正项级数, 如果
\[
  \lim_{n\to+\infty}\sqrt[n]{a_{n}}=q\in\left[0,1\right)
  \quad\text{或}\quad
  \varlimsup_{n\to+\infty}\frac{a_{n+1}}{a_{n}}=q\in\left[0,1\right),
\]
则
\[
  \varlimsup_{n\to+\infty}\sqrt[n]{na_{n}}=q\in[0,1)
  \quad\text{或}\quad
  \varlimsup_{n\to+\infty}\frac{(n+1)a_{n+1}}{na_{n}}=q\in[0,1),
\]
因此 $\sum_{n=1}^{\infty}na_{n}$ 也收敛. 此外, 由
\[
  \varlimsup_{n\to+\infty}\frac{a_{n}}{na_{n}}=\varlimsup_{n\to+\infty}\frac{1}{n}=0
\]
知 $\sum_{n=1}^{\infty}na_{n}$ 比 $\sum_{n=1}^{\infty}a_{n}$ 收敛得更慢.
\end{example}

\begin{example}\label{ex:12.2.4}
证明: 级数 $\sum_{n=1}^{\infty}\frac{1}{n!}$ 收敛.
\end{example}

\begin{proof}
证法 1 当 $n\geqslant2$ 时, $\frac{1}{n!}\leqslant\frac{1}{(n-1)n}$, 由 $\sum_{n=2}^{\infty}\frac{1}{(n-1)n}$ 收敛及比较判别法定理~\ref{thm:12.2.3}(1) 知, $\sum_{n=1}^{\infty}\frac{1}{n!}$ 也收敛.

证法 2 因为 $\frac{a_{n+1}}{a_{n}}=\frac{n!}{(n+1)!}=\frac{1}{n+1}\leqslant\frac{1}{2}<1$, 故由 d'Alembert 判别法得到 $\sum_{n=1}^{\infty}\frac{1}{n!}$ 收敛.

证法 3 因为几何平均数小于算术平均数, 故
\[
  0<\sqrt[n]{\frac{1}{n!}}
  \leqslant\frac{1+\frac{1}{2}+\cdots+\frac{1}{n}}{n}
  \rightarrow\lim_{n\to+\infty}\frac{1}{n}=0\quad(n\to+\infty),
\]
所以 $\lim_{n\to+\infty}\sqrt[n]{\frac{1}{n!}}=0<1$. 根据 Cauchy 判别法的极限形式知, $\sum_{n=1}^{\infty}\frac{1}{n!}$ 收敛.
\end{proof}

\begin{example}\label{ex:12.2.5}
证明: 级数 $\sum_{n=1}^{\infty}\frac{1}{\sqrt{n(n^{2}+1)}}$ 收敛.
\end{example}

\begin{proof}
因为
\[
  \frac{1}{\sqrt{n(n^{2}+1)}}<\frac{1}{\sqrt{n\cdot n^{2}}}=\frac{1}{n^{3/2}},
\]
又 $\sum_{n=1}^{\infty}\frac{1}{n^{3/2}}$ 收敛, 根据比较判别法定理~\ref{thm:12.2.3}(1) 知, 级数 $\sum_{n=1}^{\infty}\frac{1}{\sqrt{n(n^{2}+1)}}$ 收敛.
\end{proof}

\begin{example}\label{ex:12.2.6}
证明: 级数 $\sum_{n=2}^{\infty}\frac{1}{(\ln n)^{\ln n}}$ 收敛.
\end{example}

\begin{proof}
因为当 $n>\mathrm{e}^{\mathrm{e}^{2}}$ 时, 有
\[
  (\ln n)^{\ln n}=\mathrm{e}^{\ln n\ln\ln n}
  =\left(\mathrm{e}^{\ln n}\right)^{\ln\ln n}
  =n^{\ln\ln n}>n^{2},
\]
所以,
\[
  \frac{1}{(\ln n)^{\ln n}}<\frac{1}{n^{2}}.
\]
再由 $\sum_{n=2}^{\infty}\frac{1}{n^{2}}$ 收敛及比较判别法定理~\ref{thm:12.2.3}(1) 得到级数 $\sum_{n=2}^{\infty}\frac{1}{(\ln n)^{\ln n}}$ 收敛.
\end{proof}

\begin{example}\label{ex:12.2.7}
研究级数 $\sum_{n=1}^{\infty}\left(1-\cos\frac{x}{n}\right)$ 的敛散性.
\end{example}

\begin{proof}
因为
\begin{align*}
  \lim_{n\to+\infty}\frac{1-\cos\dfrac{x}{n}}{\dfrac{1}{n^{2}}}
  &=\lim_{n\to+\infty}\frac{2\sin^{2}\dfrac{x}{2n}}{1} \\
  &=\left\{\begin{array}{ll}
  0, & x=0,\\[
2pt]
  \displaystyle\lim_{n\to+\infty}
  \left(\frac{\sin\dfrac{x}{2n}}{\dfrac{x}{2n}}\right)^{2}\cdot\frac{x^{2}}{2}, & x\neq0,
  \end{array}\right. \\
  &=\frac{x^{2}}{2}\geqslant0,
\end{align*}
所以由 $\sum_{n=1}^{\infty}\frac{1}{n^{2}}$ 收敛及比较判别法的极限形式(定理~\ref{thm:12.2.3p}(1)、(2)) 知, 级数 $\sum_{n=1}^{\infty}\left(1-\cos\frac{x}{n}\right)$ 收敛.
\end{proof}
\begin{example}\label{ex:12.2.8}
讨论级数 $\sum_{n=1}^{\infty}\left[\frac{1}{n}-\ln\left(1+\frac{1}{n}\right)\right]$ 的敛散性.
\end{example}
\begin{proof}
从\extref{定理 1.4.2}(3) 中不等式
\[
  \frac{1}{n+1}<\ln\left(1+\frac{1}{n}\right)<\frac{1}{n},
\]
得到
\[
  0<\frac{1}{n}-\ln\left(1+\frac{1}{n}\right)
  <\frac{1}{n}-\frac{1}{n+1}=\frac{1}{n(n+1)}<\frac{1}{n^{2}}.
\]
再由 $\sum_{n=1}^{\infty}\frac{1}{n^{2}}$ 收敛及比较判别法定理~\ref{thm:12.2.3}(1) 知, $\sum_{n=1}^{\infty}\left[\frac{1}{n}-\ln\left(1+\frac{1}{n}\right)\right]$ 收敛, 其和记为
\begin{align*}
  C&=\sum_{n=1}^{\infty}\left[\frac{1}{n}-\ln\left(1+\frac{1}{n}\right)\right]
   =\lim_{n\to+\infty}\sum_{k=1}^{n}\left[\frac{1}{k}-\ln\left(1+\frac{1}{k}\right)\right] \\
  &=\lim_{n\to+\infty}\left[\sum_{k=1}^{n}\frac{1}{k}-\sum_{k=1}^{n}\bigl(\ln(k+1)-\ln k\bigr)\right] \\
  &=\lim_{n\to+\infty}\left[\sum_{k=1}^{n}\frac{1}{k}-\ln(n+1)\right]
   =\lim_{n\to+\infty}\left[\sum_{k=1}^{n}\frac{1}{k}-\ln n-\ln\left(1+\frac{1}{n}\right)\right] \\
  &=\lim_{n\to+\infty}\left(\sum_{k=1}^{n}\frac{1}{k}-\ln n\right),
\end{align*}
并称它为 Euler 常数, 可证 $C=0.577215\cdots$. 记
\[
  \varepsilon_{n}=1+\frac{1}{2}+\cdots+\frac{1}{n}-\ln n-C,
\]
则 $\lim_{n\to+\infty}\varepsilon_{n}=C-C=0$, 且有
\[
  1+\frac{1}{2}+\cdots+\frac{1}{n}=\ln n+C+\varepsilon_{n},
\]
这是调和级数 $\sum_{n=1}^{\infty}\frac{1}{n}$ 部分和的一个渐近表达式. 由此得到
\[
  \lim_{n\to+\infty}\frac{1+\frac{1}{2}+\cdots+\frac{1}{n}}{\ln n}
  =\lim_{n\to+\infty}\frac{\ln n+C+\varepsilon_{n}}{\ln n}
  =\lim_{n\to+\infty}\left(1+\frac{C+\varepsilon_{n}}{\ln n}\right)=1+0=1.
\]
\end{proof}

\begin{example}\label{ex:12.2.9}
判别级数 $\sum_{n=1}^{\infty}\frac{1}{2^{n}}\left(1+\frac{1}{n}\right)^{n^{2}}$ 的敛散性.
\end{example}

\begin{proof}
解法 1 因为 $\lim_{n\to+\infty}\sqrt[n]{a_{n}}=\lim_{n\to+\infty}\frac{1}{2}\left(1+\frac{1}{n}\right)^{n}=\frac{\mathrm{e}}{2}>1$, 故应用 Cauchy 判别法的极限形式(定理~\ref{thm:12.2.4p}(2)) 知, 级数 $\sum_{n=1}^{\infty}\frac{1}{2^{n}}\left(1+\frac{1}{n}\right)^{n^{2}}$ 发散.

解法 2 因为 $\lim_{n\to+\infty}a_{n}=\lim_{n\to+\infty}\frac{1}{2^{n}}\left(1+\frac{1}{n}\right)^{n}=\lim_{n\to+\infty}\left[\frac{(1+\frac{1}{n})^{n}}{2}\right]^{n}=+\infty\neq0$, 故由推论~\ref{cor:12.1.1} 推得 $\sum_{n=1}^{\infty}\frac{1}{2^{n}}\left(1+\frac{1}{n}\right)^{n^{2}}$ 发散.
\end{proof}

\begin{example}\label{ex:12.2.10}
讨论级数 $\sum_{n=1}^{\infty}n!x^{n}\ (x\geqslant0)$ 的敛散性.
\end{example}

\begin{proof}
因为当 $x>0$ 时,
\[
  \lim_{n\to+\infty}\frac{a_{n+1}}{a_{n}}
  =\lim_{n\to+\infty}\frac{(n+1)!x^{n+1}}{n!x^{n}}
  =\lim_{n\to+\infty}(n+1)x=+\infty,
\]
故根据 d'Alembert 判别法的极限形式(定理~\ref{thm:12.2.5p}(2)) 知, $\sum_{n=1}^{\infty}n!x^{n}\ (x\geqslant0)$ 仅当 $x=0$ 时收敛.
\end{proof}

\begin{example}\label{ex:12.2.11}
讨论级数 $\sum_{n=1}^{\infty}\frac{n!}{n^{n}}x^{n}\ (x\geqslant0)$ 的敛散性.
\end{example}

\begin{proof}
因为当 $x>0$ 时,
\[
  \lim_{n\to+\infty}\frac{a_{n+1}}{a_{n}}
  =\lim_{n\to+\infty}
  \frac{\dfrac{(n+1)!}{(n+1)^{n+1}}x^{n+1}}{\dfrac{n!}{n^{n}}x^{n}}
  =\lim_{n\to+\infty}\frac{x}{\left(1+\frac{1}{n}\right)^{n}}
  =\frac{x}{\mathrm{e}}
  \begin{cases}
  <1, & 0<x<\mathrm{e},\\
  =1, & x=\mathrm{e},\\
  >1, & x>\mathrm{e},
\end{cases}
\]
故根据 d'Alembert 判别法的极限形式(定理~\ref{thm:12.2.5p}(1)、(2)) 知, 级数 $\sum_{n=1}^{\infty}\frac{n!}{n^{n}}x^{n}$ 当 $0<x<\mathrm{e}$ 时收敛; 当 $x>\mathrm{e}$ 时此级数发散.

当 $x=\mathrm{e}$ 时, $\lim_{n\to+\infty}\frac{a_{n+1}}{a_{n}}=\frac{\mathrm{e}}{\mathrm{e}}=1$, 不能应用定理~\ref{thm:12.2.5} 判定. 但由于
\[
  \frac{a_{n+1}}{a_{n}}=\frac{x}{\left(1+\frac{1}{n}\right)^{n}}=\frac{\mathrm{e}}{\left(1+\frac{1}{n}\right)^{n}}>1,
\]
根据 d'Alembert 判别法(定理~\ref{thm:12.2.5}(2)) 知, 级数 $\sum_{n=1}^{\infty}\frac{n!}{n^{n}}\mathrm{e}^{n}$ 发散.

综上得到级数 $\sum_{n=1}^{\infty}\frac{n!}{n^{n}}x^{n}$ 仅当 $0\leqslant x<\mathrm{e}$ 时收敛.
\end{proof}

\begin{example}\label{ex:12.2.12}
讨论级数 $\sum_{n=1}^{\infty}\frac{n!}{(x+1)(x+2)\cdots(x+n)}$ 当 $x\geqslant0$ 时的敛散性.
\end{example}

\begin{proof}
由于 $\lim_{n\to+\infty}\frac{a_{n+1}}{a_{n}}=\lim_{n\to+\infty}\frac{n+1}{n+1+x}=1$($x$ 固定), 故不能用 d'Alembert 判别法判定. 因为
\[
  \lim_{n\to+\infty}n\left(\frac{a_{n}}{a_{n+1}}-1\right)=\lim_{n\to+\infty}\frac{nx}{n+1}=x,
\]
所以根据 Raabe 判别法的极限形式(定理~\ref{thm:12.2.6p}(1)、(2)) 知, 当 $x>1$ 时, 级数 $\sum_{n=1}^{\infty}\frac{n!}{(x+1)(x+2)\cdots(x+n)}$ 收敛; 当 $0\leqslant x<1$ 时此级数发散.

当 $x=1$ 时, 虽然用 Raabe 判别法不能判定, 但是原级数
\[
  \sum_{n=1}^{\infty}\frac{n!}{(1+1)(1+2)\cdots(1+n)}
  =\sum_{n=1}^{\infty}\frac{1}{n+1}
\]
发散.

综上得到级数 $\sum_{n=1}^{\infty}\frac{n!}{(x+1)(x+2)\cdots(x+n)}$ 仅当 $x>1$ 时收敛.
\end{proof}

\begin{example}\label{ex:12.2.13}
设 $p>0$, $q>0$, 问: $p,q$ 取何值时, 级数
\[
  \sum_{n=1}^{\infty}\frac{p(p+1)\cdots(p+n-1)}{n!}\frac{1}{n^{q}}
\]
收敛.
\end{example}

\begin{proof}
因为
\begin{align*}
  \frac{a_{n}}{a_{n+1}}
  &=\frac{n+1}{n+p}\left(1+\frac{1}{n}\right)^{q}
   =\left(1+\frac{1-p}{n+p}\right)\left(1+\frac{q}{n}+o\left(\frac{1}{n}\right)\right) \\
  &=1+\frac{q}{n}+\frac{1-p}{n+p}+o\left(\frac{1}{n}\right)
   =1+\frac{q+1-p}{n}+o\left(\frac{1}{n}\right),
\end{align*}
故
\[
  n\left(\frac{a_{n}}{a_{n+1}}-1\right)
  =q-p+1+o\left(\frac{1}{n}\right)\to q-p+1\quad(n\to+\infty).
\]
所以, 根据 Raabe 判别法, 级数当 $q>p$ 时收敛; 当 $q<p$ 时此级数发散.
\end{proof}

\begin{example}\label{ex:12.2.14}
证明: 级数
\[
  \frac{1}{2}+\frac{1}{3}+\frac{1}{2^{2}}+\frac{1}{3^{2}}+\frac{1}{2^{3}}+\frac{1}{3^{3}}
  +\frac{1}{2^{4}}+\frac{1}{3^{4}}+\cdots
\]
收敛.
\end{example}

\begin{proof}
因为
\[
  \varlimsup_{n\to+\infty}\sqrt[n]{a_{n}}
  =\lim_{n\to+\infty}\sqrt[2n-1]{\frac{1}{2^{n}}}
  =\lim_{n\to+\infty}\left(\frac{1}{2}\right)^{\frac{n}{2n-1}}
  =\left(\frac{1}{2}\right)^{\frac{1}{2}}
  =\frac{1}{\sqrt{2}}<1,
\]
所以由 Cauchy 判别法的极限形式(定理~\ref{thm:12.2.4p}(1)) 知, 该级数收敛. 但是,
\[
  \lim_{n\to+\infty}\frac{a_{n+1}}{a_{n}}
  =\lim_{n\to+\infty}\frac{\left(\frac{1}{3}\right)^{n}}{\left(\frac{1}{2}\right)^{n}}
  =\lim_{n\to+\infty}\left(\frac{2}{3}\right)^{n}=0,
\]
\[
  \varlimsup_{n\to+\infty}\frac{a_{n+1}}{a_{n}}
  =\lim_{n\to+\infty}\frac{\dfrac{1}{3^{n+1}}}{\dfrac{1}{2^{n}}}
  =\frac{1}{2}\lim_{n\to+\infty}\left(\frac{3}{2}\right)^{n}=+\infty,
\]
故用 d'Alembert 判别法的极限形式不能判定此级数的敛散性.
\end{proof}

\begin{example}\label{ex:12.2.15}
讨论级数 $\sum_{n=1}^{\infty}\frac{x^{n}}{(1+x)(1+x^{2})\cdots(1+x^{n})}\ (x\geqslant0)$ 的敛散性.
\end{example}

\begin{proof}
当 $x>0$ 时, 由
\begin{align*}
  \lim_{n\to+\infty}\frac{a_{n+1}}{a_{n}}
  &=\lim_{n\to+\infty}\frac{x}{1+x^{n+1}}
   =\begin{cases}
   x, & 0<x<1,\\[
2pt]
   \dfrac{1}{2}, & x=1,\\[2pt]
   0, & x>1
   \end{cases} <1,
\end{align*}
根据 d'Alembert 判别法的极限形式(定理~\ref{thm:12.2.5p}(1)) 知, 原级数当 $x\geqslant0$ 时都收敛(其中 $x=0$ 时, 原级数为 $\sum_{n=1}^{\infty}0$ 收敛).
\end{proof}
\begin{example}\label{ex:12.2.16}
研究超几何级数
\[
  F(\alpha,\beta,\gamma,x)=1+\sum_{n=1}^{\infty}
  \frac{\alpha(\alpha+1)\cdots(\alpha+n-1)\beta(\beta+1)\cdots(\beta+n-1)}
       {n!\gamma(\gamma+1)\cdots(\gamma+n-1)}x^{n}
\]
的敛散性, 其中 $\alpha,\beta,\gamma$ 皆为正数, $x\geqslant0$.
\end{example}

\begin{proof}
令
\[
  a_{n}=\frac{\alpha(\alpha+1)\cdots(\alpha+n-1)\beta(\beta+1)\cdots(\beta+n-1)}
             {n!\gamma(\gamma+1)\cdots(\gamma+n-1)}x^{n},
\]
则
\[
  \lim_{n\to+\infty}\frac{a_{n+1}}{a_{n}}
  =\lim_{n\to+\infty}\frac{(n+\alpha)(n+\beta)x}{(n+1)(n+\gamma)}=x.
\]
根据 d'Alembert 判别法的极限形式(定理~\ref{thm:12.2.5p}(1)、(2)):
\begin{enumerate}
  \item[(1)] 当 $0\leqslant x<1$ 时, $F(\alpha,\beta,\gamma,x)$ 收敛;
  \item[(2)] 当 $x>1$ 时, $F(\alpha,\beta,\gamma,x)$ 发散;
  \item[(3)] 当 $x=1$ 时, 有
\end{enumerate}
\begin{align*}
  \lim_{n\to+\infty}n\left(\frac{a_{n}}{a_{n+1}}-1\right)
  &=\lim_{n\to+\infty}n\left[\frac{(n+1)(n+\gamma)}{(n+\alpha)(n+\beta)}-1\right] \\
  &=\lim_{n\to+\infty}
   \frac{(\gamma+1-\alpha-\beta)n^{2}+(\gamma-\alpha\beta)n}
        {n^{2}+(\alpha+\beta)n+\alpha\beta} \\
  &=\gamma+1-\alpha-\beta.
\end{align*}
应用 Raabe 判别法的极限形式(定理~\ref{thm:12.2.6p}(1)、(2)) 推得:
\begin{enumerate}
  \item[①] 当 $\gamma+1-\alpha-\beta>1$(即 $\gamma>\alpha+\beta$)时, $F(\alpha,\beta,\gamma,x)$ 收敛;
  \item[②] 当 $\gamma+1-\alpha-\beta<1$(即 $\gamma<\alpha+\beta$)时, $F(\alpha,\beta,\gamma,x)$ 发散;
  \item[③] 当 $\gamma+1-\alpha-\beta=1$(即 $\gamma=\alpha+\beta$)时, 应进一步讨论.
\end{enumerate}

因为 $\alpha,\beta$ 皆为正数, 所以
\begin{align*}
  n\left(\frac{a_{n}}{a_{n+1}}-1\right)-1
  &=\frac{(\gamma+1-\alpha-\beta)n^{2}+(\gamma-\alpha\beta)n}{n^{2}+(\alpha+\beta)n+\alpha\beta}-1 \\
  &=\frac{n^{2}+(\gamma-\alpha\beta)n}{n^{2}+(\alpha+\beta)n+\alpha\beta}-1
   =\frac{(\gamma-\alpha-\beta-\alpha\beta)n-\alpha\beta}{n^{2}+(\alpha+\beta)n+\alpha\beta} \\
  &=\frac{-\alpha\beta n-\alpha\beta}{n^{2}+(\alpha+\beta)n+\alpha\beta}<0.
\end{align*}
于是, $\forall n\in\mathbb{N}$, 恒有
\[
  n\left(\frac{a_{n}}{a_{n+1}}-1\right)<1.
\]
根据 Raabe 判别法(定理~\ref{thm:12.2.6}) 可知 $F(\alpha,\beta,\gamma,x)$ 发散.

或由
\begin{align*}
  \frac{a_{n}}{a_{n+1}}
  &=\frac{(n+1)(n+\gamma)}{(n+\alpha)(n+\beta)}
   =\frac{n^{2}+(\gamma+1)n+\gamma}{n^{2}+(\alpha+\beta)n+\alpha\beta} \\
  &=1+\frac{(\gamma+1-\alpha-\beta)n+\gamma-\alpha\beta}{n^{2}+(\alpha+\beta)n+\alpha\beta} \\
  &=1+\frac{n+\gamma-\alpha\beta}{n^{2}+(\alpha+\beta)n+\alpha\beta} \\
  &=1+\frac{1}{n}+\frac{0}{n\ln n}
   +\frac{-\alpha\beta n-\alpha\beta}{n\bigl[n^{2}+(\alpha+\beta)n+\alpha\beta\bigr]} \\
  &=1+\frac{1}{n}+\frac{0}{n\ln n}+o\left(\frac{1}{n\ln n}\right),
\end{align*}
根据 Gauss 判别法(定理~\ref{thm:12.2.7}(2)), $F(\alpha,\beta,\gamma,x)$ 发散.
\end{proof}

\begin{example}\label{ex:12.2.17}
设 $\sum_{n=1}^{\infty}a_{n}$ 为一个收敛的正项级数, 试作一收敛的严格正项级数 $\sum_{n=1}^{\infty}b_{n}$, 使得
\[
  \lim_{n\to+\infty}\frac{a_{n}}{b_{n}}=0.
\]
这表明收敛得最慢的正项级数是不存在的.
\end{example}

\begin{proof}
第 1 步 设 $\sum_{n=1}^{\infty}a_{n}$ 为严格正项级数, 记 $\sum_{n=1}^{\infty}a_{n}=S$, $S_{0}=0$, $S_{n}=\sum_{k=1}^{n}a_{k}$, $\beta_{n}=S-S_{n-1}$, 则 $\beta_{n}\ (>0)$ 严格减且 $\lim_{n\to+\infty}\beta_{n}=\lim_{n\to+\infty}(S-S_{n-1})=S-S=0$. 令 $b_{n}=\sqrt{\beta_{n}}-\sqrt{\beta_{n+1}}$, 则 $b_{n}>0$ 且
\[
  \sum_{k=1}^{n}b_{k}
  =\sum_{k=1}^{n}\left(\sqrt{\beta_{k}}-\sqrt{\beta_{k+1}}\right)
  =\sqrt{\beta_{1}}-\sqrt{\beta_{n+1}}
  \rightarrow\sqrt{\beta_{1}}-0=\sqrt{\beta_{1}}\quad(n\to+\infty),
\]
即 $\sum_{n=1}^{\infty}b_{n}=\sqrt{\beta_{1}}=\sqrt{S-S_{0}}=\sqrt{S}$. 此外,
\begin{align*}
  \frac{a_{n}}{b_{n}}
  &=\frac{a_{n}}{\sqrt{\beta_{n}}-\sqrt{\beta_{n+1}}}
   =\frac{a_{n}\left(\sqrt{\beta_{n}}+\sqrt{\beta_{n+1}}\right)}{\beta_{n}-\beta_{n+1}} \\
  &\quad =\frac{a_{n}\left(\sqrt{\beta_{n}}+\sqrt{\beta_{n+1}}\right)}{S_{n}-S_{n-1}}
   =\sqrt{\beta_{n}}+\sqrt{\beta_{n+1}}\rightarrow0\quad(n\to+\infty).
\end{align*}
所以, $\sum_{n=1}^{\infty}b_{n}$ 就是要找的正项的级数.

第 2 步 设 $\sum_{n=1}^{\infty}a_{n}$ 为正项级数, 则 $a_{n}\geqslant0$. 令
\[
  \tilde{a}_{n}=
  \begin{cases}
  a_{n}, & a_{n}>0,\\[2pt]
  \dfrac{1}{2^{n}}, & a_{n}=0,
\end{cases}
\]
则 $\sum_{n=1}^{\infty}\tilde{a}_{n}$ 为收敛的严格正项级数. 根据第 1 步, 存在收敛的严格正项级数 $\sum_{n=1}^{\infty}b_{n}$, 使得,
\[
  \lim_{n\to+\infty}\frac{\tilde{a}_{n}}{b_{n}}=0.
\]
由于
\[
  0\leqslant\frac{a_{n}}{b_{n}}\leqslant\frac{\tilde{a}_{n}}{b_{n}}\to0\quad(n\to+\infty)
\]
及夹逼定理知,
\[
  \lim_{n\to+\infty}\frac{b_{n}}{a_{n}}=0.
\]
因此, $\sum_{n=1}^{\infty}b_{n}$ 就是比 $\sum_{n=1}^{\infty}a_{n}$ 收敛得慢的严格正项级数.
\end{proof}

\begin{exblock}{练习题 12.2}
\begin{enumerate}
  \item 用比较判别法讨论下列级数的敛散性:
  \begin{enumerate}
    \item $\displaystyle\sum_{n=1}^{\infty}\frac{1}{3n^{2}+5}$;
    \item $\displaystyle\sum_{n=1}^{\infty}\frac{1}{n2^{n}}$;
    \item $\displaystyle\sum_{n=1}^{\infty}\left(\frac{n^{2}}{3n^{2}+1}\right)^{n}$;
    \item $\displaystyle\sum_{n=1}^{\infty}\frac{1}{n}\sin\frac{1}{n}$;
    \item $\displaystyle\sum_{n=1}^{\infty}\frac{n+1}{n(n+2)}$;
    \item $\displaystyle\sum_{n=1}^{\infty}\frac{1}{n^{1+\frac{1}{n}}}$;
    \item $\displaystyle\sum_{n=2}^{\infty}\frac{1}{(\ln n)^{\ln n}}$;
    \item $\displaystyle\sum_{n=2}^{\infty}\frac{1}{(\ln n)^{\ln\ln n}}$;
    \item $\displaystyle\sum_{n=1}^{\infty}\left(n\ln\frac{2n+1}{2n-1}-1\right)$;
    \item $\displaystyle\sum_{n=1}^{\infty}\frac{(-1)^{n}\sin\sqrt{n^{2}+1}\,\pi}{n^{\alpha}}$.
  \end{enumerate}
  \item 讨论下列级数的敛散性:
  \begin{enumerate}
    \item $\displaystyle\sum_{n=1}^{\infty}n\tan\frac{\pi}{2^{n+1}}$;
    \item $\displaystyle\sum_{n=1}^{\infty}\frac{n^{2}}{3^{n}}$;
    \item $\displaystyle\sum_{n=1}^{\infty}\frac{n^{5}}{3^{n}}\left[\sqrt{3}+(-1)^{n}\right]^{n}$;
    \item $\displaystyle\sum_{n=1}^{\infty}\frac{n^{2}}{\left(1+\frac{1}{n}\right)^{n}}$;
    \item $\displaystyle\sum_{n=1}^{\infty}\frac{n^{n+\frac{1}{n}}}{\left(n+\frac{1}{n}\right)^{n}}$;
    \item $\displaystyle\sum_{n=1}^{\infty}\frac{x^{n}}{n!}\ (x\geqslant0)$;
    \item $\displaystyle\sum_{n=1}^{\infty}\left(\frac{n-4}{3n+1}\right)^{n}$;
    \item $\displaystyle\sum_{n=2}^{\infty}\frac{n^{\ln n}}{(\ln n)^{n}}$.
  \end{enumerate}
  \item 应用 Raabe 判别法, 讨论下列级数的敛散性:
  \begin{enumerate}
    \item $\displaystyle\sum_{n=1}^{\infty}\frac{\sqrt{n!}}{(a+\sqrt{1})(a+\sqrt{2})\cdots(a+\sqrt{n})},\ a>0$;
    \item $\displaystyle\sum_{n=1}^{\infty}\frac{n!n^{-p}}{q(q+1)\cdots(q+n)},\ p>0,\ q>0$.
  \end{enumerate}
  \item 证明: 级数 $\sum_{n=2}^{\infty}\frac{1}{n\ln n\ln\ln n}$ 发散.
  \item 问 $p,q$ 取何值时, 级数 $\sum_{n=3}^{\infty}\frac{1}{n(\ln n)^{p}(\ln\ln n)^{q}}$ 收敛.
  \item 设正项级数 $\sum_{n=1}^{\infty}a_{n}$ 收敛, 证明: $\sum_{n=1}^{\infty}a_{n}^{2}$ 也收敛. 但反之不然, 举例说明之. 构造一个收敛级数 $\sum_{n=1}^{\infty}a_{n}$, 但 $\sum_{n=1}^{\infty}a_{n}^{2}$ 发散.
  \item 设 $\sum_{n=1}^{\infty}a_{n}^{2}$, $\sum_{n=1}^{\infty}b_{n}^{2}$ 收敛, 证明: $\sum_{n=1}^{\infty}\left|a_{n}b_{n}\right|$, $\sum_{n=1}^{\infty}(a_{n}+b_{n})^{2}$ 也收敛.
  \item 设正项级数 $\sum_{n=1}^{\infty}a_{n}$ 收敛. 证明: 级数 $\sum_{n=1}^{\infty}\sqrt{a_{n}a_{n+1}}$ 也收敛. 举例说明逆命题不成立. 但若 $\{a_{n}\}$ 为单调减数列, 则逆命题也成立.
  \item 设 $\sum_{n=1}^{\infty}a_{n}$ 为一个收敛的正项级数. 证明: $\forall\delta>0$, 级数 $\sum_{n=1}^{\infty}n^{-\frac{1+\delta}{2}}\sqrt{a_{n}}<+\infty$. 问: 当 $\delta=0$ 时, 结论是否还成立?
  \item 设 $\delta>0$, $a_{n}>0$, $n=1,2,\cdots$. 证明:
  \begin{enumerate}
    \item[(1)] 如果当 $n>N$ 时, 有 $\left(\ln\frac{1}{a_{n}}\right)(\ln n)^{-1}\geqslant1+\delta$, 则 $\sum_{n=1}^{\infty}a_{n}$ 收敛;
    \item[(2)] 如果当 $n>N$ 时, 有 $\left(\ln\frac{1}{a_{n}}\right)(\ln n)^{-1}\leqslant1$, 则 $\sum_{n=1}^{\infty}a_{n}$ 发散;
    \item[(3)] 应用上面的结果, 证明 $\sum_{n=2}^{\infty}\frac{1}{(\ln n)^{\ln n}}$, $\sum_{n=1}^{\infty}\frac{1}{3^{\ln n}}$ 都收敛.
  \end{enumerate}
  \item 设 $p\geqslant0$. 证明: 级数
  \[
    \sum_{n=1}^{\infty}\frac{n!}{(p+1)(p+2)\cdots(p+n)}
\]
  当 $p>1$ 时收敛; 当 $0\leqslant p\leqslant1$ 时发散.
  \item 设 $a_{n}\geqslant0$, 证明: 级数 $\sum_{n=1}^{\infty}\frac{a_{n}}{(1+a_{1})(1+a_{2})\cdots(1+a_{n})}$ 收敛.
  \item 设 $a=A_{0}<A_{1}<A_{2}<\cdots<A_{n}<\cdots$, $A_{n}\to+\infty\ (n\to+\infty)$ 为任意给定的序列, $f(x)>0$, 则 $\int_{a}^{+\infty}f(x)\,\mathrm{d}x$ 与 $\sum_{n=1}^{\infty}\int_{A_{n-1}}^{A_{n}}f(x)\,\mathrm{d}x$ 同敛散, 且当收敛时, 有
  \[
    \int_{a}^{+\infty}f(x)\,\mathrm{d}x=\sum_{n=1}^{\infty}\int_{A_{n-1}}^{A_{n}}f(x)\,\mathrm{d}x.
\]
  \item 设 $a_{n}\neq0,\ n=1,2,\cdots$, 且 $\lim_{n\to+\infty}a_{n}=a\neq0$. 证明: 级数 $\sum_{n=1}^{\infty}\left|a_{n+1}-a_{n}\right|$ 与 $\sum_{n=1}^{\infty}\left|\frac{1}{a_{n+1}}-\frac{1}{a_{n}}\right|$ 同时敛散.
\end{enumerate}
\end{exblock}

\begin{exblock}{思考题 12.2}
\begin{enumerate}
  \item 证明: 级数 $\sum_{n=1}^{\infty}x^{1+\frac{1}{2}+\cdots+\frac{1}{n}}$ 当 $0<x<\mathrm{e}^{-1}$ 时收敛; 当 $x\geqslant\mathrm{e}^{-1}$ 时, 发散.
  \item 设 $a_{n}>0$, $S_{n}=\sum_{k=1}^{n}a_{k}$. 证明:
  \begin{enumerate}
    \item[(1)] 若 $\sum_{n=1}^{\infty}a_{n}$ 收敛, 则 $\sum_{n=1}^{\infty}\frac{a_{n}}{S_{n}}$ 也收敛;
    \item[(2)] 若 $\sum_{n=1}^{\infty}a_{n}$ 发散, 则 $\frac{a_{k}}{S_{n}^{2}}\leqslant\frac{1}{S_{k-1}}-\frac{1}{S_{k}}$ 且 $\sum_{n=1}^{\infty}\frac{a_{n}}{S_{n}^{2}}$ 收敛;
    \item[(3)] 若 $\sum_{n=1}^{\infty}a_{n}$ 发散, 则当 $\alpha>1$ 时, $\sum_{n=1}^{\infty}\frac{a_{n}}{S_{n}}$ 收敛; 当 $\alpha\leqslant1$ 时, $\sum_{n=1}^{\infty}\frac{a_{n}}{S_{n}^{\alpha}}$ 发散.
  \end{enumerate}
  \item 设 $\{a_{n}\}$ 为单调减的正数列. 证明: 级数 $\sum_{n=1}^{\infty}a_{n}$ 与 $\sum_{n=1}^{\infty}2^{n}a_{2^{n}}$ 同敛散. 应用上述结果证明:
  \begin{enumerate}
    \item[(1)] $\sum_{n=1}^{\infty}\frac{1}{n^{1+\alpha}}$ 当 $\alpha>0$ 时收敛, 当 $\alpha\leqslant0$ 时发散;
    \item[(2)] $\sum_{n=2}^{\infty}\frac{1}{n\ln n}$ 发散.
  \end{enumerate}
  \item 设 $\sum_{n=1}^{\infty}a_{n}$ 为一个发散的正项级数.
  \begin{enumerate}
    \item[(1)] 证明: $\sum_{n=1}^{\infty}\frac{a_{n}}{1+a_{n}}$ 发散, $\sum_{n=1}^{\infty}\frac{a_{n}}{1+n^{2}a_{n}}$ 收敛.
    \item[(2)] 研究 $\sum_{n=1}^{\infty}\frac{a_{n}}{1+a_{n}^{2}}$, $\sum_{n=1}^{\infty}\frac{a_{n}}{1+na_{n}}$ 的敛散性.
  \end{enumerate}
  \item 证明: 从调和级数 $\sum_{n=1}^{\infty}\frac{1}{n}$ 中划去所有分母中含有数字 9 的那些项(例如: $\frac{1}{9}$, $\frac{1}{19}$, $\frac{1}{29}$, $\cdots$)之后, 所得的新级数是收敛的, 且其和不超过 80.
  \item 设 $\sum_{n=1}^{\infty}a_{n}$ 为一个正项级数, 如果 $\lim_{n\to+\infty}\frac{a_{n+1}}{a_{n}}=q<1$ 或 $\lim_{n\to+\infty}\sqrt[n]{a_{n}}=q<1$, 则 $\forall r\in(q,1)$, 总有 $\lim_{n\to+\infty}\frac{a_{n}}{r^{n}}=0$. 这说明能用 d'Alembert 判别法的极限形式或 Cauchy 判别法的极限形式判定收敛的级数 $\sum_{n=1}^{\infty}a_{n}$ 总比等比(几何)级数 $\sum_{n=1}^{\infty}\frac{1}{r^{n}}$ 收敛得快.
  \item 研究下列级数的敛散性:
  \begin{enumerate}
    \item $\displaystyle\sum_{n=1}^{\infty}\left(\frac{1}{\sqrt{n}}-\sqrt{\ln\frac{n+1}{n}}\right)$;
    \item $\displaystyle\sum_{n=1}^{\infty}\left(a^{\frac{1}{n}}-\frac{b^{\frac{1}{n}}+c^{\frac{1}{n}}}{2}\right),\ a>0,\ b>0,\ c>0$;
    \item $\displaystyle\sum_{n=1}^{\infty}\left(n^{\frac{1}{n^{2}+1}}-1\right)$.
  \end{enumerate}
  \item 设 $a_{n}\geqslant0$, $n=1,2,\cdots$, $\sum_{n=1}^{\infty}a_{n}$ 收敛. 证明: $\lim_{m\to+\infty}\frac{1}{m}\sum_{n\leqslant m}1=0$, 其中 $\sum_{n\leqslant m}1$ 表示 $a_{1},\cdots,a_{m}$ 中满足 $a_{n}>\frac{1}{n}$ 的 $a_{n}$ 的个数.
  \item 设 $a_{n}>0\ (n=1,2,\cdots)$, $\sum_{n=1}^{\infty}a_{n}$ 收敛, 令 $r_{n}=\sum_{k=n}^{\infty}a_{k}$, 证明: 级数 $\sum_{k=1}^{\infty}\frac{a_{k}}{r_{k}}$ 发散.
  \item 设 $\lim_{n\to+\infty}\left(n^{2n\sin\frac{1}{n}}a_{n}\right)=1$, 证明: 级数 $\sum_{n=1}^{\infty}a_{n}$ 收敛.
  \item 设 $a_{n}=\left(l-\frac{p\ln n}{n}\right)^{n}$, 证明:
  \begin{enumerate}
    \item[(1)] $a_{n}\sim\frac{1}{n^{p}}\ (n\to+\infty)$;
    \item[(2)] 级数 $\sum_{n=1}^{\infty}a_{n}$ 当 $p>1$ 时收敛, 当 $p\leqslant1$ 时发散.
  \end{enumerate}
  \item 设 $0<p_{1}<p_{2}<\cdots<p_{n}<\cdots$, 证明:
  \[
    \sum_{n=1}^{\infty}\frac{1}{p_{n}}\text{ 收敛}\ \Leftrightarrow\
    \sum_{n=1}^{\infty}\frac{n}{p_{1}+p_{2}+\cdots+p_{n}}\text{ 收敛}.
\]
  \item 设 $a_{n}\geqslant0$, 且 $\sum_{n=1}^{\infty}a_{n}$ 收敛, 证明: 当 $p>\frac{1}{2}$ 时, 级数 $\sum_{n=1}^{\infty}\frac{\sqrt{a_{n}}}{n^{p}}$ 收敛.
  \item 设 $\sum_{n=1}^{\infty}a_{n}$ 为正项级数, 满足:
  \begin{enumerate}
    \item[(1)] $\sum_{k=1}^{n}(a_{k}-a_{n})$ 对 $n$ 有界;
    \item[(2)] 当 $n$ 趋于 $+\infty$ 时, $\{a_{n}\}$ 单调减趋于 $0$.
  \end{enumerate}
  证明: 级数 $\sum_{n=1}^{\infty}a_{n}$ 收敛.
  \item 应用 d'Alembert 判别法证明级数
  \[
    \sqrt{2-\sqrt{2}}+\sqrt{2-\sqrt{2}}\sqrt{2+\sqrt{2}}
    +\sqrt{2-\sqrt{2+\sqrt{2}}}\sqrt{2+\sqrt{2+\sqrt{2}}}+\cdots
\]
  收敛.
  \item 设级数 $\sum_{n=1}^{\infty}a_{n}$ 收敛, $a_{n}\geqslant0$, $\{a_{n}\}$ 单调减, 证明: $\lim_{n\to+\infty}na_{n}=0$, 即 $a_{n}=o\left(\frac{1}{n}\right)\ (n\to+\infty)$. 举例说明, 删去条件"$\{a_{n}\}$ 单调减", 上述结论不成立. 进而给出 $a_{n}>0\ (n=1,2,\cdots)$ 的反例.
  \item 设 $\{a_{n}\}$ 为正的单调减数列, $\sum_{n=1}^{\infty}a_{n}$ 发散. 证明:
  \[
    \lim_{n\to+\infty}\frac{a_{2}+a_{4}+\cdots+a_{2n}}{a_{1}+a_{3}+\cdots+a_{2n-1}}=1.
\]
  \item 设 $\{a_{n}\}$ 为正的单调增加的数列. 证明: 级数 $\sum_{n=1}^{\infty}\left(1-\frac{a_{n}}{a_{n+1}}\right)$ 当 $\{a_{n}\}$ 有界时收敛; 而当 $\{a_{n}\}$ 无界时发散.
\end{enumerate}
\end{exblock}

% ==================================================================
%  徐森林、金亚东、薛春华《数学分析》（第三册）
%  S12.3  一般级数
%  源：原书印刷页 37-53（PDF 物理页 43-59）
%  底本：PaddleOCR-VL-1.6，见 chapters/_ch12_12_3_raw.md
%  约定：全西文标点；行间公式用一对方括号；多行公式用 align/aligned；
%        跨册引用一律 \extref；数学模式内的中文必须包 \text{}
%  已订正的 OCR 错误见 tools/_make_ch12_3a.py 头部说明
% ==================================================================
\section{一般级数}\label{sec:12.3}

如果级数 $\sum_{n=1}^{\infty}a_{n}$ 的通项 $a_{n}$ 可正可负时, 称它为一般级数. 例如: $\sum_{n=1}^{\infty}\frac{\sin nx}{n}$, $\sum_{n=1}^{\infty}\frac{(-1)^{n-1}}{n}=1-\frac{1}{2}+\frac{1}{3}-\cdots+(-1)^{n-1}\frac{1}{n}+\cdots$ 都为一般级数.

对于一般级数, 可以用级数收敛的定义及级数的 Cauchy 收敛判别法（收敛准则）来判定其敛散性.

这一节我们将给出一般级数的几种特殊判别法: Leibniz 判别法, Dirichlet 判别法, Abel 判别法等.

\begin{theorem}[Leibniz 判别法]\label{thm:12.3.1}
如果 $\{a_{n}\}$ 单调减趋于 $0$, 则交错级数 $\sum_{n=1}^{\infty}(-1)^{n-1}a_{n}$（$a_{n}\geqslant0$）收敛.

设 $S_{n}$ 为级数 $\sum_{n=1}^{\infty}(-1)^{n-1}a_{n}$ 的第 $n$ 个部分和, 则 $\{S_{2k}\}$ 单调增趋于 $S$, $\{S_{2k-1}\}$ 单调减趋于 $S$. 于是, 用 $S_{n}$ 作为 $S=\sum_{n=1}^{\infty}(-1)^{n-1}a_{n}$ 的近似值, 其绝对误差为
\[
  |S_{n}-S|\leqslant a_{n+1}.
\]
\end{theorem}

\begin{proof}
因为 $\{a_{n}\}$ 单调减趋于 $0$, 故 $a_{n}-a_{n+1}\geqslant0$, 从而
\[
  S_{2k}=(a_{1}-a_{2})+(a_{3}-a_{4})+\cdots+(a_{2k-3}-a_{2k-2})+(a_{2k-1}-a_{2k})
  =S_{2k-2}+(a_{2k-1}-a_{2k})\geqslant S_{2k-2}.
\]
这就表明 $\{S_{2k}\}$ 单调增.

另一方面, 由
\[
  S_{2k}=a_{1}-(a_{2}-a_{3})-\cdots-(a_{2k-2}-a_{2k-1})-a_{2k}<a_{1}
\]
知 $\{S_{2k}\}$ 有上界 $a_{1}$. 因此, $\{S_{2k}\}$ 收敛于实数 $S$. 又
\[
  \lim_{k\to+\infty}S_{2k+1}=\lim_{k\to+\infty}(S_{2k}+a_{2k+1})=S+0=S.
\]
根据\extref{定理 1.1.1}, $\{S_{n}\}$ 收敛, 即 $\sum_{n=1}^{\infty}(-1)^{n-1}a_{n}$ 收敛, 记 $\sum_{n=1}^{\infty}(-1)^{n-1}a_{n}=\lim\limits_{n\to+\infty}S_{n}=S$.

从
\[
  S_{2k-1}=S_{2k-3}-(a_{2k-2}-a_{2k-1})\leqslant S_{2k-3}
\]
知 $\{S_{2k-1}\}$ 单调减, 由上段的结论可知它趋于 $S$. 于是,
\[
  |S_{n}-S|\leqslant|S_{n}-S_{n+1}|=a_{n+1}.
\]
\end{proof}

$\{a_{n}\}$ 单调减趋于 $0$ 的交错级数 $\sum_{n=1}^{\infty}(-1)^{n-1}a_{n}$ 常称为 Leibniz 级数.

为给出 Dirichlet 判别法与 Abel 判别法, 我们先证两个引理.

\begin{lemma}[分部求和法]\label{lem:12.3.1}
设 $S_{k}=a_{1}+\cdots+a_{k}\ (k=1,2,\cdots)$, 则有
\[
  \sum_{k=1}^{n}a_{k}b_{k}=\sum_{k=1}^{n-1}S_{k}(b_{k}-b_{k+1})+S_{n}b_{n}.
\]
\end{lemma}

\begin{proof}
设 $S_{0}=0$, 则 $a_{k}=S_{k}-S_{k-1}$, $k=1,2,\cdots$. 于是,
\[
\begin{aligned}
  \sum_{k=1}^{n}a_{k}b_{k}
  &=\sum_{k=1}^{n}(S_{k}-S_{k-1})b_{k}
   =\sum_{k=1}^{n}S_{k}b_{k}-\sum_{k=1}^{n}S_{k-1}b_{k}\\
  &=\sum_{k=1}^{n-1}S_{k}b_{k}-\sum_{k=1}^{n-1}S_{k}b_{k+1}+S_{n}b_{n}
   =\sum_{k=1}^{n-1}S_{k}(b_{k}-b_{k+1})+S_{n}b_{n}.
\end{aligned}
\]
\end{proof}

\begin{lemma}[Abel 引理]\label{lem:12.3.2}
(1) 设 $b_{1}\geqslant b_{2}\geqslant\cdots\geqslant b_{n}\geqslant0$,
\[
  m\leqslant S_{k}=\sum_{l=1}^{k}a_{l}\leqslant M\quad(k=1,2,\cdots,n),
\]
则有
\[
  b_{1}m\leqslant\sum_{k=1}^{n}a_{k}b_{k}\leqslant b_{1}M.
\]

(2) 设 $b_{1}\geqslant b_{2}\geqslant\cdots\geqslant b_{n}$ 或 $b_{1}\leqslant b_{2}\leqslant\cdots\leqslant b_{n}$, 记 $S_{k}=\sum_{l=1}^{k}a_{l}$, 如果 $|S_{k}|\leqslant M$, $k=1,\cdots,n$, 则
\[
  \Bigl|\sum_{k=1}^{n}a_{k}b_{k}\Bigr|\leqslant M\bigl(|b_{1}|+2|b_{n}|\bigr).
\]
\end{lemma}

\begin{proof}
(1) 由引理 12.3.1 及 $b_{1}\geqslant b_{2}\geqslant\cdots\geqslant b_{n}\geqslant0$, $b_{k}-b_{k+1}\geqslant0\ (k=1,2,\cdots,n-1)$, 有
\[
  \sum_{k=1}^{n}a_{k}b_{k}
  =\sum_{k=1}^{n-1}S_{k}(b_{k}-b_{k+1})+S_{n}b_{n}
  \leqslant M\sum_{k=1}^{n-1}(b_{k}-b_{k+1})+Mb_{n}=Mb_{1}.
\]
同理,
\[
  mb_{1}\leqslant\sum_{k=1}^{n}a_{k}b_{k}.
\]
结合上面的不等式得到
\[
  mb_{1}\leqslant\sum_{k=1}^{n}a_{k}b_{k}\leqslant Mb_{1}.
\]

(2) 由分部求和公式得
\[
\begin{aligned}
  \Bigl|\sum_{k=1}^{n}a_{k}b_{k}\Bigr|
  &=\Bigl|\sum_{k=1}^{n-1}S_{k}(b_{k}-b_{k+1})+S_{n}b_{n}\Bigr|
   \leqslant\sum_{k=1}^{n-1}|S_{k}|\,|b_{k}-b_{k+1}|+|S_{n}|\,|b_{n}|\\
  &\leqslant M\Bigl(\sum_{k=1}^{n-1}|b_{k}-b_{k+1}|+|b_{n}|\Bigr)
   =M\bigl(|b_{1}-b_{n}|+|b_{n}|\bigr)
   \leqslant M\bigl(|b_{1}|+2|b_{n}|\bigr).
\end{aligned}
\]
\end{proof}

从 Abel 引理立即有下面的结论.

\begin{theorem}[Dirichlet 判别法]\label{thm:12.3.2}
设
(1) $\{b_{n}\}$ 单调趋于 $0$;
(2) $S_{k}=a_{1}+\cdots+a_{k}$, $|S_{k}|\leqslant M$, $k=1,2,\cdots$, 即 $\sum_{n=1}^{\infty}a_{n}$ 的部分和有界.
则 $\sum_{n=1}^{\infty}a_{n}b_{n}$ 收敛.
\end{theorem}

\begin{proof}
证法 1 由条件 (2) 知, $-M\leqslant S_{k}\leqslant M$, $k=1,2,\cdots$, 于是
\[
  -M-S_{n}\leqslant S_{n+p}-S_{n}=a_{n+1}+\cdots+a_{n+p}\leqslant M-S_{n},
  \quad\forall p\in\mathbb{N}.
\]
由条件 (1), 且不妨设 $\{b_{n}\}$ 单调减（$\{b_{n}\}$ 单调增时用 $\{-b_{n}\}$ 代替 $\{b_{n}\}$）趋于 $0$, 即 $b_{n}\geqslant0$. 根据 Abel 引理 \ref{lem:12.3.2}(1), 有
\[
  (-M-S_{n})b_{n+1}\leqslant a_{n+1}b_{n+1}+\cdots+a_{n+p}b_{n+p}\leqslant(M-S_{n})b_{n+1},
\]
所以,
\[
  \Bigl|\sum_{k=n+1}^{n+p}a_{k}b_{k}\Bigr|
  \leqslant\max\{|-M-S_{n}|,\,|M-S_{n}|\}\,b_{n+1}\leqslant 2Mb_{n+1}.
\]
因为 $\{b_{n}\}$ 单调趋于 $0$, 所以 $\forall\varepsilon>0$, $\exists N\in\mathbb{N}$, 当 $n>N$ 时, $b_{n+1}<\frac{\varepsilon}{2M+1}$, 于是
\[
  \Bigl|\sum_{k=n+1}^{n+p}a_{k}b_{k}\Bigr|<2M\cdot\frac{\varepsilon}{2M+1}<\varepsilon.
\]
根据级数的 Cauchy 收敛准则（定理~\ref{thm:12.1.1}）, $\sum_{n=1}^{\infty}a_{n}b_{n}$ 收敛.

证法 2 由条件 (2) 知, $|S_{k}|\leqslant M$, $k=1,2,\cdots$. 再由 $\{b_{n}\}$ 单调减及 Abel 引理 \ref{lem:12.3.2}(2) 得到
\[
  \Bigl|\sum_{k=n+1}^{n+p}a_{k}\Bigr|=|S_{n+p}-S_{n}|\leqslant|S_{n+p}|+|S_{n}|\leqslant 2M,
\]
\[
  \Bigl|\sum_{k=n+1}^{n+p}a_{k}b_{k}\Bigr|\leqslant 2M\bigl(|b_{n+1}|+2|b_{n+p}|\bigr).
\]
因为 $\{b_{n}\}$ 趋于 $0$, 所以 $\forall\varepsilon>0$, $\exists N=N(\varepsilon)\in\mathbb{N}$, 当 $n>N$ 时, $|b_{n}|<\frac{\varepsilon}{6M+1}$. 于是
\[
  \Bigl|\sum_{k=n+1}^{n+p}a_{k}b_{k}\Bigr|
  \leqslant 2M\bigl(|b_{n+1}|+2|b_{n+p}|\bigr)
  <2M\cdot\frac{3\varepsilon}{6M+1}<\varepsilon.
\]
根据级数的 Cauchy 收敛准则（定理~\ref{thm:12.1.1}）, $\sum_{n=1}^{\infty}a_{n}b_{n}$ 收敛.
\end{proof}

\begin{theorem}[Abel 判别法]\label{thm:12.3.3}
设
(1) $\{b_{n}\}$ 单调有界;
(2) $\sum_{n=1}^{\infty}a_{n}$ 收敛.
则 $\sum_{n=1}^{\infty}a_{n}b_{n}$ 收敛.
\end{theorem}

\begin{proof}
证法 1 由于 $\{b_{n}\}$ 单调有界, 故 $\{b_{n}\}$ 收敛, 记 $\lim\limits_{n\to+\infty}b_{n}=b$. 不妨设 $\{b_{n}\}$ 单调减, 则 $b_{n}\geqslant b$. 令 $b_{n}'=b_{n}-b$, 则 $\{b_{n}'\}$ 单调减趋于 $0$.

又因 $\sum_{n=1}^{\infty}a_{n}$ 收敛, 故部分和有界, 根据 Dirichlet 判别法（定理~\ref{thm:12.3.2}）知 $\sum_{n=1}^{\infty}a_{n}b_{n}'$ 收敛. 从而,
\[
  \sum_{n=1}^{\infty}a_{n}b_{n}
  =\sum_{n=1}^{\infty}a_{n}(b_{n}'+b)
  =\sum_{n=1}^{\infty}a_{n}b_{n}'+b\sum_{n=1}^{\infty}a_{n}
\]
也收敛.

证法 2 由条件 $|b_{n}|\leqslant B$, $\forall n\in\mathbb{N}$. 又因 $\sum_{n=1}^{\infty}a_{n}$ 收敛, 故根据级数的 Cauchy 收敛准则（定理~\ref{thm:12.1.1}）, $\forall\varepsilon>0$, $\exists N\in\mathbb{N}$, 当 $n>N$ 时, 有
\[
  |\widetilde{S}_{p}|<\frac{\varepsilon}{2B+1},\quad\forall p\in\mathbb{N},
\]
其中 $\widetilde{S}_{p}=a_{n+1}+\cdots+a_{n+p}$, 则 $\{\widetilde{S}_{p}\}\ (p\in\mathbb{N})$ 有界, 且
\[
  -\frac{\varepsilon}{2B+1}<\widetilde{S}_{p}=a_{n+1}+\cdots+a_{n+p}<\frac{\varepsilon}{2B+1}.
\]
又因 $\{b_{n}\}$ 单调有界. 不妨设 $\{b_{n}\}$ 单调减（否则用 $\{-b_{n}\}$ 代替）, 且以 $0$ 为下界（否则用 $\{b_{n}+B\}$ 代替）. 根据 Abel 引理 \ref{lem:12.3.2}(1) 得到
\[
  -\frac{\varepsilon}{2B+1}b_{n+1}\leqslant a_{n+1}b_{n+1}+\cdots+a_{n+p}b_{n+p}
  \leqslant\frac{\varepsilon}{2B+1}b_{n+1},
\]
所以,
\[
  \Bigl|\sum_{k=n+1}^{n+p}a_{k}b_{k}\Bigr|
  \leqslant\frac{\varepsilon}{2B+1}\cdot b_{n+1}
  \leqslant\frac{\varepsilon}{2B+1}\cdot B<\varepsilon.
\]
根据级数的 Cauchy 收敛准则（定理~\ref{thm:12.1.1}）, $\sum_{n=1}^{\infty}a_{n}b_{n}$ 收敛.

证法 3 由条件 $|b_{n}|\leqslant B$, $\forall n\in\mathbb{N}$. 又因 $\sum_{n=1}^{\infty}a_{n}$ 收敛, 故 $\forall\varepsilon>0$, 根据级数的 Cauchy 收敛准则（定理~\ref{thm:12.1.1}）, $\exists N\in\mathbb{N}$, 当 $n>N$ 时,
\[
  |a_{n+1}+\cdots+a_{n+l}|<\frac{\varepsilon}{3B+1},\quad\forall l\in\mathbb{N}.
\]
根据 Abel 引理 \ref{lem:12.3.2}(2) 得到
\[
\begin{aligned}
  \Bigl|\sum_{k=n+1}^{n+p}a_{k}b_{k}\Bigr|
  &\leqslant\frac{\varepsilon}{3B+1}\bigl(|b_{n+1}|+2|b_{n+p}|\bigr)
   \leqslant\frac{\varepsilon}{3B+1}\cdot 3B<\varepsilon.
\end{aligned}
\]
再根据级数的 Cauchy 收敛准则（定理~\ref{thm:12.1.1}）, $\sum_{n=1}^{\infty}a_{n}b_{n}$ 收敛.
\end{proof}

\begin{remark}\label{rem:12.3.1}
Leibniz 判别法（定理~\ref{thm:12.3.1}）是 Dirichlet 判别法（定理~\ref{thm:12.3.2}）的特殊情形. 事实上, 由于 $\sum_{n=1}^{\infty}(-1)^{n-1}$ 的部分和 $\sum_{k=1}^{n}(-1)^{k-1}$ 满足 $\bigl|\sum_{k=1}^{n}(-1)^{k-1}\bigr|\leqslant2$, 即部分和有界. 又因为 $\{a_{n}\}$ 单调减趋于 $0$, 根据 Dirichlet 判别法知, $\sum_{n=1}^{\infty}(-1)^{n-1}a_{n}$ 收敛.
\end{remark}

\begin{example}\label{ex:12.3.1}
证明: $\sum_{n=1}^{\infty}\frac{(-1)^{n-1}}{n}$ 与 $\sum_{n=2}^{\infty}\frac{(-1)^{n}}{\ln n}$ 收敛.
\end{example}

\begin{proof}
因为 $\bigl\{\frac{1}{n}\bigr\}$ 与 $\bigl\{\frac{1}{\ln n}\bigr\}$ 都单调减趋于 $0$, 所以根据 Leibniz 判别法知, $\sum_{n=1}^{\infty}\frac{(-1)^{n-1}}{n}$ 与 $\sum_{n=2}^{\infty}\frac{(-1)^{n}}{\ln n}$ 都收敛.
\end{proof}

\begin{example}\label{ex:12.3.2}
研究 $\sum_{n=1}^{\infty}\frac{\cos nx}{n}$ 的敛散性（此级数既非正项级数又非 Leibniz 级数）.
\end{example}

\begin{proof}
解 设 $a_{n}=\cos nx$, $b_{n}=\frac{1}{n}$, $n=1,2,\cdots$. 显然, $\{b_{n}\}$ 单调减趋于 $0$, 而
\[
\begin{aligned}
  |S_{n}|
  &=|\cos x+\cos 2x+\cdots+\cos nx|
   =\Bigl|\frac{1}{2\sin\frac{x}{2}}\sum_{k=1}^{n}2\sin\frac{x}{2}\cos kx\Bigr|\\
  &=\Bigl|\frac{1}{2\sin\frac{x}{2}}\sum_{k=1}^{n}\bigl[\sin\bigl(k+\tfrac{1}{2}\bigr)x-\sin\bigl(k-\tfrac{1}{2}\bigr)x\bigr]\Bigr|
   =\Bigl|\frac{\sin\bigl(n+\tfrac{1}{2}\bigr)x-\sin\frac{x}{2}}{2\sin\frac{x}{2}}\Bigr|
   \leqslant\Bigl|\frac{1}{\sin\frac{x}{2}}\Bigr|,
\end{aligned}
\]
其中 $x\neq 2m\pi$, $m\in\mathbb{Z}$. 此时, $S_{n}=\sum_{k=1}^{n}a_{k}$ 有界. 根据 Dirichlet 判别法（定理~\ref{thm:12.3.2}）, 当 $x\neq 2m\pi\ (m\in\mathbb{Z})$ 时收敛; 当 $x=2m\pi\ (m\in\mathbb{Z})$ 时, $\sum_{n=1}^{\infty}\frac{\cos n(2m\pi)}{n}=\sum_{n=1}^{\infty}\frac{1}{n}$ 发散.
\end{proof}

\begin{example}\label{ex:12.3.3}
研究级数 $\sum_{n=1}^{\infty}\frac{\cos 3n}{n}\bigl(1+\frac{1}{n}\bigr)^{n}$ 的敛散性.
\end{example}

\begin{proof}
解 由例 \ref{ex:12.3.2}, $\sum_{n=1}^{\infty}\frac{\cos 3n}{n}$ 收敛, $\bigl(1+\frac{1}{n}\bigr)^{n}$ 单调增有上界 $3$（或 $\mathrm{e}$）. 根据 Abel 判别法知, $\sum_{n=1}^{\infty}\frac{\cos 3n}{n}\bigl(1+\frac{1}{n}\bigr)^{n}$ 收敛.
\end{proof}

\begin{example}\label{ex:12.3.4}
研究级数 $\sum_{n=1}^{\infty}(-1)^{n}\frac{\sin^{2}n}{n}$ 的敛散性.
\end{example}

\begin{proof}
解 由 Leibniz 判别法知, $\sum_{n=1}^{\infty}\frac{(-1)^{n}}{n}$ 收敛. 再根据例 \ref{ex:12.3.2},
\[
  \sum_{n=1}^{\infty}(-1)^{n}\frac{\cos 2n}{n}
  =\sum_{n=1}^{\infty}\frac{\cos n(2+\pi)}{n}\ \text{收敛}
  \quad(\text{因}\ 2+\pi\neq 2m\pi,\ m\in\mathbb{Z}).
\]
于是,
\[
  \sum_{n=1}^{\infty}(-1)^{n}\frac{\sin^{2}n}{n}
  =\sum_{n=1}^{\infty}(-1)^{n}\frac{1-\cos 2n}{2n}
  =\frac{1}{2}\biggl[\sum_{n=1}^{\infty}\frac{(-1)^{n}}{n}
  -\sum_{n=1}^{\infty}\frac{(-1)^{n}}{n}\cos 2n\biggr]
\]
是收敛的.

但是,
\[
  \sum_{n=1}^{\infty}\Bigl|(-1)^{n}\frac{\sin^{2}n}{n}\Bigr|
  =\sum_{n=1}^{\infty}\frac{\sin^{2}n}{n}
  =\frac{1}{2}\sum_{n=1}^{\infty}\frac{1-\cos 2n}{n}
\]
是发散的.

（反证）假设 $\sum_{n=1}^{\infty}\frac{1-\cos 2n}{n}$ 是收敛的, 则由例 \ref{ex:12.3.2} 及定理~\ref{thm:12.1.3} 推得
\[
  \sum_{n=1}^{\infty}\frac{1}{n}
  =\sum_{n=1}^{\infty}\biggl(\frac{1-\cos 2n}{n}+\frac{\cos 2n}{n}\biggr)
\]
也收敛. 这与例 \ref{ex:12.2.1} 中 $\sum_{n=1}^{\infty}\frac{1}{n}$ 发散相矛盾.
\end{proof}

\begin{example}\label{ex:12.3.5}
(1) 设 $\{a_{n}\}$ 为单调减的正数列, 如果 $\sum_{n=1}^{\infty}a_{n}$ 收敛, 则 $a_{n}=o\bigl(\frac{1}{n}\bigr)$;

(2) 举出例子: $a_{n}\geqslant0$, $\{a_{n}\}$ 单调减且 $a_{n}=o\bigl(\frac{1}{n}\bigr)$, $n\to+\infty$, 但 $\sum_{n=1}^{\infty}a_{n}$ 发散;

(3) 举出例子: $a_{n}\geqslant 0$, $\sum_{n=1}^{\infty}a_{n}$ 收敛, 但 $a_{n}\neq o\bigl(\frac{1}{n}\bigr)$, $n\to+\infty$.
\end{example}

\begin{proof}
证明 (1) 因为 $\sum_{n=1}^{\infty}a_{n}$ 收敛, 根据级数的 Cauchy 收敛准则, $\forall\varepsilon>0$, $\exists N\in\mathbb{N}$, 当 $n>N$ 时,
\[
  |a_{n+1}+\cdots+a_{n+p}|<\frac{\varepsilon}{2},\quad\forall p\in\mathbb{N}.
\]
特别取 $p=n$, 并由 $\{a_{n}\}$ 单调减得到
\[
  na_{2n}\leqslant a_{n+1}+\cdots+a_{2n}<\frac{\varepsilon}{2},
\]
\[
  0\leqslant 2na_{2n}<2\cdot\frac{\varepsilon}{2}=\varepsilon,
\]
\[
  \lim_{n\to+\infty}2na_{2n}=0.
\]
又因为
\[
  0\leqslant(2n+1)a_{2n+1}\leqslant 2na_{2n}+a_{2n+1}\to 0+0=0\quad(n\to+\infty),
\]
所以
\[
  \lim_{n\to+\infty}(2n+1)a_{2n+1}=0.
\]
根据\extref{定理 1.1.1} 得到
\[
  \lim_{n\to+\infty}\frac{a_{n}}{1/n}=\lim_{n\to+\infty}na_{n}=0,
  \quad\text{即}\ a_{n}=o\Bigl(\frac{1}{n}\Bigr),\ n\to+\infty.
\]

(2) 显然,
\[
  a_{n}=\frac{1}{(n+1)\ln(n+1)}=o\Bigl(\frac{1}{n}\Bigr)\quad(n\to+\infty),
\]
且 $\{a_{n}\}$ 单调减, $a_{n}\geqslant0$. 但是, $\sum_{n=1}^{\infty}\frac{1}{(n+1)\ln(n+1)}$ 发散.

(3) 设
\[
  a_{n}=\left\{\begin{array}
{ll}
    0, & n\neq m^{2},\\[2pt]
    \dfrac{1}{m^{2}}, & n=m^{2},
\end{array}\right.\quad m=1,2,\cdots
\]
或
\[
  a_{n}=\left\{\begin{array}
{ll}
    \displaystyle\frac{1}{2^{n}}, & n\neq m^{2},\\[6pt]
    \displaystyle\frac{1}{m^{2}}, & n=m^{2},
\end{array}\right.\quad m=1,2,\cdots.
\]
则 $a_{n}\geqslant0$（后者 $a_{n}>0$）, $\sum_{n=1}^{\infty}a_{n}$ 收敛, 但 $a_{n}\neq o\bigl(\frac{1}{n}\bigr)$, $n\to+\infty$.
\end{proof}
% ===== PART-B =====

% ===== PART-B =====
% ---------- 绝对收敛与条件收敛 ----------

\begin{theorem}\label{thm:12.3.4}
$\sum_{n=1}^{\infty}|a_{n}|$ 收敛蕴涵 $\sum_{n=1}^{\infty}a_{n}$ 收敛. 反之不真.
\end{theorem}

\begin{proof}
因为 $\sum_{n=1}^{\infty}|a_{n}|$ 收敛, 故由级数的 Cauchy 收敛准则, $\forall\varepsilon>0$, $\exists N\in\mathbb{N}$, 当 $n>N$ 时, 有
\[
  |a_{n+1}|+\cdots+|a_{n+p}|<\varepsilon,\quad\forall p\in\mathbb{N},
\]
则
\[
  |a_{n+1}+\cdots+a_{n+p}|\leqslant|a_{n+1}|+\cdots+|a_{n+p}|<\varepsilon.
\]
再由级数的 Cauchy 收敛准则, $\sum_{n=1}^{\infty}a_{n}$ 收敛.

但反之不真. 反例: $\sum_{n=1}^{\infty}\frac{(-1)^{n}}{n}$.
\end{proof}

\begin{definition}\label{def:12.3.1}
如果 $\sum_{n=1}^{\infty}|a_{n}|$ 收敛, 则称级数 $\sum_{n=1}^{\infty}a_{n}$ 绝对收敛. 根据定理~\ref{thm:12.3.4}, 绝对收敛的级数 $\sum_{n=1}^{\infty}a_{n}$, 它本身必收敛. 如果 $\sum_{n=1}^{\infty}a_{n}$ 收敛, 而 $\sum_{n=1}^{\infty}|a_{n}|$ 发散, 则称级数 $\sum_{n=1}^{\infty}a_{n}$ 是条件收敛的.

例如: $\sum_{n=1}^{\infty}\frac{(-1)^{[\sqrt{n}]}}{n^{2}}$, $\sum_{n=1}^{\infty}\frac{\cos nx}{n^{2}}$ 都是绝对收敛的. 而 $\sum_{n=1}^{\infty}\frac{(-1)^{n-1}}{n}$, $\sum_{n=1}^{\infty}\frac{(-1)^{n}}{n}\sin^{2}n$（例 \ref{ex:12.3.4}）都是条件收敛的.
\end{definition}

\begin{definition}\label{def:12.3.2}
设 $\sum_{n=1}^{\infty}a_{n}$ 为一般级数, 令
\[
  a_{n}^{+}=\left\{\begin{aligned}
    &a_{n}, & a_{n}\geqslant0,\\
    &0, & a_{n}<0,
  \end{aligned}\right.
  \qquad
  a_{n}^{-}=\left\{\begin{aligned}
    &0, & a_{n}\geqslant0,\\
    &-a_{n}, & a_{n}<0,
  \end{aligned}\right.
\]
则 $\sum_{n=1}^{\infty}a_{n}^{+}$, $\sum_{n=1}^{\infty}a_{n}^{-}$ 都是正项级数, 且
\[
  \sum_{n=1}^{\infty}a_{n}=\sum_{n=1}^{\infty}(a_{n}^{+}-a_{n}^{-}),
  \qquad
  \sum_{n=1}^{\infty}|a_{n}|=\sum_{n=1}^{\infty}(a_{n}^{+}+a_{n}^{-}).
\]
\end{definition}

\begin{theorem}\label{thm:12.3.5}
(1) $\sum_{n=1}^{\infty}a_{n}$ 绝对收敛 $\Leftrightarrow$ $\sum_{n=1}^{\infty}a_{n}^{+}$ 与 $\sum_{n=1}^{\infty}a_{n}^{-}$ 都收敛;

(2) $\sum_{n=1}^{\infty}a_{n}$ 条件收敛 $\Rightarrow$ $\sum_{n=1}^{\infty}a_{n}^{+}=\sum_{n=1}^{\infty}a_{n}^{-}=+\infty$, 但反之不真.
\end{theorem}

\begin{proof}
(1) 易见, $|a_{n}|=a_{n}^{+}+a_{n}^{-}$, $a_{n}^{+}\leqslant|a_{n}|$, $a_{n}^{-}\leqslant|a_{n}|$,
\[
  \sum_{n=1}^{N}|a_{n}|=\sum_{n=1}^{N}a_{n}^{+}+\sum_{n=1}^{N}a_{n}^{-}.
\]
因此,
$\sum_{n=1}^{N}|a_{n}|$ 有上界 $\Leftrightarrow$ $\sum_{n=1}^{N}a_{n}^{+}$ 与 $\sum_{n=1}^{N}a_{n}^{-}$ 同时有上界.

即
$\sum_{n=1}^{\infty}|a_{n}|$ 收敛 $\Leftrightarrow$ $\sum_{n=1}^{\infty}a_{n}^{+}$ 与 $\sum_{n=1}^{\infty}a_{n}^{-}$ 同时收敛.

此时,
\[
  \sum_{n=1}^{\infty}|a_{n}|=\sum_{n=1}^{\infty}(a_{n}^{+}+a_{n}^{-})
  =\sum_{n=1}^{\infty}a_{n}^{+}+\sum_{n=1}^{\infty}a_{n}^{-},
\]
\[
  \sum_{n=1}^{\infty}a_{n}=\sum_{n=1}^{\infty}(a_{n}^{+}-a_{n}^{-})
  =\sum_{n=1}^{\infty}a_{n}^{+}-\sum_{n=1}^{\infty}a_{n}^{-}.
\]

(2) 因为 $\sum_{n=1}^{\infty}a_{n}$ 条件收敛, 即 $\sum_{n=1}^{\infty}a_{n}$ 收敛, 但 $\sum_{n=1}^{\infty}|a_{n}|$ 发散. 根据 (1), $\sum_{n=1}^{\infty}a_{n}^{+}$ 与 $\sum_{n=1}^{\infty}a_{n}^{-}$ 至少有一个发散. （反证）假设有一个收敛, 一个发散, 则
\[
  \sum_{n=1}^{\infty}a_{n}=\sum_{n=1}^{\infty}a_{n}^{+}-\sum_{n=1}^{\infty}a_{n}^{-}
\]
发散, 这与已知 $\sum_{n=1}^{\infty}a_{n}$ 收敛相矛盾. 这就证明了 $\sum_{n=1}^{\infty}a_{n}^{+}$ 与 $\sum_{n=1}^{\infty}a_{n}^{-}$ 都发散, 即
\[
  \sum_{n=1}^{\infty}a_{n}^{+}=\sum_{n=1}^{\infty}a_{n}^{-}=+\infty.
\]
但反之不真. 例如: $a_{n}=(-1)^{n}$,
\[
  a_{n}^{+}=\left\{\begin{matrix}
    1, & n=2k,\\
    0, & n=2k-1,
  \end{matrix}\right.
  \qquad
  a_{n}^{-}=\left\{\begin{matrix}
    0, & n=2k,\\
    1, & n=2k-1.
  \end{matrix}\right.
\]
$\sum_{n=1}^{\infty}a_{n}^{+}=\sum_{n=1}^{\infty}a_{n}^{-}=+\infty$, 而 $\sum_{n=1}^{\infty}a_{n}$ 不收敛, 当然它也不是条件收敛的.
\end{proof}

\begin{theorem}[绝对收敛级数的特性]\label{thm:12.3.6}
绝对收敛级数中交换任意（或无穷）多项的次序所得的新级数仍然绝对收敛, 其和不变.
\end{theorem}

\begin{proof}
第 1 步, 先设 $\sum_{n=1}^{\infty}a_{n}$ 为收敛的正项级数. 交换任意项的次序后所得新级数为 $\sum_{n=1}^{\infty}b_{n}$. 因此, 新级数的部分和 $\sum_{n=1}^{N}b_{n}$ 是从 $\sum_{n=1}^{\infty}a_{n}$ 中挑出来的有限项构成的和. 于是, 有
\[
  \sum_{n=1}^{N}b_{n}\leqslant\sum_{n=1}^{\infty}a_{n}<+\infty,
\]
即 $\sum_{n=1}^{\infty}b_{n}$ 的部分和有界, 故由定理~\ref{thm:12.2.1} 知, $\sum_{n=1}^{\infty}b_{n}$ 收敛, 且 $\sum_{n=1}^{\infty}b_{n}\leqslant\sum_{n=1}^{\infty}a_{n}$. 反之, $\sum_{n=1}^{\infty}a_{n}$ 也可看作 $\sum_{n=1}^{\infty}b_{n}$ 交换任意项的次序得到的新级数. 于是, 又有
\[
  \sum_{n=1}^{\infty}a_{n}\leqslant\sum_{n=1}^{\infty}b_{n}.
\]
综上得到 $\sum_{n=1}^{\infty}a_{n}=\sum_{n=1}^{\infty}b_{n}$.

第 2 步, 再考虑一般的绝对收敛级数 $\sum_{n=1}^{\infty}a_{n}$. 由定理~\ref{thm:12.3.5}(1),
\[
  \sum_{n=1}^{\infty}a_{n}=\sum_{n=1}^{\infty}a_{n}^{+}-\sum_{n=1}^{\infty}a_{n}^{-}.
\]
设交换 $\sum_{n=1}^{\infty}a_{n}$ 的任意项次序所得新级数为 $\sum_{n=1}^{\infty}b_{n}$. 显然, $\sum_{n=1}^{\infty}b_{n}^{+}$, $\sum_{n=1}^{\infty}b_{n}^{-}$ 分别是由 $\sum_{n=1}^{\infty}a_{n}^{+}$, $\sum_{n=1}^{\infty}a_{n}^{-}$ 交换任意项次序所得的新级数. 由于 $\sum_{n=1}^{\infty}a_{n}^{+}$, $\sum_{n=1}^{\infty}a_{n}^{-}$ 都是正项级数, 根据第 1 步知, $\sum_{n=1}^{\infty}b_{n}^{+}=\sum_{n=1}^{\infty}a_{n}^{+}$, $\sum_{n=1}^{\infty}b_{n}^{-}=\sum_{n=1}^{\infty}a_{n}^{-}$, 并且都是收敛的. 再由定理~\ref{thm:12.3.5}(1),
\[
  \sum_{n=1}^{\infty}b_{n}=\sum_{n=1}^{\infty}b_{n}^{+}-\sum_{n=1}^{\infty}b_{n}^{-}
  =\sum_{n=1}^{\infty}a_{n}^{+}-\sum_{n=1}^{\infty}a_{n}^{-}=\sum_{n=1}^{\infty}a_{n}
\]
收敛, 且其和不变.
\end{proof}

\begin{example}\label{ex:12.3.6}
举例说明: 对于条件收敛级数, 交换它的任意多项的次序, 定理~\ref{thm:12.3.6} 中的结论不成立.
\end{example}

\begin{proof}
解 反例: $\sum_{n=1}^{\infty}\frac{(-1)^{n-1}}{n}$ 是条件收敛的, $a_{n}=\frac{(-1)^{n-1}}{n}$,
\[
  a_{n}^{+}=\left\{\begin{matrix}
    \dfrac{1}{2k-1}, & n=2k-1,\\[6pt]
    0, & n=2k,
  \end{matrix}\right.
  \qquad
  a_{n}^{-}=\left\{\begin{matrix}
    0, & n=2k-1,\\[6pt]
    \dfrac{1}{2k}, & n=2k.
  \end{matrix}\right.
\]
\[
  \sum_{n=1}^{\infty}a_{n}^{+}=\sum_{k=1}^{\infty}\frac{1}{2k-1}=+\infty,
  \qquad
  \sum_{n=1}^{\infty}a_{n}^{-}=\sum_{k=1}^{\infty}\frac{1}{2k}=+\infty.
\]
从
\[
\begin{aligned}
  S&=1-\frac{1}{2}+\frac{1}{3}-\frac{1}{4}+\frac{1}{5}-\frac{1}{6}
     +\frac{1}{7}-\frac{1}{8}+\frac{1}{9}-\frac{1}{10}+\frac{1}{11}-\frac{1}{12}+\cdots\\
  +\Bigr)\ \frac{1}{2}S&=\ \ \ \ \ \ \frac{1}{2}\ \ \ \ -\frac{1}{4}\ \ \ \ +\frac{1}{6}\ \ \ \ -\frac{1}{8}\ \ \ \ +\frac{1}{10}\ \ \ \ -\frac{1}{12}+\cdots\\
  \hline
  \frac{3}{2}S&=1\ \ \ \ +\frac{1}{3}-\frac{1}{2}+\frac{1}{5}\ \ \ \ +\frac{1}{7}-\frac{1}{4}
     +\frac{1}{9}\ \ \ \ +\frac{1}{11}-\frac{1}{6}+\cdots
\end{aligned}
\]
可看出上述级数是由原级数 $\sum_{n=1}^{\infty}(-1)^{n-1}\frac{1}{n}$ 交换了无穷多项的次序得到的. 它的和为 $\frac{3}{2}S$, 与原级数的和 $S$ 不相等.
\end{proof}

更一般地, 有下面的 Riemann 定理.

\begin{theorem}[条件收敛级的特性, Riemann]\label{thm:12.3.7}
设级数 $\sum_{n=1}^{\infty}a_{n}$ 条件收敛, 则适当交换其项的次序, 可使其收敛于任一事先指定的实数 $S$, 也可使其发散.
\end{theorem}

\begin{proof}
(1) 设 $|S|<+\infty$.

因为 $\sum_{n=1}^{\infty}a_{n}$ 条件收敛, 由定理~\ref{thm:12.3.5}(2) 知, $\sum_{n=1}^{\infty}a_{n}^{+}=\sum_{n=1}^{\infty}a_{n}^{-}=+\infty$. 于是, 可取 $n_{1}\in\mathbb{N}$ 使
\[
  a_{1}^{+}+a_{2}^{+}+\cdots+a_{n_{1}-1}^{+}\leqslant S
  <a_{1}^{+}+a_{2}^{+}+\cdots+a_{n_{1}-1}^{+}+a_{n_{1}}^{+}
\]
（$n_{1}=1$ 表示 $S<a_{n_{1}}^{+}=a_{1}^{+}$）, 这时, 有
\[
  0<(a_{1}^{+}+a_{2}^{+}+\cdots+a_{n_{1}-1}^{+}+a_{n_{1}}^{+})-S\leqslant a_{n_{1}}^{+}.
\]
再取 $n_{1}'\in\mathbb{N}$, 使
\[
  a_{1}^{+}+\cdots+a_{n_{1}}^{+}-a_{1}^{-}-\cdots-a_{n_{1}'-1}^{-}-a_{n_{1}'}^{-}
  <S\leqslant a_{1}^{+}+\cdots+a_{n_{1}}^{+}-a_{1}^{-}-\cdots-a_{n_{1}'-1}^{-},
\]
这时, 有
\[
  0<S-\biggl(\sum_{k=1}^{n_{1}}a_{k}^{+}-\sum_{k=1}^{n_{1}'}a_{k}^{-}\biggr)\leqslant a_{n_{1}'}^{-}.
\]
再取 $n_{2}\in\mathbb{N}$, 使
\[
  \sum_{k=1}^{n_{1}}a_{k}^{+}-\sum_{k=1}^{n_{1}'}a_{k}^{-}+\sum_{k=n_{1}+1}^{n_{2}-1}a_{k}^{+}
  \leqslant S
  <\sum_{k=1}^{n_{1}}a_{k}^{+}-\sum_{k=1}^{n_{1}'}a_{k}^{-}+\sum_{k=n_{1}+1}^{n_{2}}a_{k}^{+}.
\]
于是, 有
\[
  0<\sum_{k=1}^{n_{1}}a_{k}^{+}-\sum_{k=1}^{n_{1}'}a_{k}^{-}
  +\sum_{k=n_{1}+1}^{n_{2}}a_{k}^{+}-S\leqslant a_{n_{2}}^{+}.
\]
依次类推, 可得到一无穷级数, 它是由原级数交换无穷多项的次序得到的新级数. 下面证该级数收敛于 $S$. 事实上, 总有
\[
  0<\biggl(\sum_{k=1}^{n_{1}}a_{k}^{+}-\sum_{k=1}^{n_{1}'}a_{k}^{-}
  +\cdots+\sum_{k=n_{l-1}+1}^{n_{l}}a_{k}^{+}\biggr)-S\leqslant a_{n_{l}}^{+},
\]
\[
  0<S-\biggl(\sum_{k=1}^{n_{1}}a_{k}^{+}-\sum_{k=1}^{n_{1}'}a_{k}^{-}
  +\cdots+\sum_{k=n_{l-1}+1}^{n_{l}}a_{k}^{+}
  -\sum_{k=n_{l-1}'+1}^{n_{l}'}a_{k}^{-}\biggr)\leqslant a_{n_{l}'}^{-}.
\]
由于 $\sum_{n=1}^{\infty}a_{n}$ 收敛, 根据定理~\ref{thm:12.1.2}, $\lim\limits_{n\to+\infty}a_{n}=0$, 从而有
\[
  \lim_{l\to+\infty}a_{n_{l}}^{+}=0=\lim_{l\to+\infty}a_{n_{l}'}^{-}.
\]
这就证明了
\[
  (a_{1}^{+}+\cdots+a_{n_{1}}^{+})-(a_{1}^{-}+\cdots+a_{n_{1}'}^{-})
  +(a_{n_{1}+1}^{+}+\cdots+a_{n_{2}}^{+})-(a_{n_{1}'+1}^{-}+\cdots+a_{n_{2}'}^{-})+\cdots
\]
的和为 $S$. 由于上述级数的每个括号内有相同的符号, 根据定理~\ref{thm:12.1.5}, 打开括号后的级数也收敛, 且和仍为 $S$.

(2) $S=+\infty$.

取 $n_{1}\in\mathbb{N}$, 使 $\sum_{k=1}^{n_{1}}a_{k}^{+}>1$. 再取 $n_{2}$ 使 $\sum_{k=1}^{n_{1}}a_{k}^{+}-a_{1}^{-}+\sum_{k=n_{1}+1}^{n_{2}}a_{k}^{+}>2$, $\cdots$, 依次可取 $n_{l}$ 使 $\sum_{k=1}^{n_{1}}a_{k}^{+}-a_{1}^{-}+\sum_{k=n_{1}+1}^{n_{2}}a_{k}^{+}-a_{2}^{-}+\cdots-a_{l-1}^{-}+\sum_{k=n_{l-1}+1}^{n_{l}}a_{k}^{+}>l$. 由于 $\lim\limits_{n\to+\infty}a_{n}^{-}=0$, 故当 $l$ 充分大时, 有 $0\leqslant a_{n_{l}-1}^{-}<\frac{1}{2}$. 于是, 由数列极限为 $+\infty$ 的定义知, 适当交换原级数项的次序可使得到的新级数发散到 $+\infty$.

类似可作新级数使其发散到 $-\infty$. 也可作新级数使其部分和无极限.
\end{proof}

让我们用一个例子来重温一下上述的 Riemann 定理 \ref{thm:12.3.7}.

\begin{example}\label{ex:12.3.7}
将级数
\[
  1-\frac{1}{\sqrt{2}}+\frac{1}{\sqrt{3}}-\frac{1}{\sqrt{4}}
  +\frac{1}{\sqrt{5}}-\frac{1}{\sqrt{6}}+\cdots
\]
的项重新安排如下: 先依次取 $p$ 个正项, 接着依次取 $q$ 个负项, 再接着依次取 $p$ 个正项, 如此下去. 证明: 所得新级数收敛的充分必要条件是 $p=q$; 当 $p>q$ 时, 新级数发散到 $+\infty$; 当 $p<q$ 时, 新级数发散到 $-\infty$.
\end{example}

\begin{proof}
证明 由\extref{例 1.4.2} 知, 级数
\[
  1+\frac{1}{\sqrt{2}}+\frac{1}{\sqrt{3}}+\cdots+\frac{1}{\sqrt{n}}+\cdots
  =\sum_{n=1}^{\infty}\frac{1}{\sqrt{n}}=\sum_{n=1}^{\infty}\frac{1}{n^{\frac{1}{2}}}
\]
发散. 再由 Leibniz 判别法知
\[
  1-\frac{1}{\sqrt{2}}+\frac{1}{\sqrt{3}}-\frac{1}{\sqrt{4}}+\cdots
  +(-1)^{n-1}\frac{1}{\sqrt{n}}+\cdots
  =\sum_{n=1}^{\infty}\frac{(-1)^{n-1}}{\sqrt{n}}
\]
收敛, 因而它是条件收敛的级数.

设重排以后的新级数为 $\sum_{n=1}^{\infty}a_{n}$. 对于任给的正整数 $N$, 记 $m=\bigl[\frac{N}{p+q}\bigr]$, 则当 $N\to+\infty$ 时有 $m\to+\infty$, 且 $m(p+q)\leqslant N<(m+1)(p+q)$. 将 $\sum_{n=1}^{\infty}a_{n}$ 的部分和写成
\[
  \sum_{n=1}^{N}a_{n}=\sum_{n=1}^{m(p+q)}a_{n}+\sum_{n=m(p+q)+1}^{N}a_{n}.
\]
因为 $N-m(p+q)<p+q$, 这说明等式右端第 2 个和式的项数不超过 $p+q$. 所以, 当 $N\to+\infty$ 时, 有
\[
\begin{aligned}
  \biggl|\sum_{n=m(p+q)+1}^{N}a_{n}\biggr|
  &\leqslant\sum_{n=m(p+q)+1}^{N}|a_{n}|
   \leqslant(p+q)\frac{1}{\sqrt{m(p+q)}}
   =\frac{\sqrt{p+q}}{\sqrt{m}}\to 0
\end{aligned}
\]
（$N\to+\infty$, $m\to+\infty$）. 于是,\[
\begin{aligned}
  \sum_{n=1}^{\infty}a_{n}
  &=\lim_{N\to+\infty}\sum_{n=1}^{N}a_{n}
   =\lim_{m\to+\infty}\sum_{n=1}^{m(p+q)}a_{n}\\
  &=\lim_{m\to+\infty}\biggl(\sum_{n=1}^{mp}\frac{1}{\sqrt{2n-1}}
   -\sum_{n=1}^{mq}\frac{1}{\sqrt{2n}}\biggr)\\
  &=\lim_{m\to+\infty}\Bigl\{\bigl[\sqrt{2mp}+\bigl(1-\tfrac{1}{\sqrt{2}}\bigr)\beta
   +O\bigl(\tfrac{1}{\sqrt{m}}\bigr)\bigr]
   -\bigl[\sqrt{2mq}+\tfrac{1}{\sqrt{2}}\beta+O\bigl(\tfrac{1}{\sqrt{m}}\bigr)\bigr]\Bigr\}\\
  &=\lim_{m\to+\infty}\bigl[\sqrt{2m}\,(\sqrt{p}-\sqrt{q})
   +(1-\sqrt{2})\beta+O\bigl(\tfrac{1}{\sqrt{m}}\bigr)\bigr]\\
  &=\left\{\begin{matrix}
    (1-\sqrt{2})\beta, & \text{当}\ p=q\ \text{时},\\
    +\infty, & \text{当}\ p>q\ \text{时},\\
    -\infty, & \text{当}\ p<q\ \text{时}.
  \end{matrix}\right.
\end{aligned}
\]
其中第 3 个等式是在引理 \ref{lem:12.3.3} 中取 $f(x)=\frac{1}{\sqrt{x}}$, 立即可得
\[
  \sum_{k=1}^{n}\frac{1}{\sqrt{k}}=2\sqrt{n}+\beta+O\Bigl(\frac{1}{\sqrt{n}}\Bigr),
  \quad n\to+\infty,
\]
这里 $\beta$ 为某个常数. 于是,
\[
\begin{aligned}
  \sum_{n=1}^{mp}\frac{1}{\sqrt{2n-1}}
  &=\sum_{n=1}^{2mp}\frac{1}{\sqrt{n}}-\frac{1}{\sqrt{2}}\sum_{n=1}^{mp}\frac{1}{\sqrt{n}}\\
  &=\bigl[2\sqrt{2mp}+\beta+O\bigl(\tfrac{1}{\sqrt{2mp}}\bigr)\bigr]
   -\frac{1}{\sqrt{2}}\bigl[2\sqrt{mp}+\beta+O\bigl(\tfrac{1}{\sqrt{mp}}\bigr)\bigr]\\
  &=\sqrt{2mp}+\bigl(1-\tfrac{1}{\sqrt{2}}\bigr)\beta+O\bigl(\tfrac{1}{\sqrt{m}}\bigr),
\end{aligned}
\]
\[
\begin{aligned}
  \sum_{n=1}^{mq}\frac{1}{\sqrt{2n}}
  &=\frac{1}{\sqrt{2}}\sum_{n=1}^{mq}\frac{1}{\sqrt{n}}
   =\frac{1}{\sqrt{2}}\biggl[2\sqrt{mq}+\beta+O\biggl(\frac{1}{\sqrt{mq}}\biggr)\biggr]\\
  &=\sqrt{2mq}+\frac{1}{\sqrt{2}}\beta+O\biggl(\frac{1}{\sqrt{m}}\biggr).
\end{aligned}
\]
\end{proof}

\begin{lemma}\label{lem:12.3.3}
设 $a\in\mathbb{N}$, 当 $x\geqslant a$ 时函数 $f$ 非负单调减, 则极限
\[
  \lim_{n\to+\infty}\biggl(\sum_{k=a}^{n}f(k)-\int_{a}^{n}f(x)\,\mathrm{d}x\biggr)=\alpha
\]
存在, 且 $0\leqslant\alpha\leqslant f(a)$. 更进一步, 如果 $\lim\limits_{x\to+\infty}f(x)=0$, 则
\[
  \biggl|\sum_{k=a}^{[\xi]}f(k)-\int_{a}^{\xi}f(x)\,\mathrm{d}x-\alpha\biggr|
  \leqslant f(\xi-1),
\]
其中 $\xi\geqslant a+1$.

(1) 取 $f(x)=\frac{1}{\sqrt{x}}$, $a=1$, $\alpha=\beta+2$ 得到
\[
\begin{aligned}
  \biggl|\sum_{k=1}^{n}\frac{1}{\sqrt{k}}-2\sqrt{n}-\beta\biggr|
  &=\biggl|\sum_{k=1}^{n}\frac{1}{\sqrt{k}}-2\sqrt{n}+2-\alpha\biggr|\\
  &=\biggl|\sum_{k=1}^{n}\frac{1}{\sqrt{k}}-\int_{1}^{n}\frac{\mathrm{d}x}{\sqrt{x}}-\alpha\biggr|
   \leqslant\frac{1}{\sqrt{n-1}},
\end{aligned}
\]
即
\[
  \sum_{k=1}^{n}\frac{1}{\sqrt{k}}=2\sqrt{n}+\beta+O\Bigl(\frac{1}{\sqrt{n}}\Bigr),
  \quad n\to+\infty.
\]

(2) 取 $f(x)=\frac{1}{x}$, $a=1$, $\alpha=C$ 得到
\[
  \biggl|\sum_{k=1}^{n}\frac{1}{k}-\ln n-C\biggr|
  =\biggl|\sum_{k=1}^{n}\frac{1}{k}-\ln n-\alpha\biggr|
  \leqslant\frac{1}{n-1},
\]
即
\[
  \sum_{k=1}^{n}\frac{1}{k}=\ln n+C+O\Bigl(\frac{1}{n}\Bigr),
  \quad n\to+\infty,
\]
这里的常数 $C$ 称为 Euler 常数（参阅\extref{定理 1.4.2}(4)）. 可证 $C=0.57721\cdots$.
\end{lemma}

\begin{proof}
证明 令
\[
  g(\xi)=\sum_{k=a}^{[\xi]}f(k)-\int_{a}^{\xi}f(x)\,\mathrm{d}x,
\]
这里 $a\in\mathbb{N}$, 有
\[
\begin{aligned}
  g(n)-g(n+1)
  &=-f(n+1)+\int_{n}^{n+1}f(x)\,\mathrm{d}x\\
  &\geqslant-f(n+1)+f(n+1)=0.
\end{aligned}
\]
另一方面,
\[
  g(n)=\sum_{k=a}^{n-1}\biggl(f(k)-\int_{k}^{k+1}f(x)\,\mathrm{d}x\biggr)+f(n)
  \geqslant\sum_{k=a}^{n-1}(f(k)-f(k))+f(n)=f(n)\geqslant0.
\]
这说明 $\{g(n)\}$ 是一个非负单调减的数列且有下界 $0$. 因此,
$\alpha=\lim\limits_{n\to+\infty}g(n)
=\lim\limits_{n\to+\infty}\bigl[\sum_{k=a}^{n}f(k)-\int_{a}^{n}f(x)\,\mathrm{d}x\bigr]$
存在且有限. 由于 $0\leqslant g(n)\leqslant g(a)=f(a)$, 故 $0\leqslant\alpha\leqslant f(a)$.

如果进一步假定 $\lim\limits_{x\to+\infty}f(x)=0$, 则
\[
\begin{aligned}
  g(\xi)-\alpha
  &=\sum_{k=a}^{[\xi]}f(k)-\int_{a}^{\xi}f(x)\,\mathrm{d}x
   -\lim_{n\to+\infty}\biggl[\sum_{k=a}^{n}f(k)-\int_{a}^{n}f(x)\,\mathrm{d}x\biggr]\\
  &=\sum_{k=a}^{[\xi]}f(k)-\int_{a}^{[\xi]}f(x)\,\mathrm{d}x
   -\int_{[\xi]}^{\xi}f(x)\,\mathrm{d}x
   -\lim_{n\to+\infty}\biggl[\sum_{k=a}^{n}f(k)-\int_{a}^{n}f(x)\,\mathrm{d}x\biggr]\\
  &=-\int_{[\xi]}^{\xi}f(x)\,\mathrm{d}x
   -\lim_{n\to+\infty}\biggl(\sum_{k=[\xi]+1}^{n}f(k)-\int_{[\xi]}^{n}f(x)\,\mathrm{d}x\biggr)\\
  &=-\int_{[\xi]}^{\xi}f(x)\,\mathrm{d}x
   +\lim_{n\to+\infty}\sum_{k=[\xi]+1}^{n}\int_{k-1}^{k}\bigl[f(x)-f(k)\bigr]\,\mathrm{d}x\\
  &\leqslant\lim_{n\to+\infty}\sum_{k=[\xi]+1}^{n}\int_{k-1}^{k}\bigl[f(k-1)-f(k)\bigr]\,\mathrm{d}x
   =f([\xi])\leqslant f(\xi-1).
\end{aligned}
\]
类似地, 有
\[
\begin{aligned}
  g(\xi)-\alpha
  &=-\int_{[\xi]}^{\xi}f(x)\,\mathrm{d}x
   +\lim_{n\to+\infty}\sum_{k=[\xi]+1}^{n}\int_{k-1}^{k}\bigl[f(x)-f(k)\bigr]\,\mathrm{d}x\\
  &\geqslant-\int_{[\xi]}^{\xi}f(x)\,\mathrm{d}x
   \geqslant-\bigl(\xi-[\xi]\bigr)f\bigl([\xi]\bigr)
   \geqslant-f\bigl([\xi]\bigr)\geqslant-f(\xi-1).
\end{aligned}
\]
至此, 我们已经证得了
\[
  \biggl|\sum_{k=a}^{[\xi]}f(k)-\int_{a}^{\xi}f(x)\,\mathrm{d}x-\alpha\biggr|
  =|g(\xi)-\alpha|\leqslant f(\xi-1).
\]
\end{proof}

\begin{example}\label{ex:12.3.8}
(1) 取 $f(x)=\frac{1}{\sqrt{x}}$, $a=1$, $\alpha=\beta+2$ 得到
\[
\begin{aligned}
  \biggl|\sum_{k=1}^{n}\frac{1}{\sqrt{k}}-2\sqrt{n}-\beta\biggr|
  &=\biggl|\sum_{k=1}^{n}\frac{1}{\sqrt{k}}-2\sqrt{n}+2-\alpha\biggr|\\
  &=\biggl|\sum_{k=1}^{n}\frac{1}{\sqrt{k}}-\int_{1}^{n}\frac{\mathrm{d}x}{\sqrt{x}}-\alpha\biggr|
   \leqslant\frac{1}{\sqrt{n-1}},
\end{aligned}
\]
即
\[
  \sum_{k=1}^{n}\frac{1}{\sqrt{k}}=2\sqrt{n}+\beta+O\Bigl(\frac{1}{\sqrt{n}}\Bigr),
  \quad n\to+\infty.
\]

(2) 取 $f(x)=\frac{1}{x}$, $a=1$, $\alpha=C$ 得到
\[
  \biggl|\sum_{k=1}^{n}\frac{1}{k}-\ln n-C\biggr|
  =\biggl|\sum_{k=1}^{n}\frac{1}{k}-\ln n-\alpha\biggr|
  \leqslant\frac{1}{n-1},
\]
即
\[
  \sum_{k=1}^{n}\frac{1}{k}=\ln n+C+O\Bigl(\frac{1}{n}\Bigr),
  \quad n\to+\infty,
\]
这里的常数 $C$ 称为 Euler 常数（参阅\extref{定理 1.4.2}(4)）. 可证 $C=0.57721\cdots$.
\end{example}

\begin{exblock}{练习题 12.3}
\begin{enumerate}
  \item 应用 Cauchy 收敛原理讨论下列级数的敛散性:
  \begin{enumerate}
    \item $\displaystyle\sum_{n=1}^{\infty}\frac{1}{2n-1}$;
    \item $\displaystyle\sum_{n=1}^{\infty}\frac{\sin n}{2^{n}}$;
    \item $\displaystyle\sum_{n=1}^{\infty}\frac{\cos n!}{n(n+1)}$;
    \item $\displaystyle\sum_{n=1}^{\infty}\frac{a\cos n+b\sin n}{n(n+\sin n!)}$.
  \end{enumerate}
  \item 应用 Cauchy 收敛原理及 $\bigl|\sum_{n=N+1}^{N+p}(-1)^{n-1}a_{n}\bigr|\leqslant a_{N}$ 证明交错级数的 Leibniz 判别法.
  \item 设 $a_{n}\leqslant c_{n}\leqslant b_{n}$, $n=1,2,\cdots$, 如果 $\sum_{n=1}^{\infty}a_{n}$, $\sum_{n=1}^{\infty}b_{n}$ 都是收敛级数. 证明: $\sum_{n=1}^{\infty}c_{n}$ 也收敛. 如果 $\sum_{n=1}^{\infty}a_{n}$, $\sum_{n=1}^{\infty}b_{n}$ 都发散, 结论如何?
  \item 在交错级数的 Leibniz 判别法中, 如果去掉 $\{a_{n}\}$ 单调减这个条件, 结论可能不成立. 试以下例
  \[
    \sum_{n=3}^{\infty}\frac{1}{\sqrt{\bigl[\frac{n+1}{2}\bigr]+(-1)^{n}}}
    =\frac{1}{\sqrt{2}-1}-\frac{1}{\sqrt{2}+1}
    +\frac{1}{\sqrt{3}-1}-\frac{1}{\sqrt{3}+1}+\cdots
\]
  说明之.
  \item 讨论下列级数的敛散性:
  \begin{enumerate}
    \item $\displaystyle\sum_{n=1}^{\infty}(-1)^{n-1}\frac{\sqrt{n}}{n+1}$;
    \item $\displaystyle\sum_{n=1}^{\infty}(-1)^{n-1}\frac{1}{\sqrt[n]{n}}$;
    \item $\displaystyle\sum_{n=1}^{\infty}(-1)^{n-1}\sin\frac{1}{n}$;
    \item $\displaystyle\sum_{n=1}^{\infty}\frac{\sin nx}{n}$.
  \end{enumerate}
  \item 设 $\{a_{n}\}$ 单调减趋于 $0$, 讨论级数 $\sum_{n=1}^{\infty}a_{n}\cos nx$, $\sum_{n=1}^{\infty}a_{n}\sin nx$ 的敛散性.
  \item 设 $\sum_{n=1}^{\infty}a_{n}$ 为一个收敛级数, 如果 $\lim\limits_{n\to+\infty}\frac{b_{n}}{a_{n}}=1$, 能否断言 $\sum_{n=1}^{\infty}b_{n}$ 也为收敛级数? 请研究级数:
  \[
    \sum_{n=1}^{\infty}\frac{(-1)^{n-1}}{\sqrt{n}},
    \qquad
    \sum_{n=1}^{\infty}\biggl\{\frac{(-1)^{n-1}}{\sqrt{n}}+\frac{1}{n}\biggr\}.
\]
  \item 设 $\sum_{n=1}^{\infty}a_{n}$ 收敛, 应用分部求和公式证明:
  \[
    \lim_{n\to+\infty}\frac{a_{1}+2a_{2}+\cdots+na_{n}}{n}=0.
\]
  \item 下列级数, 哪些是绝对收敛的? 哪些是条件收敛的?
  \begin{enumerate}
    \item $\displaystyle\sum_{n=1}^{\infty}\frac{\sin nx}{n^{a}}$, $a>1$;
    \item $\displaystyle\sum_{n=1}^{\infty}\frac{(-1)^{n-1}}{2^{n}}\cos nx$;
    \item $\displaystyle\sum_{n=2}^{\infty}(-1)^{n}\frac{1}{n\ln n}$;
    \item $\displaystyle\sum_{n=1}^{\infty}(-1)^{\frac{1}{2}n(n-1)}\frac{n^{10}}{a^{n}}$, $a>1$;
    \item $\displaystyle\sum_{n=1}^{\infty}(-1)^{n}\frac{n-1}{n+1}\frac{1}{\sqrt[3]{n}}$;
    \item $\displaystyle\sum_{n=2}^{\infty}\frac{\sin\frac{n\pi}{4}}{\ln n}$.
  \end{enumerate}
  \item 讨论下列级数的绝对收敛性与条件收敛性:
  \begin{enumerate}
    \item $\displaystyle\sum_{n=1}^{\infty}\frac{(-1)^{n-1}}{n^{p+\frac{1}{n}}}$;
    \item $\displaystyle\sum_{n=1}^{\infty}(-1)^{n}\frac{\cos 2n}{n^{p}}$.
  \end{enumerate}
  \item 设 $\sum_{n=1}^{\infty}a_{n}$ 条件收敛, 证明:
  \begin{enumerate}
    \item $\displaystyle\sum_{n=1}^{\infty}a_{n}^{+}=\sum_{n=1}^{\infty}a_{n}^{-}=+\infty$, 逆命题是否成立?
    \item 记 $S_{N}^{+}=\sum_{n=1}^{N}a_{n}^{+}$, $S_{N}^{-}=\sum_{n=1}^{N}a_{n}^{-}$, 则
          $\displaystyle\lim_{N\to+\infty}\frac{S_{N}^{+}}{S_{N}^{-}}=1$.
  \end{enumerate}
\end{enumerate}
\end{exblock}

\begin{exblock}{思考题 12.3}
\begin{enumerate}
  \item 讨论下列级数的敛散性:
  \begin{enumerate}
    \item $\displaystyle\sum_{n=1}^{\infty}\sin\bigl(\pi\sqrt{n^{2}+1}\bigr)$;
    \item $\displaystyle\sum_{n=1}^{\infty}\biggl(1+\frac{1}{2}+\cdots+\frac{1}{n}\biggr)\frac{\sin nx}{n}$.
  \end{enumerate}
  \item 设级数 $\sum_{n=1}^{\infty}\frac{a_{n}}{n^{\alpha}}$ 收敛, 则 $\forall\beta>\alpha$, 级数 $\sum_{n=1}^{\infty}\frac{a_{n}}{n^{\beta}}$ 也收敛.
  \item 设 $\lim\limits_{n\to+\infty}n\bigl(\frac{a_{n}}{a_{n+1}}-1\bigr)=\lambda>0$, 则交错级数 $\sum_{n=1}^{\infty}(-1)^{n-1}a_{n}$ 收敛.
  \item 设 $a>0$, $b>0$. 应用 Raabe 判别法证明: 当 $b-1>a$ 时, 级数
  \[
    1+\sum_{n=1}^{\infty}\frac{a(a+1)\cdots(a+n-1)}{b(b+1)\cdots(b+n-1)}
\]
  收敛, 其和为 $\frac{b-1}{b-a-1}$. 当 $b-1\leqslant a$ 时, 收敛情况如何?
  \item 设 $\{a_{n}\}$ 为正的单调增数列, 证明: 级数 $\sum_{n=1}^{\infty}\bigl(\frac{a_{n+1}}{a_{n}}-1\bigr)$ 收敛 $\Leftrightarrow$ $\{a_{n}\}$ 有界.
  \item 应用 Leibniz 判别法证明:
  \begin{enumerate}
    \item $\displaystyle\sum_{n=1}^{\infty}\frac{(-1)^{[\sqrt{n}]}}{n}$ 收敛;
    \item $\displaystyle 1-\frac{1}{3}\biggl(1+\frac{1}{2}\biggr)
          +\frac{1}{5}\biggl(1+\frac{1}{2}+\frac{1}{3}\biggr)
          -\frac{1}{7}\biggl(1+\frac{1}{2}+\frac{1}{3}+\frac{1}{4}\biggr)+\cdots$ 收敛.
  \end{enumerate}
  \item 讨论下列级数的绝对收敛性与条件收敛性:
  \begin{enumerate}
    \item $\displaystyle\sum_{n=2}^{\infty}\ln\biggl(1+\frac{(-1)^{n}}{n^{p}}\biggr)$, $p>0$;
    \item $\displaystyle\sum_{n=1}^{\infty}(-1)^{n-1}
          \biggl(\frac{1\cdot 3\cdot 5\cdot\cdots\cdot(2n-1)}
          {2\cdot 4\cdot 6\cdot\cdots\cdot(2n)}\biggr)^{p}$;
    \item $\displaystyle\frac{1}{1^{p}}-\frac{1}{2^{q}}+\frac{1}{3^{p}}-\frac{1}{4^{q}}
          +\frac{1}{5^{p}}-\frac{1}{6^{q}}+\cdots$.
  \end{enumerate}
  \item 若将级数
  \[
    1-\frac{1}{2}+\frac{1}{3}-\frac{1}{4}+\frac{1}{5}-\frac{1}{6}+\cdots
\]
  的项重新这样安排: 先依次取 $p$ 个正项, 接着依次取 $q$ 个负项, 再接着依次取 $p$ 个正项, 如此继续下去. 应用公式
  $1+\frac{1}{2}+\cdots+\frac{1}{n}=\ln n+C+o(1)$（参阅\extref{定理 1.4.2}(5) 证明）证明所得新级数的和为
  \[
    \ln 2+\frac{1}{2}\ln\frac{p}{q}.
\]
  \item 将级数
  \[
    1-\frac{1}{2^{\alpha}}+\frac{1}{3^{\alpha}}-\frac{1}{4^{\alpha}}
    +\frac{1}{5^{\alpha}}-\frac{1}{6^{\alpha}}+\cdots,\quad 0<\alpha<1
\]
  的项重新安排如下: 先依次取 $p$ 个正项, 接着依次取 $q$ 个负项, 再接着依次取 $p$ 个正项, 如此继续下去. 证明: 所得新级数收敛 $\Leftrightarrow p=q$; 当 $p>q$ 时, 新级数发散到 $+\infty$; 当 $p<q$ 时, 新级数发散到 $-\infty$.
  \item 如果改变调和级数
  \[
    1+\frac{1}{2}+\frac{1}{3}+\cdots+\frac{1}{n}+\cdots
\]
  中某些项的符号, 使得 $p$ 个正项后, 跟 $q$ 个负项, 但不改变原来的次序. 证明: 新得的级数仅当 $p=q$ 时收敛.
  \item 设 $S_{n}=q\cos\alpha+q^{2}\cos 2\alpha+\cdots+q^{n}\cos n\alpha
        =\sum_{k=1}^{n}q^{k}\cos k\alpha$.
  \begin{enumerate}
    \item 证明:
    \[
      2qS_{n}\cos\alpha
      =\bigl[q^{n+1}\cos(n+1)\alpha+S_{n}-q\cos\alpha\bigr]
      +\bigl[q^{2}+q^{2}S_{n}-q^{n+2}\cos n\alpha\bigr];
\]
    \item $\displaystyle\sum_{k=1}^{\infty}q^{k}\cos k\alpha
          =\lim_{n\to+\infty}S_{n}
          =\frac{q\cos\alpha-q^{2}}{1+q^{2}-2q\cos\alpha}$, $|q|<1$;
    \item 利用复数和式 $\displaystyle\sum_{k=0}^{n}z^{k}=\frac{1-z^{n+1}}{1-z}$,
          取实部推得 (1) 与 (2).
  \end{enumerate}
  \item 证明:
  \[
    1+\frac{1}{2}+\biggl(\frac{1}{3}-1\biggr)+\frac{1}{4}+\frac{1}{5}
    +\biggl(\frac{1}{6}-\frac{1}{2}\biggr)+\frac{1}{7}+\frac{1}{8}
    +\biggl(\frac{1}{9}-\frac{1}{3}\biggr)+\cdots=\ln 3.
\]
  \item 设 $S_{n}(x)=\sum_{k=1}^{n}\frac{\sin kx}{k}$.
  \begin{enumerate}
    \item 证明: $\displaystyle S_{n}'(x)=\frac{\sin\bigl(n+\frac{1}{2}\bigr)x}{2\sin\frac{x}{2}}-\frac{1}{2}$,
    \[
      S_{n}(x)=-\frac{1}{2}\int_{x}^{\pi}\frac{1}{\sin\frac{t}{2}}
      \sin\bigl(n+\tfrac{1}{2}\bigr)t\,\mathrm{d}t+\frac{1}{2}(\pi-x);
\]
    \item $\displaystyle S(x)=\sum_{k=1}^{\infty}\frac{\sin kx}{k}
          =\left\{\begin{array}
{ll}
            0, & x=0,\\[
2pt]
            \dfrac{\pi-x}{2}, & 0<x\leqslant\pi.
\end{array}\right.$
          （对照例 \ref{ex:16.1.2}）.
  \end{enumerate}
  \item 设级数 $\sum_{n=1}^{\infty}a_{n}$ 的前 $n$ 项和 $S_{n}$ 分成两项, 即 $S_{n}=S_{n}^{+}+S_{n}^{-}$, 其中 $S_{n}^{+}$ 与 $S_{n}^{-}$ 分别为正项之和与负项之和. 如果 $\sum_{n=1}^{\infty}a_{n}$ 条件收敛, 证明: $\lim\limits_{n\to+\infty}\frac{S_{n}^{+}}{S_{n}^{-}}=-1$.
  \item 举出一个发散的交错级数, 使其通项趋于 $0$. 考察级数
  \[
    1-\frac{1}{2}+\frac{1}{3^{2}}-\frac{1}{4}+\frac{1}{5^{2}}-\cdots
    +\frac{1}{(2n-1)^{2}}-\frac{1}{2n}+\cdots.
\]
  \item 设 $a_{n}>0$, $a_{n}>a_{n+1}$（$n=1,2,\cdots$）, 且 $\lim\limits_{n\to+\infty}a_{n}=0$. 证明: 级数
  \[
    \sum_{n=1}^{\infty}(-1)^{n-1}\frac{a_{1}+a_{2}+\cdots+a_{n}}{n}
\]
  是收敛的.
\end{enumerate}
\end{exblock}

% ==================================================================
%  徐森林、金亚东、薛春华《数学分析》（第三册）
%  S12.4  级数的乘法
%  源：原书印刷页 55-62（PDF 物理页 61-68）
%  底本：PaddleOCR-VL-1.6，见 chapters/_ch12_12_4_raw.md
%  约定：全西文标点；行间公式用一对方括号；多行公式用 align/aligned；
%        跨册引用一律 \extref；数学模式内的中文必须包 \text{}
%  已订正的 OCR 错误见 tools/_make_ch12_4.py 头部说明
% ==================================================================
\section{级数的乘法}\label{sec:12.4}

设有两个收敛级数 $\sum_{n=1}^{\infty}a_{n}=A$, $\sum_{n=1}^{\infty}b_{n}=B$. 这两个无穷级数如何相乘? 对于有限和的乘法, 有
\[
  \begin{aligned}
    \sum_{i=1}^{m}a_{i}\cdot\sum_{j=1}^{n}b_{j}&=(a_{1}+a_{2}+\cdots+a_{m})(b_{1}+b_{2}+\cdots+b_{n})\\
    &=a_{1}b_{1}+a_{2}b_{1}+\cdots+a_{m}b_{1}\\
    &\quad+a_{1}b_{2}+a_{2}b_{2}+\cdots+a_{m}b_{2}+\cdots\\
    &\quad+a_{1}b_{n}+a_{2}b_{n}+\cdots+a_{m}b_{n}.
  \end{aligned}
\]

形式地套到无穷级数的乘积上, 可以写作:
\[
  \begin{aligned}
    &a_{1}b_{1}+a_{2}b_{1}+a_{3}b_{1}+\cdots+a_{m}b_{1}+\cdots\\
    &+a_{1}b_{2}+a_{2}b_{2}+a_{3}b_{2}+\cdots+a_{m}b_{2}+\cdots\\
    &+a_{1}b_{3}+a_{2}b_{3}+a_{3}b_{3}+\cdots+a_{m}b_{3}+\cdots\\
    &+a_{1}b_{n}+a_{2}b_{n}+a_{3}b_{n}+\cdots+a_{m}b_{n}+\cdots\\
    &=\sum_{i,j=1}^{\infty}a_{i}b_{j}.
  \end{aligned}
\]

这些项如何求它的部分和? 通常有两种方法:

(1) 沿对角线相加:
\[
  \begin{aligned}
    &a_{1}b_{1}+(a_{1}b_{2}+a_{2}b_{1})+(a_{1}b_{3}+a_{2}b_{2}+a_{3}b_{1})+(a_{1}b_{4}+a_{2}b_{3}+a_{3}b_{2}+a_{4}b_{1})+\cdots\\
    &+(a_{1}b_{n}+a_{2}b_{n-1}+\cdots+a_{n-1}b_{2}+a_{n}b_{1})+\cdots=\sum_{n=1}^{\infty}c_{n},
  \end{aligned}
\]
\[
  c_{n}=\sum_{i=1}^{n}a_{i}b_{n-i+1}=a_{1}b_{n}+a_{2}b_{n-1}+\cdots+a_{n-1}b_{2}+a_{n}b_{1}.
\]

(2) 沿方块相加:
\[
  \begin{aligned}
    &a_{1}b_{1}+(a_{1}b_{2}+a_{2}b_{2}+a_{2}b_{1})+(a_{1}b_{3}+a_{2}b_{3}+a_{3}b_{3}+a_{3}b_{2}+a_{3}b_{1})\\
    &+\left(a_{1}b_{4}+a_{2}b_{4}+a_{3}b_{4}+a_{4}b_{4}+a_{5}b_{4}+a_{6}b_{3}+a_{7}b_{3}+a_{8}b_{2}+a_{9}b_{1}\right)+\cdots=\sum_{n=1}^{\infty}d_{n},
  \end{aligned}
\]
其中 $d_{n}=a_{1}b_{n}+a_{2}b_{n}+\cdots+a_{n}b_{n}+a_{n}b_{n-1}+\cdots+a_{n}b_{1}$（共 $2n-1$ 项）.

用得较多的是对角线法. 问题是 $\sum_{n=1}^{\infty}c_{n}=AB$? 先观察例子.

\begin{example}\label{ex:12.4.1}
设 $a_{n}=b_{n}=\frac{(-1)^{n-1}}{\sqrt{n}}$, 由 Leibniz 判别法知, $\sum_{n=1}^{\infty}a_{n}=\sum_{n=1}^{\infty}b_{n}=\sum_{n=1}^{\infty}\frac{(-1)^{n-1}}{\sqrt{n}}$ 收敛. 而从
\[
  \begin{aligned}
    c_{n}&=\sum_{i=1}^{n}a_{i}b_{n-i+1}=a_{1}b_{n}+a_{2}b_{n-1}+a_{3}b_{n-2}+\cdots+a_{n}b_{1}\\
    &=1\cdot\frac{(-1)^{n-1}}{\sqrt{n}}+\frac{-1}{\sqrt{2}}\cdot\frac{(-1)^{n-2}}{\sqrt{n-1}}+\frac{(-1)^{2}}{\sqrt{3}}\cdot\frac{(-1)^{n-3}}{\sqrt{n-2}}+\cdots+\frac{(-1)^{n-1}}{\sqrt{n}}\cdot1\\
    &=(-1)^{n-1}\left[\frac{1}{\sqrt{1\cdot n}}+\frac{1}{\sqrt{2(n-1)}}+\frac{1}{\sqrt{3(n-2)}}+\cdots+\frac{1}{\sqrt{n\cdot1}}\right],
  \end{aligned}
\]
\[
  |c_{n}|=\sum_{k=1}^{n}\frac{1}{\sqrt{k(n-k+1)}}\geqslant\frac{n}{\frac{k+(n-k+1)}{2}}=\frac{2n}{n+1}\to2\neq0\quad(n\to+\infty),
\]
即 $\lim\limits_{n\to+\infty}c_{n}\neq0$ 以及定理~\ref{thm:12.1.2} 立即推得 $\sum_{n=1}^{\infty}c_{n}$ 发散.
\end{example}

这表明收敛级数的乘积按对角线求和未必收敛. 但是, 如果 $\sum_{n=1}^{\infty}a_{n}$ 与 $\sum_{n=1}^{\infty}b_{n}$ 都是绝对收敛的级数, 情形就不一样了.

\begin{theorem}[Cauchy]\label{thm:12.4.1}
设级数 $\sum_{n=1}^{\infty}a_{n}$ 与 $\sum_{n=1}^{\infty}b_{n}$ 都绝对收敛, 其和分别为 $A$ 与 $B$, 则 $\sum_{n=1}^{\infty}c_{n}$ 也是绝对收敛的级数, 且 $\sum_{n=1}^{\infty}c_{n}=AB$, 其中 $c_{n}=\sum_{i=1}^{n}a_{i}b_{n-i+1}$.
\end{theorem}

\begin{proof}
设 $A_{N}=\sum_{n=1}^{N}a_{n}$, $B_{N}=\sum_{n=1}^{N}b_{n}$ 为部分和. 由条件知
\[
  \lim_{N\to+\infty}A_{N}=A,\quad\lim_{N\to+\infty}B_{N}=B.
\]
再设
\[
  A_{N}^{*}=\sum_{n=1}^{N}|a_{n}|,\quad B_{N}^{*}=\sum_{n=1}^{N}|b_{n}|,\quad A^{*}=\lim_{N\to+\infty}A_{N}^{*},\quad B^{*}=\lim_{N\to+\infty}B_{N}^{*}.
\]

(1)
\[
  \begin{aligned}
    C_{N}^{*}&=\sum_{n=1}^{N}|c_{n}|=\sum_{n=1}^{N}\left|a_{1}b_{n}+a_{2}b_{n-1}+\cdots+a_{n}b_{1}\right|\\
    &\leqslant\sum_{n=1}^{N}\bigl[|a_{1}b_{n}|+|a_{2}b_{n-1}|+\cdots+|a_{n}b_{1}|\bigr]\\
    &=|a_{1}b_{1}|+(|a_{1}b_{2}|+|a_{2}b_{1}|)+\cdots+(|a_{1}b_{N}|+\cdots+|a_{N}b_{1}|)\\
    &=|a_{1}b_{1}|+|a_{2}b_{1}|+\cdots+|a_{N}b_{1}|\\
    &\quad+|a_{1}b_{2}|+|a_{2}b_{2}|+\cdots+|a_{N}b_{2}|+\cdots\\
    &\quad+|a_{1}b_{N}|+|a_{2}b_{N}|+\cdots+|a_{N}b_{N}|\\
    &=\sum_{n=1}^{N}|a_{n}|\cdot\sum_{n=1}^{N}|b_{n}|=A_{N}^{*}B_{N}^{*}\leqslant A^{*}B^{*},
  \end{aligned}
\]
即部分和 $\sum_{n=1}^{N}|c_{n}|$ 有上界, 因此, $\sum_{n=1}^{\infty}|c_{n}|$ 收敛, 从而 $\sum_{n=1}^{\infty}c_{n}$ 绝对收敛.

(2) 再证 $\sum_{n=1}^{\infty}c_{n}=AB$. 令
\[
  \begin{aligned}
    C_{N}&=\sum_{n=1}^{N}c_{n}=a_{1}b_{1}+(a_{1}b_{2}+a_{2}b_{1})+(a_{1}b_{3}+a_{2}b_{2}+a_{3}b_{1})+\cdots\\
    &\quad+(a_{1}b_{N}+a_{2}b_{N-1}+\cdots+a_{N}b_{1})\to C\quad(N\to+\infty).
  \end{aligned}
\]
由于 $\sum_{n=1}^{\infty}|c_{n}|=|a_{1}b_{1}|+(|a_{1}b_{2}|+|a_{2}b_{1}|)+\cdots+(|a_{1}b_{n}|+\cdots+|a_{n}b_{1}|)+\cdots$ 绝对收敛, 故打开括号
\[
  |a_{1}b_{1}|+|a_{1}b_{2}|+|a_{2}b_{1}|+|a_{1}b_{3}|+|a_{2}b_{2}|+|a_{3}b_{1}|+\cdots
\]
也收敛, 即级数
\[
  a_{1}b_{1}+a_{1}b_{2}+a_{2}b_{1}+a_{1}b_{3}+a_{2}b_{2}+a_{3}b_{1}+\cdots
\]
绝对收敛. 根据定理~\ref{thm:12.3.6}, 可以任意交换它的无穷多项, 且任意结合所得新级数也绝对收敛, 且和不变. 因此, 可按方块形式相加得
\[
  C=\lim_{N\to+\infty}A_{N}B_{N}=\lim_{N\to+\infty}A_{N}\cdot\lim_{N\to+\infty}B_{N}=AB.
\]
\end{proof}

如果将定理~\ref{thm:12.4.1} 中条件减弱为有一个级数绝对收敛, 而另一个级数收敛, 则有下面的定理.

\begin{theorem}[Mertens]\label{thm:12.4.2}
设级数 $\sum_{n=1}^{\infty}a_{n}$ 绝对收敛, $\sum_{n=1}^{\infty}b_{n}$ 收敛, 且 $\sum_{n=1}^{\infty}a_{n}=A$, $\sum_{n=1}^{\infty}b_{n}=B$, 则
\[
  \sum_{n=1}^{\infty}c_{n}=AB,
\]
其中 $c_{n}=a_{1}b_{n}+a_{2}b_{n-1}+\cdots+a_{n}b_{1}=\sum_{i=1}^{n}a_{i}b_{n-i+1}$.
\end{theorem}

\begin{proof}
令 $A_{n}=\sum_{k=1}^{n}a_{k}$, $B_{n}=\sum_{k=1}^{n}b_{k}$, $C_{n}=\sum_{k=1}^{n}c_{k}$. 则 $\lim\limits_{n\to+\infty}A_{n}=A$, $\lim\limits_{n\to+\infty}B_{n}=B$. 再令 $\beta_{n}=B-B_{n}$, 显然有
\[
  \lim_{n\to+\infty}\beta_{n}=\lim_{n\to+\infty}(B-B_{n})=B-B=0.
\]
于是,
\[
  \begin{aligned}
    C_{n}&=\sum_{k=1}^{n}c_{k}=a_{1}b_{1}+(a_{1}b_{2}+a_{2}b_{1})+\cdots+(a_{1}b_{n}+a_{2}b_{n-1}+\cdots+a_{n}b_{1})\\
    &\xlongequal{\text{重新组合}}a_{1}(b_{1}+\cdots+b_{n})+a_{2}(b_{1}+\cdots+b_{n-1})+\cdots+a_{n}b_{1}\\
    &=a_{1}B_{n}+a_{2}B_{n-1}+\cdots+a_{n}B_{1}=a_{1}(B-\beta_{n})+a_{2}(B-\beta_{n-1})+\cdots+a_{n}(B-\beta_{1})\\
    &=(a_{1}+a_{2}+\cdots+a_{n})B-(a_{1}\beta_{n}+a_{2}\beta_{n-1}+\cdots+a_{n}\beta_{1})\\
    &\xrightarrow{}AB-0=AB\quad(n\to+\infty).
  \end{aligned}
\]
这就证明了 $\sum_{n=1}^{\infty}c_{n}=\lim\limits_{N\to+\infty}C_{N}=AB$.

最后, 我们来证明 $\lim\limits_{n\to+\infty}(a_{1}\beta_{n}+\cdots+a_{n}\beta_{1})=0$. 事实上, 由于 $\lim\limits_{n\to+\infty}\beta_{n}=0$, 故 $|\beta_{n}|<M$, $\forall n\in\mathbb{N}$. 且 $\forall\varepsilon>0$, $\exists N_{1}\in\mathbb{N}$, 当 $n>N_{1}$ 时,
\[
  |\beta_{n}|<\frac{\varepsilon}{2A^{*}+1},\quad\text{其中}\quad A^{*}=\sum_{n=1}^{\infty}|a_{n}|.
\]
又由 $\sum_{n=1}^{\infty}|a_{n}|$ 绝对收敛, 故从级数的 Cauchy 收敛准则, $\exists N_{2}\in\mathbb{N}$, 当 $n>N_{2}$ 时,
\[
  |a_{n+1}|+\cdots+|a_{n+p}|<\frac{\varepsilon}{2M}.
\]
取 $N=N_{1}+N_{2}$, 当 $n>N$ 时, $n-N_{1}>N-N_{1}=N_{2}$, 有
\[
  \begin{aligned}
    &\left|a_{1}\beta_{n}+a_{2}\beta_{n-1}+\cdots+a_{n}\beta_{1}\right|\\
    &\leqslant\left|a_{1}\beta_{n}\right|+\cdots+\left|a_{n-N_{1}}\beta_{N_{1}+1}\right|+\left|a_{n-N_{1}+1}\beta_{N_{1}}\right|+\cdots+\left|a_{n}\beta_{1}\right|\\
    &<\left(\left|a_{1}\right|+\cdots+\left|a_{n-N_{1}}\right|\right)\frac{\varepsilon}{2A^{*}+1}+M\bigl(\left|a_{n-N_{1}+1}\right|+\cdots+\left|a_{n}\right|\bigr)\\
    &<A^{*}\cdot\frac{\varepsilon}{2A^{*}+1}+M\cdot\frac{\varepsilon}{2M}<\frac{\varepsilon}{2}+\frac{\varepsilon}{2}=\varepsilon.
  \end{aligned}
\]
\end{proof}

\begin{example}\label{ex:12.4.2}
证明: $\sum_{n=1}^{\infty}nx^{n-1}=\frac{1}{(1-x)^{2}}$, $|x|<1$.
\end{example}

证法 1 因为
\[
  1+x+x^{2}+\cdots+x^{n-1}+\cdots=\frac{1}{1-x},\quad|x|<1
\]
绝对收敛, 所以
\[
  \begin{aligned}
    \frac{1}{(1-x)^{2}}&=\left(\frac{1}{1-x}\right)^{2}=\left(\sum_{n=1}^{\infty}x^{n-1}\right)^{2}\xlongequal[\text{或定理~\ref{thm:12.4.2}}]{\text{定理~\ref{thm:12.4.1}}}\sum_{n=1}^{\infty}c_{n}\\
    &=\sum_{n=1}^{\infty}(1\cdot x^{n-1}+x\cdot x^{n-2}+\cdots+x^{n-1}\cdot1)=\sum_{n=1}^{\infty}nx^{n-1}.
  \end{aligned}
\]

证法 2 根据幂级数的逐项求导定理 定理~\ref{thm:14.1.4} 得到.
\[
  \frac{1}{(1-x)^{2}}=\left(\frac{1}{1-x}\right)'=\left(\sum_{n=0}^{\infty}x^{n}\right)'\xlongequal{\text{逐项求导}}\sum_{n=0}^{\infty}\left(x^{n}\right)'=\sum_{n=1}^{\infty}nx^{n-1},\quad|x|<1.
\]

\begin{example}\label{ex:12.4.3}
考察幂级数
\[
  E(x)\xlongequal{\mathrm{def}}\sum_{n=0}^{\infty}\frac{x^{n}}{n!}=1+x+\frac{x^{2}}{2!}+\frac{x^{3}}{3!}+\cdots+\frac{x^{n}}{n!}+\cdots,
\]
证明: $E(x)$ 收敛, 因而定义了 $\mathbb{R}$ 上的一个函数, 并且
\[
  E(x)E(y)=E(x+y).
\]
如果记 $E(x)=\mathrm{e}^{x}$, 则 $\mathrm{e}^{x}\cdot\mathrm{e}^{y}=\mathrm{e}^{x+y}$.
\end{example}

证法 1 因为
\[
  \frac{|a_{n+1}|}{|a_{n}|}=\frac{\left|\frac{x^{n}}{n!}\right|}{\left|\frac{x^{n-1}}{(n-1)!}\right|}=\frac{|x|}{n}\to0<1\quad(n\to+\infty),
\]
根据 d'Alembert 判别法知, 对 $\forall x\in\mathbb{R}$, 级数 $\sum_{n=1}^{\infty}\frac{x^{n}}{n!}$ 绝对收敛. 从而 $E(x)=\sum_{n=1}^{\infty}\frac{x^{n}}{n!}$ 为定义在 $\mathbb{R}$ 上的一个函数. 并且
\[
  \begin{aligned}
    E(x)\cdot E(y)&=\sum_{n=0}^{\infty}\frac{x^{n}}{n!}\cdot\sum_{n=0}^{\infty}\frac{y^{n}}{n!}\xlongequal[\text{或定理~\ref{thm:12.4.2}}]{\text{定理~\ref{thm:12.4.1}}}\sum_{n=0}^{\infty}c_{n}(x,y)\\
    &=\sum_{n=0}^{\infty}\frac{(x+y)^{n}}{n!}=E(x+y),
  \end{aligned}
\]
其中
\[
  c_{0}(x,y)=1,\quad c_{1}(x,y)=1\cdot y+x\cdot1=x+y,\quad\cdots,
\]
\[
  \begin{aligned}
    c_{n}(x,y)&=1\cdot\frac{y^{n}}{n!}+\frac{x}{1!}\cdot\frac{y^{n-1}}{(n-1)!}+\frac{x^{2}}{2!}\cdot\frac{y^{n-2}}{(n-2)!}+\cdots+\frac{x^{n-1}}{(n-1)!}\cdot\frac{y}{1!}+\frac{x^{n}}{n!}\cdot1\\
    &=\sum_{k=0}^{n}\frac{x^{k}\cdot y^{n-k}}{k!(n-k)!}=\frac{1}{n!}\sum_{k=0}^{n}\frac{n!}{k!(n-k)!}x^{k}y^{n-k}\\
    &=\frac{1}{n!}\sum_{k=0}^{n}C_{n}^{k}x^{k}y^{n-k}=\frac{(x+y)^{n}}{n!}.
  \end{aligned}
\]

证法 2 根据幂级数的逐项求导定理 定理~\ref{thm:14.1.4} 得到
\[
  E'(x)=\left(\sum_{n=0}^{\infty}\frac{x^{n}}{n!}\right)'=\sum_{n=0}^{\infty}\left(\frac{x^{n}}{n!}\right)'=\sum_{n=1}^{\infty}\frac{x^{n-1}}{(n-1)!}=\sum_{n=0}^{\infty}\frac{x^{n}}{n!}=E(x).
\]
于是,
\[
  \begin{aligned}
    \frac{\partial}{\partial x}\left[\frac{E(x+y)}{E(x)E(y)}\right]&=\frac{E'(x+y)E(x)E(y)-E(x+y)E'(x)E(y)}{E^{2}(x)E^{2}(y)}\\
    &=\frac{E(x+y)E(x)E(y)-E(x+y)E(x)E(y)}{E^{2}(x)E^{2}(y)}=0.
  \end{aligned}
\]
同理,
\[
  \frac{\partial}{\partial y}\left[\frac{E(x+y)}{E(x)E(y)}\right]=0.
\]
根据\extref{定理 8.2.10}, 有
\[
  \frac{E(x+y)}{E(x)E(y)}=\frac{E(0+0)}{E(0)E(0)}=\frac{1}{1\cdot1}=1,
\]
即
\[
  E(x+y)=E(x)E(y).
\]

证法 3 设 $a>0$.

当 $m,n\in\mathbb{N}$ 时,
\[
  \begin{aligned}
    a^{m}\cdot a^{n}&=(\underbrace{a\cdots a}_{\text{$m$ 个}})\cdot(\underbrace{a\cdots a}_{\text{$n$ 个}})=\underbrace{a\cdots a}_{\text{$m+n$ 个}}=a^{m+n},
  \end{aligned}
\]
\[
  a^{-m}\cdot a^{n}=\underbrace{\frac{1}{a\cdots a}}_{\text{$m$ 个}}\cdot(\underbrace{a\cdots a}_{\text{$n$ 个}})=a^{-m+n},
\]
同理 $a^{m}\cdot a^{-n}=a^{m-n}$; 此外, 还有
\[
  a^{-m}\cdot a^{-n}=\underbrace{\frac{1}{a\cdots a}}_{\text{$m$ 个}}\cdot\underbrace{\frac{1}{a\cdots a}}_{\text{$n$ 个}}=\frac{1}{a^{m+n}}=a^{-(m+n)}=a^{-m-n}.
\]

当 $m=0$, $n\in\mathbb{Z}$（整数集）时,
\[
  a^{m}\cdot a^{n}=a^{0}\cdot a^{n}=1\cdot a^{n}=a^{n}=a^{0+n}=a^{m+n}.
\]
同理, 当 $m\in\mathbb{Z}$, $n=0$ 时, 也有 $a^{m}\cdot a^{n}=a^{m+n}$.

当 $m,p\in\mathbb{Z}$, $n,q\in\mathbb{Z}$ 时,
\[
  \begin{aligned}
    a^{\frac{m}{n}}\cdot a^{\frac{p}{q}}&=a^{\frac{mq}{nq}}\cdot a^{\frac{np}{nq}}=\left(a^{\frac{1}{nq}}\right)^{mq}\cdot\left(a^{\frac{1}{nq}}\right)^{np}\\
    &=\left(a^{\frac{1}{nq}}\right)^{mq+np}=a^{\frac{mq+np}{nq}}=a^{\frac{m}{n}+\frac{p}{q}}.
  \end{aligned}
\]

当 $x,y\in\mathbb{R}$ 时, 取有理数列 $x_{n}$ 与 $y_{n}$, 使得 $\lim\limits_{n\to+\infty}x_{n}=x$, $\lim\limits_{n\to+\infty}y_{n}=y$. 于是,
\[
  a^{x}\cdot a^{y}=\lim_{n\to+\infty}a^{x_{n}}\cdot\lim_{n\to+\infty}a^{y_{n}}=\lim_{n\to+\infty}\bigl(a^{x_{n}}\cdot a^{y_{n}}\bigr)=\lim_{n\to+\infty}a^{x_{n}+y_{n}}=a^{x+y}.
\]

另一方面, 由证法 2 知, $E'(x)=E(x)$, 即 $(\ln E(x))'=\frac{E'(x)}{E(x)}=1$. 由 Newton-Leibniz 公式得到
\[
  \ln E(x)=\ln E(0)+\int_{0}^{x}(\ln E(x))'\mathrm{d}x=\ln1+\int_{0}^{x}\mathrm{d}x=x,
\]
\[
  E(x)=\mathrm{e}^{x}.
\]
于是,
\[
  E(x)E(y)=\mathrm{e}^{x}\cdot\mathrm{e}^{y}=\mathrm{e}^{x+y}=E(x+y).
\]

证法 4 由\extref{例 4.1.2} 知,
\[
  \mathrm{e}^{x}=\sum_{k=0}^{n}\frac{x^{k}}{k!}+R_{n}(x)=\sum_{k=0}^{n}\frac{x^{k}}{k!}+\frac{\mathrm{e}^{\theta x}}{(n+1)!}x^{n+1},\quad\theta\in(0,1).
\]
另一方面, 由\extref{例 1.1.5}, 有 $\lim\limits_{n\to+\infty}\frac{|x|^{n+1}}{(n+1)!}=0$, 或者由
\[
  \lim_{n\to+\infty}\frac{\frac{|x|^{n+1}}{(n+1)!}}{\frac{|x|^{n}}{n!}}=\lim_{n\to+\infty}\frac{|x|}{n+1}=0
\]
及 d'Alembert 判别法推得 $\sum_{n=0}^{\infty}\frac{|x|^{n}}{n!}$ 收敛. 再由定理~\ref{thm:12.1.2}, $\lim\limits_{n\to+\infty}\frac{|x|^{n+1}}{(n+1)!}=0$. 于是,
\[
  \left|\frac{\mathrm{e}^{\theta x}}{(n+1)!}x^{n+1}\right|\leqslant\mathrm{e}^{|x|}\frac{|x|^{n+1}}{(n+1)!}\to0\quad(n\to+\infty),
\]
\[
  \begin{aligned}
    E(x)&=\sum_{n=0}^{\infty}\frac{x^{n}}{n!}=\lim_{n\to+\infty}\sum_{k=0}^{n}\frac{x^{k}}{k!}=\lim_{n\to+\infty}\left[\mathrm{e}^{x}-\frac{\mathrm{e}^{\theta x}}{(n+1)!}x^{n+1}\right]\\
    &=\mathrm{e}^{x}-0=\mathrm{e}^{x},
  \end{aligned}
\]
\[
  E(x)E(y)=\mathrm{e}^{x}\cdot\mathrm{e}^{y}=\mathrm{e}^{x+y}=E(x+y).
\]

\begin{exblock}{练习题 12.4}
\begin{enumerate}
  \item 设 $|q|<1$, 证明:
  \[
    \Big(\sum_{n=0}^{\infty}q^{n}\Big)^{2}=\sum_{n=0}^{\infty}(n+1)q^{n}.
\]
  \item 设级数 $\sum_{n=0}^{\infty}a_{n}x^{n}$, $\sum_{n=0}^{\infty}b_{n}x^{n}$ 对 $(-R,R)$ 中的 $x$ 绝对收敛. 证明:
  \[
    \left(\sum_{n=0}^{\infty}a_{n}x^{n}\right)\left(\sum_{n=0}^{\infty}b_{n}x^{n}\right)=\sum_{n=0}^{\infty}c_{n}x^{n},
\]
  其中 $c_{n}=\sum_{k=0}^{n}a_{k}b_{n-k}$.
  \item 证明: 级数
  \[
    C(x)=\sum_{n=0}^{\infty}(-1)^{n}\frac{x^{2n}}{(2n)!},\quad S(x)=\sum_{n=0}^{\infty}(-1)^{n}\frac{x^{2n+1}}{(2n+1)!}
\]
  对所有实数 $x$ 绝对收敛, 而且
  \[
    S(2x)=2S(x)C(x).
\]
\end{enumerate}
\end{exblock}

\begin{exblock}{思考题 12.4}
\begin{enumerate}
  \item 证明: 下面两个级数
  \[
    \sum_{n=1}^{\infty}(-1)^{n-1}\frac{1}{n^{\alpha}},\quad\sum_{n=1}^{\infty}(-1)^{n-1}\frac{1}{n^{\beta}},\quad\alpha>0,\ \beta>0
\]
  的 Cauchy 乘积当 $\alpha+\beta>1$ 时收敛; 当 $\alpha+\beta\leqslant1$ 时发散.
  \item 证明: 下面两个发散级数
  \[
    1-\sum_{n=1}^{\infty}\left(\frac{3}{2}\right)^{n},\quad1+\sum_{n=1}^{\infty}\left(\frac{3}{2}\right)^{n-1}\left(2^{n}+\frac{1}{2^{n+1}}\right)
\]
  的 Cauchy 乘积是一个绝对收敛级数.
\end{enumerate}
\end{exblock}

% ==================================================================
%  徐森林、金亚东、薛春华《数学分析》（第三册）
%  S12.5  无穷乘积
%  源：原书印刷页 62-73（PDF 物理页 68-79）
%  底本：PaddleOCR-VL-1.6，见 chapters/_ch12_12_5_raw.md
%  约定：全西文标点；行间公式用一对方括号；多行公式用 align/aligned；
%        跨册引用一律 \extref；数学模式内的中文必须包 \text{}
%  已订正的 OCR 错误与原书照录清单见 tools/_make_ch12_5a.py 头部说明
% ==================================================================
\section{无穷乘积}\label{sec:12.5}

与无穷级数相对偶地可研究无穷乘积.

\begin{definition}\label{def:12.5.1}
设 $p_{1},p_{2},\cdots,p_{n},\cdots$ 为一给定的数列. 称形式乘积
\[
  \prod_{n=1}^{\infty}p_{n}=p_{1}p_{2}\cdots p_{n}\cdots
\]
为无穷乘积. $P_{n}=p_{1}p_{2}\cdots p_{n}$ 称为 $\prod_{n=1}^{\infty}p_{n}$ 的第 $n$ 个部分乘积.
\end{definition}

如果 $\lim\limits_{n\to+\infty}P_{n}=P\neq0$, 则称无穷乘积 $\prod_{n=1}^{\infty}p_{n}$ 收敛, 且表示为 $\prod_{n=1}^{\infty}p_{n}=P$. 否则（$\lim\limits_{n\to+\infty}P_{n}=0$, $+\infty$, $-\infty$, $\infty$ 或不存在）称无穷乘积 $\prod_{n=1}^{\infty}p_{n}$ 发散. 特别当 $P=0$ 时, 称无穷乘积 $\prod_{n=1}^{\infty}p_{n}$ 发散于 $0$.

\begin{theorem}[无穷乘积收敛的必要条件]\label{thm:12.5.1}
设无穷乘积 $\prod_{n=1}^{\infty}p_{n}$ 收敛, 则 $\lim\limits_{n\to+\infty}p_{n}=1$. 但反之不真.
\end{theorem}

\begin{proof}
设无穷乘积 $\prod_{n=1}^{\infty}p_{n}$ 收敛于 $P\neq0$, 即 $\lim\limits_{n\to+\infty}P_{n}=P$. 于是,
\[
  \lim_{n\to+\infty}p_{n}=\lim_{n\to+\infty}\frac{P_{n}}{P_{n-1}}=\frac{\lim\limits_{n\to+\infty}P_{n}}{\lim\limits_{n\to+\infty}P_{n-1}}=\frac{P}{P}=1.
\]
但反之不真. 例如: 考察无穷级数 $\prod_{n=1}^{\infty}\left(1+\frac{1}{n}\right)$. 显然, $\lim\limits_{n\to+\infty}p_{n}=\lim\limits_{n\to+\infty}\left(1+\frac{1}{n}\right)=1$. 而 $P_{n}=\prod_{k=1}^{n}\left(1+\frac{1}{k}\right)=\frac{2}{1}\cdot\frac{3}{2}\cdots\frac{n}{n-1}\cdot\frac{n+1}{n}=n+1$, $\prod_{n=1}^{\infty}\left(1+\frac{1}{n}\right)$ 发散于 $\lim\limits_{n\to+\infty}P_{n}=+\infty$.
\end{proof}

\begin{example}\label{ex:12.5.1}
$\sum_{n=1}^{\infty}p_{n}$ 发散于 $0$, 未必有 $\lim\limits_{n\to+\infty}p_{n}=1$.
\end{example}

(1) $p_{n}=\frac{1}{2}$, $\lim\limits_{n\to+\infty}p_{n}=\frac{1}{2}\neq1$, 但
\[
  \prod_{n=1}^{\infty}p_{n}=\prod_{n=1}^{\infty}\frac{1}{2}=\lim_{n\to+\infty}\prod_{k=1}^{n}\frac{1}{2}=\lim_{n\to+\infty}\frac{1}{2^{n}}=0.
\]

(2) $p_{n}=\frac{(-1)^{n}}{2}$, $\lim\limits_{n\to+\infty}p_{n}$ 不存在, 当然 $\lim\limits_{n\to+\infty}p_{n}\neq1$, 但
\[
  \prod_{n=1}^{\infty}p_{n}=\prod_{n=1}^{\infty}\frac{(-1)^{n}}{2}=\lim_{n\to+\infty}\prod_{k=1}^{n}\frac{(-1)^{k}}{2}=\lim_{n\to+\infty}\frac{(-1)^{\frac{n(n+1)}{2}}}{2^{n}}=0.
\]

(3) 设 $p_{n_{0}}=0$, 则当 $n>n_{0}$ 时, $P_{n}=p_{1}\cdots p_{n_{0}}\cdots p_{n}=0$, 因而
\[
  \prod_{n=1}^{\infty}p_{n}=\lim_{n\to+\infty}P_{n}=\lim_{n\to+\infty}0=0.
\]

\begin{theorem}[无穷乘积与无穷级数敛散性的对偶定理]\label{thm:12.5.2}
设 $p_{n}>0$, 则
\begin{enumerate}
  \item $\prod_{n=1}^{\infty}p_{n}$ 收敛于 $P>0 \Leftrightarrow \sum_{n=1}^{\infty}\ln p_{n}$ 收敛于 $S=\ln P$（$P=\mathrm{e}^{S}$）;
  \item $\prod_{n=1}^{\infty}p_{n}$ 发散于 $0 \Leftrightarrow \sum_{n=1}^{\infty}\ln p_{n}$ 发散于 $-\infty$.
\end{enumerate}
\end{theorem}

\begin{proof}
(1) $\prod_{n=1}^{\infty}p_{n}$ 收敛于 $P>0$
$\Leftrightarrow$ 数列 $P_{n}=\prod_{k=1}^{n}p_{k}$ 收敛于 $P>0$
$\Leftrightarrow$ 数列 $\ln P_{n}=\sum_{k=1}^{n}\ln p_{k}$ 收敛于 $\ln P=S$
$\Leftrightarrow$ $\sum_{n=1}^{\infty}\ln p_{n}$ 收敛于 $\ln P=S$.

(2) $\prod_{n=1}^{\infty}p_{n}$ 发散于 $0$
$\Leftrightarrow$ 数列 $P_{n}=\prod_{k=1}^{n}p_{k}$ 收敛于 $0$
$\Leftrightarrow$ 数列 $\ln P_{n}=\sum_{k=1}^{n}\ln p_{k}$ 发散于 $-\infty$.
\end{proof}

如果 $\prod_{n=1}^{\infty}p_{n}$ 收敛, 由定理~\ref{thm:12.5.1} 知, $\lim\limits_{n\to+\infty}p_{n}=1$, 则 $\exists N\in\mathbb{N}$, 当 $n>N$ 时, $p_{n}>0$. 令 $p_{n}=1+a_{n}$, $a_{n}>-1$. 记
\[
  \prod_{n=1}^{\infty}p_{n}=\prod_{n=1}^{\infty}(1+a_{n}).
\]

此无穷乘积收敛 $\Rightarrow\lim\limits_{n\to+\infty}p_{n}=1\Leftrightarrow\lim\limits_{n\to+\infty}a_{n}=0$.

这样改变无穷乘积的形式, 能否带来什么好结果呢? 以下诸定理给出了肯定的回答. 并且将判别无穷乘积敛散的问题转化为判别无穷级数敛散的问题.

\begin{theorem}\label{thm:12.5.3}
如果 $a_{n}\geqslant0$（或 $a_{n}\leqslant0$ 且 $1+a_{n}>0$,）, $n=1,2,\cdots$, 则 $\prod_{n=1}^{\infty}(1+a_{n})$ 与 $\sum_{n=1}^{\infty}a_{n}$ 同敛散.
\end{theorem}

\begin{proof}
因为不论 $\prod_{n=1}^{\infty}(1+a_{n})$ 与 $\sum_{n=1}^{\infty}a_{n}$ 何者收敛, 都有 $\lim\limits_{n\to+\infty}a_{n}=0$. 此时, 若 $a_{n}\neq0$ 有
\[
  \lim_{n\to+\infty}\frac{\ln(1+a_{n})}{a_{n}}=\lim_{n\to+\infty}\ln(1+a_{n})^{\frac{1}{a_{n}}}=\ln\mathrm{e}=1,
\]
所以, $\sum_{n=1}^{\infty}\ln(1+a_{n})$ 与 $\sum_{n=1}^{\infty}a_{n}$ 同敛散（同为正项级数或同为负项级数, 再由比较判别法的极限形式. 注意, $a_{n}=0$ 不影响敛散性的证明）. 根据定理~\ref{thm:12.5.2} 有, $\prod_{n=1}^{\infty}(1+a_{n})$ 与 $\sum_{n=1}^{\infty}\ln(1+a_{n})$ 同敛散, 所以, $\prod_{n=1}^{\infty}(1+a_{n})$ 与 $\sum_{n=1}^{\infty}a_{n}$ 同敛散.
\end{proof}

\begin{remark}\label{rem:12.5.1}
在定理~\ref{thm:12.5.3} 中, 如果 $a_{n}\geqslant0$, $\sum_{n=1}^{\infty}a_{n}=+\infty$, 则
\[
  \prod_{k=1}^{n}(1+a_{k})\geqslant a_{1}+a_{2}+\cdots+a_{n}=\sum_{k=1}^{n}a_{k}\to+\infty\quad(n\to+\infty),
\]
从而 $\prod_{n=1}^{\infty}(1+a_{n})=+\infty$.

如果 $a_{n}\leqslant0$, $1+a_{n}\geqslant0$ 且 $\sum_{n=1}^{\infty}a_{n}=-\infty$, 则 $P_{n}=\prod_{k=1}^{n}(1+a_{k})$ 单调减有下界 $0$. 此时 $\prod_{n=1}^{\infty}(1+a_{n})$ 必发散于 $0$.

如果 $\{a_{n}\}$ 不保号, 就有下面的结论.
\end{remark}

\begin{theorem}\label{thm:12.5.4}
(1) 设 $\sum_{n=1}^{\infty}a_{n}^{2}$ 收敛, 则 $\prod_{n=1}^{\infty}(1+a_{n})$ 与 $\sum_{n=1}^{\infty}a_{n}$ 同敛散.

(2) 设 $\sum_{n=1}^{\infty}a_{n}$ 收敛, $\sum_{n=1}^{\infty}a_{n}^{2}=+\infty$, 则 $\prod_{n=1}^{\infty}(1+a_{n})$ 发散于 $0$.
\end{theorem}

\begin{proof}
(1) 因为 $\sum_{n=1}^{\infty}a_{n}^{2}$ 收敛, 根据定理~\ref{thm:12.1.2}, $\lim\limits_{n\to+\infty}a_{n}^{2}=0$, 从而 $\lim\limits_{n\to+\infty}a_{n}=0$. 于是, 当 $a_{n}\neq0$ 时, 有
\[
  \begin{aligned}
    \lim_{n\to+\infty}\frac{a_{n}-\ln(1+a_{n})}{a_{n}^{2}}&=\lim_{x\to0}\frac{x-\ln(1+x)}{x^{2}}\xlongequal{\text{L'Hospital 法则}}\lim_{x\to0}\frac{1-\frac{1}{1+x}}{2x}\\
    &=\lim_{x\to0}\frac{1}{2(1+x)}=\frac{1}{2}.
  \end{aligned}
\]
从而, $\sum_{n=1}^{\infty}[a_{n}-\ln(1+a_{n})]$ 也收敛（注意, 若 $a_{n}=0$, 不影响收敛性的证明）. 由此推得 $\sum_{n=1}^{\infty}\ln(1+a_{n})$ 与 $\sum_{n=1}^{\infty}a_{n}$ 同敛散. 再由定理~\ref{thm:12.5.2} 知, $\prod_{n=1}^{\infty}(1+a_{n})$ 与 $\sum_{n=1}^{\infty}\ln(1+a_{n})$ 同敛散, 所以 $\prod_{n=1}^{\infty}(1+a_{n})$ 与 $\sum_{n=1}^{\infty}a_{n}$ 同敛散.

(2) 从 $\sum_{n=1}^{\infty}a_{n}^{2}=+\infty$ 与
\[
  \lim_{n\to+\infty}\frac{a_{n}-\ln(1+a_{n})}{a_{n}^{2}}=\frac{1}{2}\quad(a_{n}\neq0)
\]
便知 $\sum_{n=1}^{\infty}[a_{n}-\ln(1+a_{n})]=+\infty$. 再由 $\sum_{n=1}^{\infty}a_{n}$ 收敛推得 $\sum_{n=1}^{\infty}\ln(1+a_{n})=-\infty$, 因而 $\prod_{n=1}^{\infty}(1+a_{n})$ 发散到 $0$.
\end{proof}

\begin{theorem}\label{thm:12.5.5}
设 $1+a_{n}>0$, 则 $\prod_{n=1}^{\infty}(1+|a_{n}|)$ 收敛蕴涵着 $\prod_{n=1}^{\infty}(1+a_{n})$ 收敛. 但反之不真.
\end{theorem}

\begin{proof}
显然,
\[
  \begin{gathered}
    \prod_{n=1}^{\infty}(1+|a_{n}|)\ \text{收敛}
    \xLongleftrightarrow{\text{定理~\ref{thm:12.5.2}(1)}}
    \sum_{n=1}^{\infty}\ln(1+|a_{n}|)\ \text{收敛}\\[3pt]
    \xLongleftrightarrow{\text{定理~\ref{thm:12.5.3}}}
    \sum_{n=1}^{\infty}|a_{n}|\ \text{收敛}
    \xLongleftarrow[\text{(*)}]{\text{比较判别法}}
    \sum_{n=1}^{\infty}|\ln(1+a_{n})|\ \text{收敛}\\[3pt]
    \Rightarrow\sum_{n=1}^{\infty}\ln(1+a_{n})\ \text{收敛}
    \xLongleftrightarrow{\text{定理~\ref{thm:12.5.2}(1)}}
    \prod_{n=1}^{\infty}(1+a_{n})\ \text{收敛}.
  \end{gathered}
\]
其中 $(*)$ 是当 $a_{n}\neq0$ 时, 有
\[
  \lim_{n\to+\infty}\frac{\bigl|\ln(1+a_{n})\bigr|}{|a_{n}|}=\lim_{n\to+\infty}\left|\frac{\ln(1+a_{n})}{a_{n}}\right|=1.
\]

但是, 其逆命题不真. 反例: 设
\[
  a_{n}=\left\{\begin{array}
{ll}
    -\dfrac{1}{\sqrt{k+1}}, & n=2k-1,\\[8pt]
    \dfrac{1}{\sqrt{k+1}}+\dfrac{1}{k+1}+\dfrac{1}{(k+1)\sqrt{k+1}}, & n=2k,
\end{array}\right.
\]
则 $\sum_{n=3}^{\infty}a_{n}$, $\sum_{n=3}^{\infty}a_{n}^{2}$ 都发散, $\prod_{n=1}^{\infty}(1+|a_{n}|)$ 发散, 而 $\prod_{n=1}^{\infty}(1+a_{n})$ 收敛.

事实上, 因为
\[
  \begin{aligned}
    \sum_{n=1}^{2N}a_{n}&=\sum_{k=1}^{N}(a_{2k-1}+a_{2k})=\sum_{k=1}^{N}\left(-\frac{1}{\sqrt{k+1}}+\frac{1}{\sqrt{k+1}}+\frac{1}{k+1}+\frac{1}{(k+1)\sqrt{k+1}}\right)\\
    &=\sum_{k=1}^{N}\frac{1}{k+1}+\sum_{k=1}^{N}\frac{1}{(k+1)\sqrt{k+1}}\to+\infty\quad(N\to+\infty),
  \end{aligned}
\]
所以, $\sum_{n=1}^{\infty}a_{n}=+\infty$. 而
\[
  \sum_{n=1}^{2N}a_{n}^{2}=\sum_{k=1}^{N}(a_{2k-1}^{2}+a_{2k}^{2})\geqslant\sum_{k=2}^{N}\frac{1}{k+1}\to+\infty\quad(N\to+\infty).
\]
所以, $\sum_{n=1}^{\infty}a_{n}^{2}=+\infty$. 现在证明 $\prod_{n=1}^{\infty}(1+a_{n})$ 收敛. 因为
\[
  \begin{aligned}
    P_{2N}&=\prod_{n=1}^{2N}(1+a_{n})=\prod_{k=1}^{N}(1+a_{2k-1})(1+a_{2k})\\
    &=\prod_{k=1}^{N}\left(1-\frac{1}{\sqrt{k+1}}\right)\left(1+\frac{1}{\sqrt{k+1}}+\frac{1}{k+1}+\frac{1}{(k+1)\sqrt{k+1}}\right)\\
    &=\prod_{k=1}^{N}\left(1-\frac{1}{\sqrt{k+1}}\right)\left(1+\frac{1}{\sqrt{k+1}}\right)\left(1+\frac{1}{k+1}\right)=\prod_{k=1}^{N}\left[1-\frac{1}{(k+1)^{2}}\right],
  \end{aligned}
\]
由于 $\sum_{k=1}^{\infty}\frac{1}{(k+1)^{2}}$ 收敛与上面得到的结果知, $\prod_{k=1}^{\infty}\left(1-\frac{1}{(k+1)^{2}}\right)$ 收敛. 这说明数列 $\{P_{2N}\}$ 当 $N\to+\infty$ 时有非零的极限, 设 $\lim\limits_{N\to+\infty}P_{2N}=P\neq0$. 于是,
\[
  \lim_{N\to+\infty}P_{2N+1}=\lim_{N\to+\infty}P_{2N}(1+a_{2N+1})=P(1+0)=P.
\]
根据\extref{定理 1.1.1} 就证明了 $\lim\limits_{n\to+\infty}P_{n}=P\neq0$, 即 $\prod_{n=1}^{\infty}(1+a_{n})$ 收敛于 $P$.

再证 $\prod_{n=1}^{\infty}(1+|a_{n}|)$ 发散.（反证）假设 $\prod_{n=1}^{\infty}(1+|a_{n}|)$ 收敛, 由于定理~\ref{thm:12.5.3}, 它应等价于 $\sum_{n=1}^{\infty}|a_{n}|$ 收敛. 由此推得 $\sum_{n=1}^{\infty}a_{n}$ 收敛, 这与上面已证得 $\sum_{n=1}^{\infty}a_{n}=+\infty$（发散于 $+\infty$）相矛盾.
\end{proof}

\begin{example}\label{ex:12.5.2}
在定理~\ref{thm:12.5.4} 中, 如果 $\sum_{n=1}^{\infty}a_{n}^{2}=+\infty$, 问: $\prod_{n=1}^{\infty}(1+a_{n})$ 与 $\sum_{n=1}^{\infty}a_{n}$ 是否仍同敛散? 回答是否定的. 定理~\ref{thm:12.5.5} 中的反例也是这里的反例.
\end{example}

\begin{definition}\label{def:12.5.2}
如果 $\prod_{n=1}^{\infty}(1+|a_{n}|)$ 收敛, 则称 $\prod_{n=1}^{\infty}(1+a_{n})$ 绝对收敛; 如果 $\prod_{n=1}^{\infty}(1+a_{n})$ 收敛, 而 $\prod_{n=1}^{\infty}(1+|a_{n}|)$ 发散, 则称 $\prod_{n=1}^{\infty}(1+a_{n})$ 条件收敛.
\end{definition}
% ===== PART-B =====
% ===== PART-B =====

显然, 定理~\ref{thm:12.5.5} 表明 $\prod_{n=1}^{\infty}(1+a_{n})$ 绝收收敛, 它必收敛. 但是, 反之不真. 例~\ref{ex:12.5.6} 中, 当 $\frac{1}{2}<\alpha\leqslant1$ 时, $\prod_{n=1}^{\infty}\left[1+(-1)^{n-1}\frac{1}{n^{\alpha}}\right]$ 也是收敛非绝对收敛的反例.

根据无穷乘积与无穷级数敛散性的对偶定理~\ref{thm:12.5.2} 以及绝对收敛级数的特性定理~\ref{thm:12.3.6}, 我们有下面绝对收敛无穷乘积的特性定理.

\begin{theorem}[绝对收敛无穷乘积的特性]\label{thm:12.5.6}
绝对收敛的无穷乘积, 可以任意改变其因子的次序而不影响其收敛性与积.
\end{theorem}

\begin{proof}
设 $\prod_{n=1}^{\infty}(1+a_{n})$ 改变因子的先后次序变成 $\prod_{n=1}^{\infty}(1+b_{n})$. 因为 $\prod_{n=1}^{\infty}(1+a_{n})$ 绝对收敛, 即 $\prod_{n=1}^{\infty}(1+|a_{n}|)$ 收敛. 从定理~\ref{thm:12.5.5} 的证明知道, 它等价于 $\sum_{n=1}^{\infty}\ln(1+|a_{n}|)$ 收敛, 等价于 $\sum_{n=1}^{\infty}|a_{n}|$ 收敛, 也等价于 $\sum_{n=1}^{\infty}\ln(1+a_{n})$ 绝对收敛. 由定理~\ref{thm:12.3.6} 得到 $\sum_{n=1}^{\infty}\ln(1+b_{n})$ 也绝对收敛, 而且 $\sum_{n=1}^{\infty}\ln(1+b_{n})=\sum_{n=1}^{\infty}\ln(1+a_{n})$. 从而, $\prod_{n=1}^{\infty}(1+b_{n})$ 绝对收敛, 而且
\[
  \prod_{n=1}^{\infty}(1+a_{n})=\prod_{n=1}^{\infty}(1+b_{n}).
\]
\end{proof}

与条件收敛级数的特性定理~\ref{thm:12.3.7}（Riemann）相呼应, 有下面的定理.

\begin{theorem}[条件收敛无穷乘积的特性]\label{thm:12.5.7}
设无穷乘积 $\prod_{n=1}^{\infty}(1+a_{n})$ 条件收敛（$1+a_{n}>0$）. 则适当交换它因子的次序, 可使其收敛于任一事先指定的正数, 也可使其发散（发散于 $0$, $+\infty$, 或 $P_{n}=\prod_{k=1}^{n}(1+a_{k})$ 无极限）.
\end{theorem}

\begin{proof}
因为 $\prod_{n=1}^{\infty}(1+a_{n})$ 条件收敛, 所以 $\prod_{n=1}^{\infty}(1+a_{n})$ 收敛, 而 $\prod_{n=1}^{\infty}(1+|a_{n}|)$ 发散. 由定理~\ref{thm:12.5.3} 知, $\sum_{n=1}^{\infty}|a_{n}|=+\infty$. 再由
\[
  \lim_{n\to+\infty}\left|\frac{\ln(1+a_{n})}{a_{n}}\right|=1\quad(a_{n}\neq0)
\]
推得 $\sum_{n=1}^{\infty}|\ln(1+a_{n})|=+\infty$. 从而 $\sum_{n=1}^{\infty}\ln(1+a_{n})$ 为一个条件收敛的级数, 设其和为 $\alpha$. 今改变 $\sum_{n=1}^{\infty}\ln(1+a_{n})$ 的次序, 得一新级数 $\sum_{n=1}^{\infty}\ln(1+b_{n})$, 使其和为预先指定的实数 $\beta\neq\alpha$, 于是,
\[
  \prod_{n=1}^{\infty}(1+b_{n})=\mathrm{e}^{\beta}\neq\mathrm{e}^{\alpha}=\prod_{n=1}^{\infty}(1+a_{n}).
\]
而 $\prod_{n=1}^{\infty}(1+b_{n})$ 正是由 $\prod_{n=1}^{\infty}(1+a_{n})$ 改变因子次序所得的无穷乘积.

如果 $\beta=+\infty$, 则 $\sum_{n=1}^{\infty}(1+b_{n})=+\infty$; 如果 $\beta=-\infty$, 则 $\sum_{n=1}^{\infty}(1+b_{n})=0$（发散于 $0$）; 如果新级数 $\sum_{n=1}^{\infty}\ln(1+b_{n})$ 无极限, 则新无穷乘积 $\prod_{n=1}^{\infty}(1+b_{n})$ 的部分积无极限. 这些是罗列的新无穷乘积发散的各种情况.
\end{proof}

\begin{example}\label{ex:12.5.3}
(1) $\prod_{n=1}^{\infty}\left(1+\frac{1}{n}\right)=+\infty$; (2) $\prod_{n=2}^{\infty}\left(1-\frac{1}{n}\right)=0$.
\end{example}

证法 1 (1) $\prod_{n=1}^{\infty}\left(1+\frac{1}{n}\right)=\lim\limits_{n\to+\infty}\prod_{k=1}^{n}\left(1+\frac{1}{k}\right)=\lim\limits_{n\to+\infty}\left(\frac{2}{1}\cdot\frac{3}{2}\cdots\frac{n+1}{n}\right)=\lim\limits_{n\to+\infty}(n+1)=+\infty$.

(2)
\[
  \prod_{n=2}^{\infty}\left(1-\frac{1}{n}\right)=\lim_{n\to+\infty}\prod_{k=2}^{n}\left(1-\frac{1}{k}\right)=\lim_{n\to+\infty}\left(\frac{1}{2}\cdot\frac{2}{3}\cdots\frac{n-1}{n}\right)=\lim_{n\to+\infty}\frac{1}{n}=0.
\]

证法 2 从 $\sum_{n=1}^{\infty}\frac{1}{n}=+\infty$ 及定理~\ref{thm:12.5.3} 可知, $\prod_{n=1}^{\infty}\left(1+\frac{1}{n}\right)$ 与 $\prod_{n=2}^{\infty}\left(1-\frac{1}{n}\right)$ 都发散. 由于 $1+\frac{1}{n}>1$, $0<1-\frac{1}{n}<1$ 可推得
\[
  \prod_{n=1}^{\infty}\left(1+\frac{1}{n}\right)=+\infty,\quad\prod_{n=2}^{\infty}\left(1-\frac{1}{n}\right)=0.
\]

\begin{example}\label{ex:12.5.4}
证明:
(1) $\prod_{n=1}^{\infty}\left[1+\frac{1}{(n+1)^{2}}\right]$, (2) $\prod_{n=1}^{\infty}\left[1-\frac{1}{(n+1)^{2}}\right]$
都绝对收敛, 当然也是收敛的.
\end{example}

证明 (1)、(2) 从 $\sum_{n=1}^{\infty}\frac{1}{(n+1)^{2}}<+\infty$（收敛）及定理~\ref{thm:12.5.3} 立知,
\[
  \prod_{n=1}^{\infty}\left[1+\frac{1}{(n+1)^{2}}\right]\quad\text{与}\quad\prod_{n=1}^{\infty}\left[1-\frac{1}{(n+1)^{2}}\right]
\]
都是绝对收敛的.

事实上, 对于 (2),
\[
  \begin{aligned}
    \prod_{n=1}^{\infty}\left[1-\frac{1}{(n+1)^{2}}\right]&=\lim_{n\to+\infty}\prod_{k=1}^{n}\left[1-\frac{1}{(k+1)^{2}}\right]\\
    &=\lim_{n\to+\infty}\prod_{k=1}^{n}\left(1-\frac{1}{k+1}\right)\left(1+\frac{1}{k+1}\right)=\lim_{n\to+\infty}\prod_{k=1}^{n}\frac{k}{k+1}\prod_{k=1}^{n}\frac{k+2}{k+1}\\
    &=\lim_{n\to+\infty}\left(\frac{1}{2}\cdot\frac{2}{3}\cdot\frac{3}{4}\cdot\frac{4}{5}\cdots\frac{n}{n+1}\right)\left(\frac{3}{2}\cdot\frac{4}{3}\cdot\frac{5}{4}\cdots\frac{n+1}{n}\cdot\frac{n+2}{n+1}\right)\\
    &=\lim_{n\to+\infty}\left(\frac{1}{2}\cdot\frac{n+2}{n+1}\right)=\frac{1}{2},
  \end{aligned}
\]
即 $\prod_{n=1}^{\infty}\left[1-\frac{1}{(n+1)^{2}}\right]$ 收敛于 $\frac{1}{2}$.

\begin{example}\label{ex:12.5.5}
证明:
(1) $\prod_{n=1}^{\infty}\left[1+\frac{1}{(n+1)^{\alpha}}\right]$, (2) $\prod_{n=1}^{\infty}\left[1-\frac{1}{(n+1)^{\alpha}}\right]$,
当 $\alpha>1$ 时都绝对收敛; 当 $\alpha<1$ 时都发散.
\end{example}

证明 \extref{例 1.4.2} 表明 $\sum_{n=1}^{\infty}\frac{1}{(n+1)^{\alpha}}$ 当 $\alpha>1$ 时收敛; 当 $\alpha<1$ 时发散. 再根据定理~\ref{thm:12.5.3} 知, $\prod_{n=1}^{\infty}\left[1+\frac{1}{(n+1)^{\alpha}}\right]$ 与 $\prod_{n=1}^{\infty}\left[1-\frac{1}{(n+1)^{\alpha}}\right]$ 当 $\alpha>1$ 时都绝对收敛; 当 $\alpha\leqslant1$ 时都发散.

\begin{example}\label{ex:12.5.6}
设 $\alpha>0$, 讨论无穷乘积
\[
  \prod_{n=1}^{\infty}\left[1+(-1)^{n-1}\frac{1}{n^{\alpha}}\right]
\]
的敛散性.
\end{example}

解 记 $a_{n}=(-1)^{n-1}\frac{1}{n^{\alpha}}$, 则由 Leibniz 判别法知 $\sum_{n=1}^{\infty}a_{n}$ 收敛. 当 $\alpha>\frac{1}{2}$ 时, $\sum_{n=1}^{\infty}a_{n}^{2}=\sum_{n=1}^{\infty}\frac{1}{n^{2\alpha}}$ 收敛. 根据定理~\ref{thm:12.5.4}(1), 无穷级数 $\prod_{n=1}^{\infty}\left[1+(-1)^{n-1}\frac{1}{n^{\alpha}}\right]$ 是收敛的. 但是, 由于当 $\frac{1}{2}<\alpha\leqslant1$ 时, $\prod_{n=1}^{\infty}(1+|a_{n}|)=\prod_{n=1}^{\infty}\left(1+\frac{1}{n^{\alpha}}\right)$ 发散立知, $\prod_{n=1}^{\infty}\left[1+(-1)^{n-1}\frac{1}{n^{\alpha}}\right]$ 是条件收敛的.

当 $0<\alpha\leqslant\frac{1}{2}$ 时, 由于 $\sum_{n=1}^{\infty}a_{n}^{2}=\sum_{n=1}^{\infty}\frac{1}{n^{2\alpha}}$ 发散, 而 $\sum_{n=1}^{\infty}a_{n}=\sum_{n=1}^{\infty}(-1)^{n-1}\frac{1}{n^{\alpha}}$ 收敛, 根据定理~\ref{thm:12.5.4}(2), 无穷级数 $\prod_{n=1}^{\infty}(1+a_{n})$ 发散于 $0$.

综合上述与例~\ref{ex:12.5.5} 的结论,
\[
  \prod_{n=1}^{\infty}\left[1+(-1)^{n-1}\frac{1}{n^{\alpha}}\right]
\]
当 $\alpha>1$ 时绝对收敛; 当 $\frac{1}{2}<\alpha\leqslant1$ 时, 条件收敛; 当 $0<\alpha\leqslant\frac{1}{2}$ 时发散于 $0$.

\begin{example}\label{ex:12.5.7}
设 $\alpha>\beta>0$. 证明:
\[
  \prod_{n=1}^{\infty}\frac{\beta+n}{\alpha+n}=\lim_{n\to+\infty}\frac{(\beta+1)(\beta+2)\cdots(\beta+n)}{(\alpha+1)(\alpha+2)\cdots(\alpha+n)}=0.
\]
\end{example}

证明 考虑无穷乘积
\[
  \prod_{n=1}^{\infty}\frac{\beta+n}{\alpha+n}=\prod_{n=1}^{\infty}\left(1+\frac{\beta-\alpha}{\alpha+n}\right).
\]
令 $a_{n}=\frac{\beta-\alpha}{\alpha+n}$, 则 $a_{n}<0$ 且 $\sum_{n=1}^{\infty}\frac{\beta-\alpha}{\alpha+n}$ 发散. 由注~\ref{rem:12.5.1} 知, 无穷乘积 $\prod_{n=1}^{\infty}\frac{\beta+n}{\alpha+n}$ 发散于 $0$, 即
\[
  \prod_{n=1}^{\infty}\frac{\beta+n}{\alpha+n}=\lim_{n\to+\infty}\frac{(\beta+1)(\beta+2)\cdots(\beta+n)}{(\alpha+1)(\alpha+2)\cdots(\alpha+n)}=0.
\]

\begin{example}\label{ex:12.5.8}
讨论 $\prod_{n=1}^{\infty}(1+x^{2^{n-1}})$ 的敛散性.
\end{example}

解 $P_{n}=(1+x)(1+x^{2})(1+x^{4})(1+x^{8})\cdots(1+x^{2^{n-1}})$.

当 $x=1$ 时, $P_{n}=2^{n}\to+\infty$（$n\to+\infty$）, 无穷乘积发散于 $+\infty$.

当 $x=-1$ 时, 因为 $1+x=0$, 所以 $P_{n}=0$（$n\in\mathbb{N}$）, $\prod_{n=1}^{\infty}(1+x^{2^{n-1}})=\lim\limits_{n\to+\infty}P_{n}=\lim\limits_{n\to+\infty}0=0$, 无穷乘积发散于 $0$.

当 $x\neq\pm1$ 时,
\[
  \begin{aligned}
    P_{n}&=\frac{(1-x)(1+x)(1+x^{2})(1+x^{4})\cdots(1+x^{2^{n-1}})}{1-x}\\
    &=\frac{1-x^{2^{n}}}{1-x}\xrightarrow{}\left\{\begin{array}
{ll}
      \dfrac{1}{1-x}, & |x|<1,\\[8pt]
      +\infty, & x>1,\\[2pt]
      -\infty, & x<-1.
\end{array}\right.
  \end{aligned}
\]
于是, $\prod_{n=1}^{\infty}(1+x^{2^{n-1}})$ 当且仅当 $|x|<1$ 时收敛.

\begin{example}\label{ex:12.5.9}
设 $\alpha>-1$,
\[
  \mathrm{C}_{\alpha}^{0}=1,\quad\mathrm{C}_{\alpha}^{n}=\frac{\alpha(\alpha-1)\cdots(\alpha-n+1)}{n!},
\]
证明: $\lim\limits_{n\to+\infty}C_{\alpha}^{n}=0$.
\end{example}

证明 易见
\[
  \begin{aligned}
    (-1)^{n}C_{\alpha}^{n}&=(-1)^{n}\frac{\alpha(\alpha-1)\cdots(\alpha-n+1)}{n!}\\
    &=\left(1-\frac{\alpha+1}{1}\right)\left(1-\frac{\alpha+1}{2}\right)\cdots\left(1-\frac{\alpha+1}{n}\right)\\
    &=\prod_{k=1}^{n}\left(1-\frac{\alpha+1}{k}\right).
  \end{aligned}
\]
因为 $\alpha+1>0$, 且 $\sum_{n=1}^{\infty}\frac{\alpha+1}{n}$ 发散, 故由注~\ref{rem:12.5.1} 知无穷乘积 $\prod_{k=1}^{\infty}\left(1-\frac{\alpha+1}{k}\right)$ 发散于 $0$, 即
\[
  \lim_{n\to+\infty}C_{\alpha}^{n}=\lim_{n\to+\infty}(-1)^{n}\prod_{k=1}^{n}\left(1-\frac{\alpha+1}{k}\right)=0.
\]

\begin{example}\label{ex:12.5.10}
证明: $\prod_{n=1}^{\infty}\frac{(2n)^{2}}{(2n-1)(2n+1)}=\frac{\pi}{2}$（Wallis 公式）.
\end{example}

证明 当 $0<\alpha<\frac{\pi}{2}$ 时, $\sin^{2n+1}x<\sin^{2n}x<\sin^{2n-1}x$, 积分得
\[
  \int_{0}^{\frac{\pi}{2}}\sin^{2n+1}x\,\mathrm{d}x<\int_{0}^{\frac{\pi}{2}}\sin^{2n}x\,\mathrm{d}x<\int_{0}^{\frac{\pi}{2}}\sin^{2n-1}x\,\mathrm{d}x,
\]
即
\[
  \frac{(2n)!!}{(2n+1)!!}<\frac{(2n-1)!!}{(2n)!!}\frac{\pi}{2}<\frac{(2n-2)!!}{(2n-1)!!},
\]
于是,
\[
  \frac{1}{2n+1}\left[\frac{(2n)!!}{(2n-1)!!}\right]^{2}<\frac{\pi}{2}<\left[\frac{(2n)!!}{(2n-1)!!}\right]^{2}\cdot\frac{1}{2n}.
\]
又因为
\[
  P_{n}=\prod_{k=1}^{n}\frac{(2k)^{2}}{(2k-1)(2k+1)}=\frac{1}{2n+1}\left[\frac{(2n)!!}{(2n-1)!!}\right]^{2},
\]
所以与上式比较得到
\[
  P_{n}<\frac{\pi}{2}<P_{n}\cdot\frac{2n+1}{2n},
\]
\[
  \frac{2n}{2n+1}\cdot\frac{\pi}{2}<P_{n}<\frac{\pi}{2}.
\]
根据 $\lim\limits_{n\to+\infty}\frac{2n}{2n+1}\cdot\frac{\pi}{2}=\frac{\pi}{2}=\lim\limits_{n\to+\infty}\frac{\pi}{2}$ 与夹逼定理推得
\[
  \prod_{n=1}^{\infty}\frac{(2n)^{2}}{(2n-1)(2n+1)}=\lim_{n\to+\infty}P_{n}=\frac{\pi}{2}.
\]

\begin{exblock}{练习题 12.5}
\begin{enumerate}
  \item 讨论下列无穷乘积的敛散性:
  \begin{enumerate}
    \item $\displaystyle\prod_{n=1}^{\infty}\frac{1}{n}$;
    \item $\displaystyle\prod_{n=1}^{\infty}\frac{(n+1)^{2}}{n(n+2)}$;
    \item $\displaystyle\prod_{n=1}^{\infty}\sqrt[n]{1+\frac{1}{n}}$.
  \end{enumerate}
  \item 证明下列等式:
  \begin{enumerate}
    \item $\displaystyle\prod_{n=2}^{\infty}\frac{n^{2}-1}{n^{2}+1}=\frac{2}{3}$;
    \item $\displaystyle\prod_{n=2}^{\infty}\left(1-\frac{2}{n(n+1)}\right)=\frac{1}{3}$;
    \item $\displaystyle\prod_{n=0}^{\infty}\left(1+\left(\frac{1}{2}\right)^{2^{n}}\right)=2$.
  \end{enumerate}
  \item 应用 Wallis 公式证明:
  \begin{enumerate}
    \item $\displaystyle\prod_{n=1}^{\infty}\left(1-\frac{1}{4n^{2}}\right)=\frac{2}{\pi}$;
    \item $\displaystyle\prod_{n=1}^{\infty}\left(1-\frac{1}{(2n+1)^{2}}\right)=\frac{\pi}{4}$.
  \end{enumerate}
  \item 设 $\sum_{n=1}^{\infty}a_{n}^{2}<+\infty$, 证明: $\prod_{n=1}^{\infty}\cos a_{n}$ 收敛.
  \item 能否由 $\prod_{n=1}^{\infty}p_{n}$, $\prod_{n=1}^{\infty}q_{n}$ 的收敛性推得下列乘积:
  \[
    \prod_{n=1}^{\infty}(p_{n}+q_{n}),\quad\prod_{n=1}^{\infty}p_{n}q_{n},\quad\prod_{n=1}^{\infty}\frac{p_{n}}{q_{n}}
\]
  的收敛性?
\end{enumerate}
\end{exblock}

\begin{exblock}{思考题 12.5}
\begin{enumerate}
  \item 讨论下列无穷乘积的敛散性:
  \begin{enumerate}
    \item $\displaystyle\prod_{n=1}^{\infty}\frac{n}{\sqrt{n^{2}+1}}$;
    \item $\displaystyle\prod_{n=2}^{\infty}\left(\frac{n^{2}-1}{n^{2}+1}\right)^{p}$, $p$ 为任意实数;
    \item $\displaystyle\prod_{n=1}^{\infty}\sqrt[n]{\ln(n+x)-\ln n}$, $x>0$.
  \end{enumerate}
  \item 证明下列等式:
  \begin{enumerate}
    \item $\displaystyle\prod_{n=1}^{\infty}\cos\frac{x}{2^{n}}=\frac{\sin x}{x}$, $x\neq0$;
    \item $\displaystyle\frac{2}{\sqrt{2}}\cdot\frac{2}{\sqrt{2+\sqrt{2}}}\cdot\frac{2}{\sqrt{2+\sqrt{2+\sqrt{2}}}}\cdots=\frac{\pi}{2}$;
    \item $\displaystyle\prod_{n=1}^{\infty}\left(1+\frac{1}{a_{n}}\right)=\mathrm{e}$, 其中 $a_{1}=1$, $a_{n}=n(a_{n-1}+1)$.
  \end{enumerate}
  （提示: 先证 $P_{n}=\prod_{k=1}^{n}\left(1+\frac{1}{a_{k}}\right)=\frac{a_{n}+1}{n!}$.）
  \item 设
  \[
    a_{n}=\left\{\begin{array}
{ll}
      -\dfrac{1}{\sqrt{k}}, & n=2k-1,\\[8pt]
      \dfrac{1}{\sqrt{k}}+\dfrac{1}{k}+\dfrac{1}{k\sqrt{k}}, & n=2k,
\end{array}\right.
\]
  证明: $\sum_{n=1}^{\infty}a_{n}$, $\sum_{n=1}^{\infty}a_{n}^{2}$ 都发散, 但 $\prod_{n=1}^{\infty}(1+a_{n})$ 收敛.
  \item 设级数 $\sum_{n=1}^{\infty}x_{n}^{2}$ 收敛, 证明: 无穷乘积 $\prod_{n=1}^{\infty}\cos x_{n}$ 也收敛.
  \item 设级数 $\sum_{n=1}^{\infty}a_{n}$ 绝对收敛, 证明: 无穷乘积 $\prod_{n=1}^{\infty}\tan\left(\frac{\pi}{4}+a_{n}\right)$ 也收敛, 其中 $|a_{n}|<\frac{\pi}{4}$.
  \item 证明: 级数
  \[
    \left(\frac{1}{\sqrt{2}}+\frac{1}{2}\right)+\left(-\frac{1}{\sqrt{2}}\right)+\left(\frac{1}{\sqrt{3}}+\frac{1}{3}\right)+\left(-\frac{1}{\sqrt{3}}\right)+\cdots
\]
  发散, 但无穷乘积
  \[
    \left(1+\frac{1}{\sqrt{2}}+\frac{1}{2}\right)\left(1-\frac{1}{\sqrt{2}}\right)\left(1+\frac{1}{\sqrt{3}}+\frac{1}{3}\right)\left(1-\frac{1}{\sqrt{3}}\right)\cdots
\]
  却收敛.
\end{enumerate}
\end{exblock}

% ==================================================================
%  徐森林、金亚东、薛春华《数学分析》（第三册）
%  复习题 12
%  源：原书印刷页 73-77（PDF 物理页 79-83）
%  底本：PaddleOCR-VL-1.6，见 chapters/_ch12_12_5_raw.md 末段
%  约定：只用 exblock，不设 section 命令（故不入目录）
%  已订正的 OCR 错误与原书照录清单见 tools/_make_ch12_ex.py 头部说明
% ==================================================================
\begin{exblock}{复习题 12}
\begin{enumerate}
  \item 设 $\sum_{n=1}^{\infty}a_{n}$ 为正项级数, 如果存在正数 $\alpha$ 与 $\beta$, 使得 $a_{n}-a_{n+1}\geqslant\beta a_{n}^{2-\alpha}$. 证明:
  \begin{enumerate}
    \item $\{a_{n}\}$ 单调减趋于 $0$, 且 $\alpha\leqslant1$;
    \item $\sum_{n=1}^{\infty}a_{n}<+\infty$, 且
    \[
      \sum_{n=N}^{\infty}a_{n}=O(a_{N}^{\alpha}),\quad N\to+\infty.
\]
  \end{enumerate}
  \item 试作一个收敛级数 $\sum_{n=1}^{\infty}a_{n}$, 使得 $\sum_{n=1}^{\infty}a_{n}^{3}$ 发散.
  \item 设 $\{b_{n}\}$ 单调趋于 $0$, 且 $\sum_{n=1}^{\infty}a_{n}b_{n}$ 收敛. 应用 Cauchy 收敛原理与 Abel 引理证明:
  \[
    \lim_{n\to+\infty}(a_{1}+a_{2}+\cdots+a_{n})b_{n}=0.
\]
  \item 设 $\sum_{n=1}^{\infty}\frac{a_{n}}{n^{\alpha}}$（$\alpha>0$）收敛. 证明:
  \[
    \lim_{n\to+\infty}\frac{a_{1}+a_{2}+\cdots+a_{n}}{n^{\alpha}}=0.
\]
  \item 设 $\{b_{n}\}$ 为正的单调增趋于 $+\infty$ 的数列, 如果 $\sum_{n=1}^{\infty}a_{n}$ 收敛. 证明:
  \[
    \lim_{n\to+\infty}\frac{a_{1}b_{1}+a_{2}b_{2}+\cdots+a_{n}b_{n}}{b_{n}}=0.
\]
  注意: 当 $b_{n}=n$ 时, 就得到练习题 12.3 中题 8 的结论. 因此, 本题是它的推广.
  \item 设 $\sum_{n=1}^{\infty}a_{n}$ 收敛, 令
  \[
    b_{n}=\frac{a_{1}+2a_{2}+\cdots+na_{n}}{n(n+1)}.
\]
  证明:
  \begin{enumerate}
    \item $\sum_{k=1}^{\infty}b_{k}=\sum_{k=1}^{n}a_{k}-\frac{a_{1}+2a_{2}+\cdots+na_{n}}{n+1}$;
    \item $\sum_{n=1}^{\infty}b_{n}$ 收敛, 且 $\sum_{n=1}^{\infty}b_{n}=\sum_{n=1}^{\infty}a_{n}$.
  \end{enumerate}
  \item 设级数 $\sum_{n=1}^{\infty}a_{n}$ 绝对收敛, 且 $\lim\limits_{n\to+\infty}b_{n}=0$. 证明:
  \[
    \lim_{n\to+\infty}(a_{1}b_{n}+a_{2}b_{n-1}+\cdots+a_{n}b_{1})=0.
\]
  \item 如果对任意一个趋于 $0$ 的数列 $\{x_{n}\}$, 级数 $\sum_{n=1}^{\infty}a_{n}x_{n}$ 都收敛. 试用反证法证明 $\sum_{n=1}^{\infty}a_{n}$ 一定绝对收敛.

  如果将条件中的“任意一个趋于 0 的数列”改为“任意一个单调减趋于 0 的数列”, 结论是否成立?
  \item 证明: $\prod_{n=1}^{\infty}\frac{1}{\mathrm{e}}\left(1+\frac{1}{n}\right)^{n}=0$, 并由此推得 $\lim\limits_{n\to+\infty}\frac{n^{n}}{n!}\mathrm{e}^{-n}=0$.
  \item 设级数
  \[
    \sum_{n=1}^{\infty}a_{n}(x^{2}-1)(x^{2}-2^{2})\cdots(x^{2}-n^{2})
\]
  在 $x=x_{0}$（非整数）处收敛, 证明: 该级数对所有的 $x$ 都收敛.
  \item 设正数列 $\{a_{n}\}$ 满足:
  \[
    \frac{a_{n}}{a_{n+1}}=1+\frac{1}{n}+O(b_{n}),\quad n\to+\infty,
\]
  其中 $\sum_{n=1}^{\infty}b_{n}$ 是一个绝对收敛级数. 证明: $\sum_{n=1}^{\infty}a_{n}$ 发散.
  \item 设 $p_{1}=2$, $p_{2}=3$, $p_{3}=5$, $\cdots$ 为从小到大排列的全体素（质）数, 证明:
  \begin{enumerate}
    \item $\displaystyle\prod_{n=1}^{\infty}\frac{1}{1-\frac{1}{p_{n}^{x}}}=\sum_{n=1}^{\infty}\frac{1}{n^{x}}=\zeta(x)$（Riemann $\zeta$ 函数）
    当 $x>1$ 时收敛; 当 $0<x\leqslant1$ 时发散到 $+\infty$.
    \item $\displaystyle\sum_{k=1}^{\infty}\frac{1}{p_{k}}=\frac{1}{2}+\frac{1}{3}+\frac{1}{5}+\frac{1}{7}+\frac{1}{11}+\cdots=+\infty$.
  \end{enumerate}
  \item 证明: 级数
  \[
    \sum_{n=1}^{\infty}(2-\mathrm{e}^{\alpha})\left(2-\mathrm{e}^{\frac{\alpha}{2}}\right)\cdots\left(2-\mathrm{e}^{\frac{\alpha}{n}}\right)
\]
  当 $\alpha=N\ln2$（$N$ 为自然数）或 $\alpha>1$ 时收敛; 当 $\alpha\leqslant1$ 且 $\alpha\neq N\ln2$（$N$ 为自然数）时发散.
  \item 如果对满足 $\sum_{n=1}^{\infty}a_{n}^{2}$ 收敛的任意 $a_{n}$, 必有 $\sum_{n=1}^{\infty}a_{n}b_{n}$ 收敛, 证明: $\sum_{n=1}^{\infty}b_{n}^{2}$ 也收敛.
  \item 如果对任意满足 $\int_{a}^{+\infty}f^{2}(x)\mathrm{d}x$ 收敛的 $f$, 必有 $\int_{a}^{+\infty}f(x)g(x)\mathrm{d}x$ 收敛, 证明: $\int_{a}^{+\infty}g^{2}(x)\mathrm{d}x$ 也收敛.
  \item 设 $a_{n}\geqslant0$, $n=1,2,\cdots$, $\sum_{n=1}^{\infty}a_{n}^{2}<+\infty$. 证明:
  \begin{enumerate}
    \item $\displaystyle\sum_{n=1}^{\infty}\frac{\sqrt{m}}{\sqrt{n}(m+n)}\leqslant\pi$, $m=1,2,\cdots$;
    \item $\displaystyle\sum_{m,n=1}^{\infty}\frac{a_{m}a_{n}}{m+n}\leqslant\pi\sum_{n=1}^{\infty}a_{n}^{2}$.
  \end{enumerate}
  \item 设 $a_{1}=1$, $a_{2}=2$, $a_{n}=a_{n-2}+a_{n-1}$（$n\geqslant3$）, 证明: $\frac{a_{n}}{a_{n+1}}\leqslant\frac{2}{3}$, 且 $\sum_{n=1}^{\infty}\frac{1}{a_{n}}$ 收敛.
  \item 设 $\{a_{n}\}$ 为正实数序列, $\sum_{n=1}^{\infty}\frac{1}{a_{n}}$ 收敛. 证明: 级数
  \[
    \sum_{n=1}^{\infty}\frac{n^{2}}{(a_{1}+a_{2}+\cdots+a_{n})^{2}}a_{n}
\]
  也收敛.
  \item 设级数 $\sum_{n=1}^{\infty}a_{n}$ 收敛, $\sum_{n=1}^{\infty}(b_{n+1}-b_{n})$ 收敛. 证明: 级数 $\sum_{n=1}^{\infty}a_{n}b_{n}$ 也收敛.
  \item 设 $\alpha>\beta>1$,
  \[
    I_{n}=\int_{0}^{\frac{\pi}{2}}\frac{\mathrm{d}t}{1+(n\pi+t)^{\alpha}\sin^{\beta}t},\quad
    J_{n}=\int_{\frac{\pi}{2}}^{\pi}\frac{\mathrm{d}t}{1+(n\pi+t)^{\alpha}\sin^{\beta}t}.
\]
  证明:
  \begin{enumerate}
    \item $\displaystyle\int_{0}^{+\infty}\frac{\mathrm{d}x}{1+x^{\alpha}|\sin x|^{\beta}}=\sum_{n=0}^{\infty}I_{n}+\sum_{n=0}^{\infty}J_{n}$;
    \item $\displaystyle\sum_{n=0}^{\infty}I_{n}$, $\sum_{n=0}^{\infty}J_{n}$, $\displaystyle\int_{0}^{+\infty}\frac{\mathrm{d}x}{1+x^{\alpha}|\sin x|^{\beta}}$ 均收敛.
  \end{enumerate}
  \item 证明:
  \begin{enumerate}
    \item $\displaystyle\left|x\sin\frac{1}{x^{2}}-\frac{1}{x}\cos\frac{1}{x^{2}}\right|\geqslant\frac{1}{x}\left|\cos\frac{1}{x^{2}}\right|-x$;
    \item $\displaystyle\int_{0}^{1}\left|x\sin\frac{1}{x^{2}}-\frac{1}{x}\cos\frac{1}{x^{2}}\right|\mathrm{d}x$ 发散.
  \end{enumerate}
  \item 证明: $\displaystyle\sum_{k=0}^{\infty}\int_{2k\pi}^{(2k+1)\pi}\mathrm{e}^{-\frac{x}{2}}\frac{|\sin x-\cos x|}{\sqrt{\sin x}}\,\mathrm{d}x=\frac{2\cdot\sqrt[4]{8}\,\mathrm{e}^{-\frac{\pi}{8}}}{1-\mathrm{e}^{-\pi}}$.
  \item 设单调函数 $f(x)$ 在 $[0,+\infty)$ 上有定义, 而广义积分 $\int_{0}^{+\infty}f(x)\,\mathrm{d}x$ 存在且有限. 证明:
  \[
    \lim_{h\to0}h\bigl[f(h)+f(2h)+\cdots\bigr]=\int_{0}^{+\infty}f(x)\,\mathrm{d}x.
\]
  \item 将题 23 应用于单调函数 $f(x)=\frac{\mathrm{e}^{-x}}{1+\mathrm{e}^{-x}}$, 证明:
  \[
    \lim_{t\to1^{-}}(1-t)\left(\frac{t}{1+t}+\frac{t^{2}}{1+t^{2}}+\cdots+\frac{t^{n}}{1+t^{n}}+\cdots\right)=\ln2.
\]
  \item 设数列 $\{a_{n}\}$ 单调减少趋于 $0$, 且
  \[
    b_{k}=a_{k}-2a_{k+1}+a_{k+2}\geqslant0,\quad k=1,2,\cdots,
\]
  证明: $\sum_{k=1}^{\infty}kb_{k}=a_{1}$.
  \item 设 $\{a_{n}\}$ 为实数列, 它满足不等式
  \[
    0\leqslant a_{k}\leqslant100a_{n},\quad\text{其中}\quad n\leqslant k\leqslant2n,\ n=1,2,\cdots.
\]
  又级数 $\sum_{n=0}^{\infty}a_{n}$ 收敛. 证明: $\lim\limits_{n\to+\infty}na_{n}=0$.
\end{enumerate}
\end{exblock}

